% Generated by shm_tables.m. % Generated by shm_tables.m. % Generated by shm_tables.m. % Generated by shm_tables.m. \input{numbers} in the preamble, then \shm{key} in the text.
\makeatletter
\providecommand{\shm}[1]{\@ifundefined{shm@#1}{\textbf{??#1??}}{\csname shm@#1\endcsname}}
\expandafter\def\csname shm@angles\endcsname{15}% switching angles per quarter cycle
\expandafter\def\csname shm@orders\endcsname{16}% odd non-triplen orders in THD49 (5..49)
\expandafter\def\csname shm@ma.lo\endcsname{0.70}% lowest modulation index of the study
\expandafter\def\csname shm@ma.hi\endcsname{1.20}% highest modulation index of the study
\expandafter\def\csname shm@ma.step\endcsname{0.0025}% modulation-index step of the traced grid
\expandafter\def\csname shm@grid.points\endcsname{201}% grid points along Ma
\expandafter\def\csname shm@anchors\endcsname{11}% multistart anchors
\expandafter\def\csname shm@nstart\endcsname{150}% random starts per anchor
\expandafter\def\csname shm@distinct.mean\endcsname{104.2}% distinct local optima per anchor, mean
\expandafter\def\csname shm@distinct.max\endcsname{112}% distinct local optima at one anchor, max
\expandafter\def\csname shm@hitrate\endcsname{1.1}% % of converged random starts that reach the best optimum
\expandafter\def\csname shm@spread\endcsname{9.41}% median minus best THD49 among the distinct optima, pp
\expandafter\def\csname shm@branches\endcsname{76}% distinct continuous branches traced
\expandafter\def\csname shm@owners\endcsname{17}% branches that own part of the per-Ma optimum
\expandafter\def\csname shm@ownerchanges\endcsname{21}% owner changes of the per-Ma optimum over the grid
\expandafter\def\csname shm@jump.mean\endcsname{27.2}% mean jump of the optimal angle vector at an owner change, deg
\expandafter\def\csname shm@jump.max\endcsname{49.2}% largest jump of the optimal angle vector, deg
\expandafter\def\csname shm@stored\endcsname{13}% branches kept in the stored subset
\expandafter\def\csname shm@stored.gap.mean\endcsname{0.0004}% mean gap of the stored subset to the full envelope, pp
\expandafter\def\csname shm@stored.gap.max\endcsname{0.0282}% worst gap of the stored subset to the full envelope, pp
\expandafter\def\csname shm@coverage\endcsname{100.0}% % of grid points where some branch exists
\expandafter\def\csname shm@env.thd.mean\endcsname{2.13}% mean THD49 of the per-Ma optimum, %
\expandafter\def\csname shm@env.thd.min\endcsname{0.22}% lowest THD49 of the per-Ma optimum, %
\expandafter\def\csname shm@env.thd.max\endcsname{3.86}% highest THD49 of the per-Ma optimum, %
\expandafter\def\csname shm@nfe.branchset\endcsname{17696578}% objective evaluations to build the branch set
\expandafter\def\csname shm@time.branchset\endcsname{29.3}% minutes to build the branch set
\expandafter\def\csname shm@effangles.min\endcsname{12.0}% fewest effective angles on any branch (collapsed pairs removed)
\expandafter\def\csname shm@traces\endcsname{91}% continuation traces before merging (seeds and probes)
\expandafter\def\csname shm@probes\endcsname{50}% modulation indices of the probe multistart
\expandafter\def\csname shm@probe.nstart\endcsname{100}% random starts per probe
\expandafter\def\csname shm@probe.traces\endcsname{11}% traces added because a probe beat the envelope
\expandafter\def\csname shm@probe.hits\endcsname{2}% probes that found an optimum below the envelope
\expandafter\def\csname shm@fold.points\endcsname{24}% grid points searched just beyond the ends of owning branches
\expandafter\def\csname shm@fold.traces\endcsname{2}% traces added by that fold-end search
\expandafter\def\csname shm@dropped\endcsname{0}% traces of one or two grid points that own nothing, dropped
\expandafter\def\csname shm@distinct.mean.0p5\endcsname{102.5}% distinct optima per anchor, clustered at 0.5 deg, mean
\expandafter\def\csname shm@distinct.mean.2\endcsname{99.5}% distinct optima per anchor, clustered at 2 deg, mean
\expandafter\def\csname shm@distinct.levels\endcsname{95.6}% distinct THD49 levels (1e-3 pp apart) per anchor, mean
\expandafter\def\csname shm@distinct.within1pp\endcsname{10.4}% distinct optima within 1 pp of the best, per anchor, mean
\expandafter\def\csname shm@jump.wave.mean\endcsname{15.1}% share of the period in which the optimal voltage changes at an owner change, mean, %
\expandafter\def\csname shm@jump.wave.max\endcsname{48.8}% share of the period in which the optimal voltage changes at an owner change, largest, %
\expandafter\def\csname shm@jump.wave.min\endcsname{4.3}% share of the period in which the optimal voltage changes at an owner change, smallest, %
\expandafter\def\csname shm@jump.wave.at.maxangle\endcsname{4.3}% the same at the owner change with the largest angle move, %
\expandafter\def\csname shm@distinct.once\endcsname{83}% share of the distinct optima reached by a single random start, per anchor, mean, %
\expandafter\def\csname shm@distinct.chao1\endcsname{480}% Chao1 estimate of the total number of local optima per anchor, mean (the distinct count is a lower bound)
\expandafter\def\csname shm@lut.excess.R10\endcsname{21.9503}% LUT at R = 10: mean THD49 excess over the optimum, pp
\expandafter\def\csname shm@lut.excessmax.R10\endcsname{63.309}% LUT at R = 10: worst THD49 excess, pp
\expandafter\def\csname shm@lut.angmax.R10\endcsname{41.19}% LUT at R = 10: worst angle error, deg
\expandafter\def\csname shm@lut.fund.R10\endcsname{0.185}% LUT at R = 10: mean fundamental error, %
\expandafter\def\csname shm@lut.excess.R25\endcsname{13.3422}% LUT at R = 25: mean THD49 excess over the optimum, pp
\expandafter\def\csname shm@lut.excessmax.R25\endcsname{62.555}% LUT at R = 25: worst THD49 excess, pp
\expandafter\def\csname shm@lut.angmax.R25\endcsname{37.65}% LUT at R = 25: worst angle error, deg
\expandafter\def\csname shm@lut.fund.R25\endcsname{0.108}% LUT at R = 25: mean fundamental error, %
\expandafter\def\csname shm@lut.excess.R50\endcsname{7.6717}% LUT at R = 50: mean THD49 excess over the optimum, pp
\expandafter\def\csname shm@lut.excessmax.R50\endcsname{62.767}% LUT at R = 50: worst THD49 excess, pp
\expandafter\def\csname shm@lut.angmax.R50\endcsname{39.05}% LUT at R = 50: worst angle error, deg
\expandafter\def\csname shm@lut.fund.R50\endcsname{0.073}% LUT at R = 50: mean fundamental error, %
\expandafter\def\csname shm@lut.excess.R100\endcsname{3.7973}% LUT at R = 100: mean THD49 excess over the optimum, pp
\expandafter\def\csname shm@lut.excessmax.R100\endcsname{59.862}% LUT at R = 100: worst THD49 excess, pp
\expandafter\def\csname shm@lut.angmax.R100\endcsname{48.42}% LUT at R = 100: worst angle error, deg
\expandafter\def\csname shm@lut.fund.R100\endcsname{0.044}% LUT at R = 100: mean fundamental error, %
\expandafter\def\csname shm@nearest.excess.R10\endcsname{0.0610}% NEAREST at R = 10: mean THD49 excess over the optimum, pp
\expandafter\def\csname shm@nearest.excessmax.R10\endcsname{1.923}% NEAREST at R = 10: worst THD49 excess, pp
\expandafter\def\csname shm@nearest.angmax.R10\endcsname{52.70}% NEAREST at R = 10: worst angle error, deg
\expandafter\def\csname shm@nearest.fund.R10\endcsname{1.567}% NEAREST at R = 10: mean fundamental error, %
\expandafter\def\csname shm@nearest.excess.R25\endcsname{0.0100}% NEAREST at R = 25: mean THD49 excess over the optimum, pp
\expandafter\def\csname shm@nearest.excessmax.R25\endcsname{0.665}% NEAREST at R = 25: worst THD49 excess, pp
\expandafter\def\csname shm@nearest.angmax.R25\endcsname{40.62}% NEAREST at R = 25: worst angle error, deg
\expandafter\def\csname shm@nearest.fund.R25\endcsname{0.637}% NEAREST at R = 25: mean fundamental error, %
\expandafter\def\csname shm@nearest.excess.R50\endcsname{0.0005}% NEAREST at R = 50: mean THD49 excess over the optimum, pp
\expandafter\def\csname shm@nearest.excessmax.R50\endcsname{0.452}% NEAREST at R = 50: worst THD49 excess, pp
\expandafter\def\csname shm@nearest.angmax.R50\endcsname{49.20}% NEAREST at R = 50: worst angle error, deg
\expandafter\def\csname shm@nearest.fund.R50\endcsname{0.364}% NEAREST at R = 50: mean fundamental error, %
\expandafter\def\csname shm@nearest.excess.R100\endcsname{-0.0082}% NEAREST at R = 100: mean THD49 excess over the optimum, pp
\expandafter\def\csname shm@nearest.excessmax.R100\endcsname{0.150}% NEAREST at R = 100: worst THD49 excess, pp
\expandafter\def\csname shm@nearest.angmax.R100\endcsname{47.63}% NEAREST at R = 100: worst angle error, deg
\expandafter\def\csname shm@nearest.fund.R100\endcsname{0.269}% NEAREST at R = 100: mean fundamental error, %
\expandafter\def\csname shm@sqp.excess.R10\endcsname{4.7744}% LUT+SQP at R = 10: mean THD49 excess over the optimum, pp
\expandafter\def\csname shm@sqp.excessmax.R10\endcsname{31.525}% LUT+SQP at R = 10: worst THD49 excess, pp
\expandafter\def\csname shm@sqp.angmax.R10\endcsname{78.09}% LUT+SQP at R = 10: worst angle error, deg
\expandafter\def\csname shm@sqp.fund.R10\endcsname{0.003}% LUT+SQP at R = 10: mean fundamental error, %
\expandafter\def\csname shm@sqp.excess.R25\endcsname{2.4875}% LUT+SQP at R = 25: mean THD49 excess over the optimum, pp
\expandafter\def\csname shm@sqp.excessmax.R25\endcsname{30.501}% LUT+SQP at R = 25: worst THD49 excess, pp
\expandafter\def\csname shm@sqp.angmax.R25\endcsname{78.48}% LUT+SQP at R = 25: worst angle error, deg
\expandafter\def\csname shm@sqp.fund.R25\endcsname{0.000}% LUT+SQP at R = 25: mean fundamental error, %
\expandafter\def\csname shm@sqp.excess.R50\endcsname{1.3759}% LUT+SQP at R = 50: mean THD49 excess over the optimum, pp
\expandafter\def\csname shm@sqp.excessmax.R50\endcsname{29.799}% LUT+SQP at R = 50: worst THD49 excess, pp
\expandafter\def\csname shm@sqp.angmax.R50\endcsname{76.64}% LUT+SQP at R = 50: worst angle error, deg
\expandafter\def\csname shm@sqp.fund.R50\endcsname{0.000}% LUT+SQP at R = 50: mean fundamental error, %
\expandafter\def\csname shm@sqp.excess.R100\endcsname{0.5158}% LUT+SQP at R = 100: mean THD49 excess over the optimum, pp
\expandafter\def\csname shm@sqp.excessmax.R100\endcsname{29.182}% LUT+SQP at R = 100: worst THD49 excess, pp
\expandafter\def\csname shm@sqp.angmax.R100\endcsname{76.44}% LUT+SQP at R = 100: worst angle error, deg
\expandafter\def\csname shm@sqp.fund.R100\endcsname{0.000}% LUT+SQP at R = 100: mean fundamental error, %
\expandafter\def\csname shm@cont.excess.R10\endcsname{2.9051}% CONT at R = 10: mean THD49 excess over the optimum, pp
\expandafter\def\csname shm@cont.excessmax.R10\endcsname{7.882}% CONT at R = 10: worst THD49 excess, pp
\expandafter\def\csname shm@cont.angmax.R10\endcsname{49.71}% CONT at R = 10: worst angle error, deg
\expandafter\def\csname shm@cont.fund.R10\endcsname{0.040}% CONT at R = 10: mean fundamental error, %
\expandafter\def\csname shm@cont.excess.R25\endcsname{1.7159}% CONT at R = 25: mean THD49 excess over the optimum, pp
\expandafter\def\csname shm@cont.excessmax.R25\endcsname{6.730}% CONT at R = 25: worst THD49 excess, pp
\expandafter\def\csname shm@cont.angmax.R25\endcsname{48.59}% CONT at R = 25: worst angle error, deg
\expandafter\def\csname shm@cont.fund.R25\endcsname{0.007}% CONT at R = 25: mean fundamental error, %
\expandafter\def\csname shm@cont.excess.R50\endcsname{1.3021}% CONT at R = 50: mean THD49 excess over the optimum, pp
\expandafter\def\csname shm@cont.excessmax.R50\endcsname{5.197}% CONT at R = 50: worst THD49 excess, pp
\expandafter\def\csname shm@cont.angmax.R50\endcsname{48.43}% CONT at R = 50: worst angle error, deg
\expandafter\def\csname shm@cont.fund.R50\endcsname{0.003}% CONT at R = 50: mean fundamental error, %
\expandafter\def\csname shm@cont.excess.R100\endcsname{1.1490}% CONT at R = 100: mean THD49 excess over the optimum, pp
\expandafter\def\csname shm@cont.excessmax.R100\endcsname{3.979}% CONT at R = 100: worst THD49 excess, pp
\expandafter\def\csname shm@cont.angmax.R100\endcsname{48.42}% CONT at R = 100: worst angle error, deg
\expandafter\def\csname shm@cont.fund.R100\endcsname{0.002}% CONT at R = 100: mean fundamental error, %
\expandafter\def\csname shm@branch.excess.R10\endcsname{1.2992}% BRANCH at R = 10: mean THD49 excess over the optimum, pp
\expandafter\def\csname shm@branch.excessmax.R10\endcsname{13.031}% BRANCH at R = 10: worst THD49 excess, pp
\expandafter\def\csname shm@branch.angmax.R10\endcsname{52.62}% BRANCH at R = 10: worst angle error, deg
\expandafter\def\csname shm@branch.fund.R10\endcsname{0.004}% BRANCH at R = 10: mean fundamental error, %
\expandafter\def\csname shm@branch.excess.R25\endcsname{0.3848}% BRANCH at R = 25: mean THD49 excess over the optimum, pp
\expandafter\def\csname shm@branch.excessmax.R25\endcsname{10.243}% BRANCH at R = 25: worst THD49 excess, pp
\expandafter\def\csname shm@branch.angmax.R25\endcsname{52.62}% BRANCH at R = 25: worst angle error, deg
\expandafter\def\csname shm@branch.fund.R25\endcsname{0.018}% BRANCH at R = 25: mean fundamental error, %
\expandafter\def\csname shm@branch.excess.R50\endcsname{0.1626}% BRANCH at R = 50: mean THD49 excess over the optimum, pp
\expandafter\def\csname shm@branch.excessmax.R50\endcsname{6.860}% BRANCH at R = 50: worst THD49 excess, pp
\expandafter\def\csname shm@branch.angmax.R50\endcsname{52.62}% BRANCH at R = 50: worst angle error, deg
\expandafter\def\csname shm@branch.fund.R50\endcsname{0.003}% BRANCH at R = 50: mean fundamental error, %
\expandafter\def\csname shm@branch.excess.R100\endcsname{0.0706}% BRANCH at R = 100: mean THD49 excess over the optimum, pp
\expandafter\def\csname shm@branch.excessmax.R100\endcsname{2.023}% BRANCH at R = 100: worst THD49 excess, pp
\expandafter\def\csname shm@branch.angmax.R100\endcsname{50.94}% BRANCH at R = 100: worst angle error, deg
\expandafter\def\csname shm@branch.fund.R100\endcsname{0.003}% BRANCH at R = 100: mean fundamental error, %
\expandafter\def\csname shm@lut.floor.ratio\endcsname{0.17}% LUT mean excess at R = 100 divided by that at R = 10 (1 = no improvement at all)
\expandafter\def\csname shm@lut.worst.ratio\endcsname{0.95}% LUT worst excess at R = 100 divided by that at R = 10
\expandafter\def\csname shm@lut.slope\endcsname{-0.76}% log-log slope of the LUT mean excess vs R (resolution error alone: -2)
\expandafter\def\csname shm@branch.slope\endcsname{-2.00}% log-log slope of the branch-aware mean excess vs stored sets
\expandafter\def\csname shm@split.slope.same\endcsname{-2.78}% log-log slope of the same-owner (resolution) error vs R, median |excess|
\expandafter\def\csname shm@split.slope.straddle\endcsname{-0.05}% log-log slope of the straddling (structural) error vs R, median |excess|
\expandafter\def\csname shm@split.slope.same.mean\endcsname{-1.74}% log-log slope of the same-owner error vs R, mean
\expandafter\def\csname shm@split.slope.straddle.mean\endcsname{-0.03}% log-log slope of the straddling error vs R, mean
\expandafter\def\csname shm@split.frac.R10\endcsname{89.0}% % of test points whose bracketing rows straddle an owner change, R = 10
\expandafter\def\csname shm@split.straddle.R10\endcsname{24.661}% mean excess of straddling intervals, R = 10, pp
\expandafter\def\csname shm@split.same.R10\endcsname{0.00210}% mean excess of same-owner intervals, R = 10, pp
\expandafter\def\csname shm@split.straddle.median.R10\endcsname{19.070}% median |excess| of straddling intervals, R = 10, pp
\expandafter\def\csname shm@split.same.median.R10\endcsname{2.19e-03}% median |excess| of same-owner intervals, R = 10, pp
\expandafter\def\csname shm@split.same.p90.R10\endcsname{3.76e-03}% 90th percentile |excess| of same-owner intervals, R = 10, pp
\expandafter\def\csname shm@split.frac.R25\endcsname{46.0}% % of test points whose bracketing rows straddle an owner change, R = 25
\expandafter\def\csname shm@split.straddle.R25\endcsname{28.989}% mean excess of straddling intervals, R = 25, pp
\expandafter\def\csname shm@split.same.R25\endcsname{0.00123}% mean excess of same-owner intervals, R = 25, pp
\expandafter\def\csname shm@split.straddle.median.R25\endcsname{27.360}% median |excess| of straddling intervals, R = 25, pp
\expandafter\def\csname shm@split.same.median.R25\endcsname{3.63e-05}% median |excess| of same-owner intervals, R = 25, pp
\expandafter\def\csname shm@split.same.p90.R25\endcsname{1.95e-03}% 90th percentile |excess| of same-owner intervals, R = 25, pp
\expandafter\def\csname shm@split.frac.R50\endcsname{27.8}% % of test points whose bracketing rows straddle an owner change, R = 50
\expandafter\def\csname shm@split.straddle.R50\endcsname{27.581}% mean excess of straddling intervals, R = 50, pp
\expandafter\def\csname shm@split.same.R50\endcsname{0.00017}% mean excess of same-owner intervals, R = 50, pp
\expandafter\def\csname shm@split.straddle.median.R50\endcsname{24.542}% median |excess| of straddling intervals, R = 50, pp
\expandafter\def\csname shm@split.same.median.R50\endcsname{1.09e-05}% median |excess| of same-owner intervals, R = 50, pp
\expandafter\def\csname shm@split.same.p90.R50\endcsname{9.76e-05}% 90th percentile |excess| of same-owner intervals, R = 50, pp
\expandafter\def\csname shm@split.frac.R100\endcsname{16.8}% % of test points whose bracketing rows straddle an owner change, R = 100
\expandafter\def\csname shm@split.straddle.R100\endcsname{22.560}% mean excess of straddling intervals, R = 100, pp
\expandafter\def\csname shm@split.same.R100\endcsname{0.00005}% mean excess of same-owner intervals, R = 100, pp
\expandafter\def\csname shm@split.straddle.median.R100\endcsname{16.736}% median |excess| of straddling intervals, R = 100, pp
\expandafter\def\csname shm@split.same.median.R100\endcsname{3.22e-06}% median |excess| of same-owner intervals, R = 100, pp
\expandafter\def\csname shm@split.same.p90.R100\endcsname{1.29e-05}% 90th percentile |excess| of same-owner intervals, R = 100, pp
\expandafter\def\csname shm@split.straddle.Rmax\endcsname{22.560}% straddling-interval excess at R = 100, pp
\expandafter\def\csname shm@split.same.Rmax\endcsname{0.00005}% same-owner-interval excess at R = 100, pp
\expandafter\def\csname shm@basin.onowner\endcsname{61.4}% % of corrected points (R = 25) that land on the owning branch
\expandafter\def\csname shm@basin.onowner.near\endcsname{14.3}% % landing on the owning branch within two grid steps of an owner change
\expandafter\def\csname shm@basin.onowner.far\endcsname{79.5}% % landing on the owning branch farther from an owner change
\expandafter\def\csname shm@basin.excess.wrong\endcsname{6.438}% mean THD49 excess where the corrector landed on another branch, pp
\expandafter\def\csname shm@basin.ok.near\endcsname{28.6}% % of corrected points within 0.01 pp of the optimum, within two grid steps of an owner change
\expandafter\def\csname shm@basin.ok.far\endcsname{79.5}% % of corrected points within 0.01 pp of the optimum, farther from an owner change
\expandafter\def\csname shm@cont.cost.mean\endcsname{0.999}% mean THD49 cost of the continuity-constrained path, pp
\expandafter\def\csname shm@cont.cost.max\endcsname{2.282}% worst THD49 cost of the continuity-constrained path, pp
\expandafter\def\csname shm@cont.breaks\endcsname{0}% places where even the continuity-constrained path must jump
\expandafter\def\csname shm@cont.jumps.env\endcsname{20}% jumps (> 5 deg) the per-Ma optimum makes over the grid
\expandafter\def\csname shm@cont.jumps.zero\endcsname{11}% fewest jumps of a path whose mean THD cost is below 1e-4 pp (the rest only swap near-equal branches)
\expandafter\def\csname shm@cont.jumps.for.1\endcsname{0}% fewest jumps that keep the mean THD49 cost at or below 1 pp
\expandafter\def\csname shm@cont.jumps.for.0p1\endcsname{5}% fewest jumps that keep the mean THD49 cost at or below 0.1 pp
\expandafter\def\csname shm@cont.jumps.for.0p01\endcsname{8}% fewest jumps that keep the mean THD49 cost at or below 0.01 pp
\expandafter\def\csname shm@cont.cost.k1\endcsname{0.688}% mean THD49 cost of the best path allowed one jump, pp
\expandafter\def\csname shm@reg1.linear.K41\endcsname{0.0101}% 1-D regression, K = 41: angle MAE of linear, deg
\expandafter\def\csname shm@reg1.linear.K6\endcsname{0.1820}% 1-D regression, K = 6: angle MAE of linear, deg
\expandafter\def\csname shm@reg1.spline.K41\endcsname{0.0077}% 1-D regression, K = 41: angle MAE of spline, deg
\expandafter\def\csname shm@reg1.spline.K6\endcsname{0.1351}% 1-D regression, K = 6: angle MAE of spline, deg
\expandafter\def\csname shm@reg1.pchip.K41\endcsname{0.0068}% 1-D regression, K = 41: angle MAE of pchip, deg
\expandafter\def\csname shm@reg1.pchip.K6\endcsname{0.1359}% 1-D regression, K = 6: angle MAE of pchip, deg
\expandafter\def\csname shm@reg1.RBF.K41\endcsname{0.0092}% 1-D regression, K = 41: angle MAE of RBF, deg
\expandafter\def\csname shm@reg1.RBF.K6\endcsname{0.1741}% 1-D regression, K = 6: angle MAE of RBF, deg
\expandafter\def\csname shm@reg1.GP.K41\endcsname{0.0114}% 1-D regression, K = 41: angle MAE of GP, deg
\expandafter\def\csname shm@reg1.GP.K6\endcsname{0.2451}% 1-D regression, K = 6: angle MAE of GP, deg
\expandafter\def\csname shm@reg1.tanhLM.K41\endcsname{0.0119}% 1-D regression, K = 41: angle MAE of tanhLM, deg
\expandafter\def\csname shm@reg1.tanhLM.K6\endcsname{0.1463}% 1-D regression, K = 6: angle MAE of tanhLM, deg
\expandafter\def\csname shm@reg1.MLP.K41\endcsname{0.0441}% 1-D regression, K = 41: angle MAE of MLP, deg
\expandafter\def\csname shm@reg1.MLP.K6\endcsname{0.5077}% 1-D regression, K = 6: angle MAE of MLP, deg
\expandafter\def\csname shm@reg1.ANFIS.K41\endcsname{0.0140}% 1-D regression, K = 41: angle MAE of ANFIS, deg
\expandafter\def\csname shm@reg1.ANFIS.K6\endcsname{0.3298}% 1-D regression, K = 6: angle MAE of ANFIS, deg
\expandafter\def\csname shm@reg1.mlp.ratio\endcsname{6.4}% MLP angle MAE divided by the best interpolant (pchip) at K = 41
\expandafter\def\csname shm@reg1.mlp.seedcv\endcsname{23}% seed-to-seed coefficient of variation of the MLP MAE, %
\expandafter\def\csname shm@reg2.linear\endcsname{0.0555}% 2-D regression, largest lattice (64 points): angle MAE of linear, deg
\expandafter\def\csname shm@reg2.spline\endcsname{0.0384}% 2-D regression, largest lattice (64 points): angle MAE of spline, deg
\expandafter\def\csname shm@reg2.RBF\endcsname{0.0468}% 2-D regression, largest lattice (64 points): angle MAE of RBF, deg
\expandafter\def\csname shm@reg2.GP\endcsname{0.0598}% 2-D regression, largest lattice (64 points): angle MAE of GP, deg
\expandafter\def\csname shm@reg2.tanhLM\endcsname{0.0485}% 2-D regression, largest lattice (64 points): angle MAE of tanhLM, deg
\expandafter\def\csname shm@reg2.MLP\endcsname{0.0774}% 2-D regression, largest lattice (64 points): angle MAE of MLP, deg
\expandafter\def\csname shm@reg2.ANFIS\endcsname{0.0553}% 2-D regression, largest lattice (64 points): angle MAE of ANFIS, deg
\expandafter\def\csname shm@reg3.linear\endcsname{0.0189}% 3-D regression, largest lattice (216 points): angle MAE of linear, deg
\expandafter\def\csname shm@reg3.spline\endcsname{0.0120}% 3-D regression, largest lattice (216 points): angle MAE of spline, deg
\expandafter\def\csname shm@reg3.RBF\endcsname{0.0119}% 3-D regression, largest lattice (216 points): angle MAE of RBF, deg
\expandafter\def\csname shm@reg3.GP\endcsname{0.0208}% 3-D regression, largest lattice (216 points): angle MAE of GP, deg
\expandafter\def\csname shm@reg3.tanhLM\endcsname{0.0246}% 3-D regression, largest lattice (216 points): angle MAE of tanhLM, deg
\expandafter\def\csname shm@reg3.MLP\endcsname{0.0197}% 3-D regression, largest lattice (216 points): angle MAE of MLP, deg
\expandafter\def\csname shm@reg3.ANFIS\endcsname{0.0265}% 3-D regression, largest lattice (216 points): angle MAE of ANFIS, deg
\expandafter\def\csname shm@dstar.sel.K9\endcsname{4.33}% dimension from which a 9-per-axis label table outgrows the classifier
\expandafter\def\csname shm@dstar.reg.K9\endcsname{3.25}% dimension from which a 9-per-axis angle table outgrows one MLP per branch
\expandafter\def\csname shm@dstar.sel.K17\endcsname{3.36}% dimension from which a 17-per-axis label table outgrows the classifier
\expandafter\def\csname shm@dstar.reg.K17\endcsname{2.52}% dimension from which a 17-per-axis angle table outgrows one MLP per branch
\expandafter\def\csname shm@dstar.sel.K25\endcsname{2.95}% dimension from which a 25-per-axis label table outgrows the classifier
\expandafter\def\csname shm@dstar.reg.K25\endcsname{2.22}% dimension from which a 25-per-axis angle table outgrows one MLP per branch
\expandafter\def\csname shm@t.lut\endcsname{34.4}% table lookup + linear interpolation, microseconds per call
\expandafter\def\csname shm@t.pchip\endcsname{71.5}% pchip evaluation inside a branch, microseconds per call
\expandafter\def\csname shm@t.cls\endcsname{86.4}% softmax classifier forward pass, microseconds per call
\expandafter\def\csname shm@t.mlp\endcsname{267.8}% MLP forward pass, microseconds per call
\expandafter\def\csname shm@t.sqp\endcsname{30771.7}% SQP corrector (fmincon, warm start), microseconds per call
\expandafter\def\csname shm@dense.points\endcsname{61}% off-grid points of the independent multistart check
\expandafter\def\csname shm@dense.nstart\endcsname{200}% fresh random starts per off-grid point
\expandafter\def\csname shm@dense.gap.mean\endcsname{-0.4249}% stored branches minus independent multistart, mean, pp
\expandafter\def\csname shm@dense.gap.worst\endcsname{0.0022}% stored branches minus independent multistart, worst, pp
\expandafter\def\csname shm@dense.better\endcsname{32}% off-grid points where the stored branches beat the multistart
\expandafter\def\csname shm@dense.over1\endcsname{0}% off-grid points where the stored branches lose by more than 1 pp
\expandafter\def\csname shm@dense.over0p1\endcsname{0}% off-grid points where the stored branches lose by more than 0.1 pp
\expandafter\def\csname shm@dense.gap.worst.interp\endcsname{0.0268}% the same without the corrector: the table's own interpolation, worst, pp
\expandafter\def\csname shm@dense.better.max\endcsname{1.76}% largest margin by which the stored branches beat the fresh multistart, pp
\expandafter\def\csname shm@dense.interp.used\endcsname{0}% off-grid points whose stored value is a branch's interpolated angles (its corrector failed or left it)
\expandafter\def\csname shm@dense.corr.fail\endcsname{1}% off-grid points where the corrector of the best interpolated branch did not converge or left the branch
\expandafter\def\csname shm@spec.ma\endcsname{1.00}% Ma of the spectrum audit quoted
\expandafter\def\csname shm@spec.thd49\endcsname{2.44}% THD49 at that Ma, %
\expandafter\def\csname shm@spec.thd99\endcsname{25.7}% THD to the 99th at that Ma, %
\expandafter\def\csname shm@spec.thdfull\endcsname{34.2}% THD to the 999th at that Ma, %
\expandafter\def\csname shm@spec.wthd\endcsname{0.499}% weighted THD (1/n) to the 999th at that Ma, %
\expandafter\def\csname shm@spec.oob\endcsname{17.1}% largest single harmonic above the 49th, % of b1
\expandafter\def\csname shm@spec.oob.order\endcsname{55}% its order
\expandafter\def\csname shm@comply\endcsname{89.6}% % of the envelope meeting the per-order limits (0.2+32.5/n above the 25th)
\expandafter\def\csname shm@comply.alt\endcsname{35.8}% % meeting them with 0.2+12.5/n above the 25th
\expandafter\def\csname shm@bound.ub.model\endcsname{89.875}% the model's upper angle bound, 90 - theta_gap/2, deg
\expandafter\def\csname shm@bound.ub.min\endcsname{80}% tightest upper angle bound tested, deg
\expandafter\def\csname shm@bound.cost.mean\endcsname{0.825}% mean rise of the best THD49 from the model's bound to the tightest one, pp
\expandafter\def\csname shm@bound.cost.max\endcsname{1.688}% largest rise of the best THD49 from the model's bound to the tightest one, pp
\expandafter\def\csname shm@bound.cost89.max\endcsname{1.030}% largest rise of the best THD49 if the angles were kept below 89 deg, pp
\expandafter\def\csname shm@sw.dev\endcsname{1.6e-15}% worst relative deviation, closed form vs exact Fourier integral of the gate pattern
\expandafter\def\csname shm@sw.fftdev\endcsname{0.0096}% worst |THD49 sampled FFT - model| at zero dead time, pp
\expandafter\def\csname shm@sw.dthd.2us\endcsname{0.092}% largest change of THD49 of the optimal patterns caused by a 2 us dead time, pp
\expandafter\def\csname shm@sw.dthd.4us\endcsname{0.089}% largest change of THD49 of the optimal patterns caused by a 4 us dead time, pp
\expandafter\def\csname shm@sw.loss.1us\endcsname{0.0}% % of the envelope whose narrowest pulse is shorter than a 1 us dead time
\expandafter\def\csname shm@sw.loss.2us\endcsname{0.0}% % of the envelope whose narrowest pulse is shorter than a 2 us dead time
\expandafter\def\csname shm@sw.loss.4us\endcsname{0.0}% % of the envelope whose narrowest pulse is shorter than a 4 us dead time
\expandafter\def\csname shm@sw.loss.8us\endcsname{0.0}% % of the envelope whose narrowest pulse is shorter than a 8 us dead time
\expandafter\def\csname shm@sim.cases\endcsname{9}% steady-state cases whose circuit ran the commanded angles
\expandafter\def\csname shm@sim.invalid\endcsname{0}% cases left out: the circuit did not run the commanded angles (must be 0)
\expandafter\def\csname shm@sim.vdc\endcsname{800}% DC-link voltage of the circuit model, V
\expandafter\def\csname shm@sim.ts.us\endcsname{1}% time step of the logged line-to-line voltage, us
\expandafter\def\csname shm@sim.opt.dev.max\endcsname{0.098}% optimal patterns: largest |circuit - analytic| THD49, pp
\expandafter\def\csname shm@sim.opt.samp.max\endcsname{0.004}% optimal patterns: largest |sampled ideal - analytic| THD49 (the time step), pp
\expandafter\def\csname shm@sim.opt.circ.max\endcsname{0.101}% optimal patterns: largest |circuit - sampled ideal| THD49 (what the circuit adds), pp
\expandafter\def\csname shm@sim.opt.ma.dev\endcsname{0.23}% optimal patterns: largest relative error of the circuit's fundamental, %
\expandafter\def\csname shm@sim.lut.rows\endcsname{100}% rows of the largest table simulated at its worst point
\expandafter\def\csname shm@sim.lut.ma\endcsname{0.7925}% Ma of that worst point
\expandafter\def\csname shm@sim.lut.model\endcsname{60.90}% that table's interpolated angles: analytic THD49, %
\expandafter\def\csname shm@sim.lut.circuit\endcsname{60.92}% that table's interpolated angles: circuit THD49, %
\expandafter\def\csname shm@sim.lut.ithd\endcsname{5.53}% that table's interpolated angles: phase-current THD49, %
\expandafter\def\csname shm@sim.branch.model\endcsname{1.04}% branch-aware table at the same Ma: analytic THD49, %
\expandafter\def\csname shm@sim.branch.circuit\endcsname{0.98}% branch-aware table at the same Ma: circuit THD49, %
\expandafter\def\csname shm@sim.branch.ithd\endcsname{0.08}% branch-aware table at the same Ma: phase-current THD49, %
\expandafter\def\csname shm@sim.np.pp.max\endcsname{13.6}% largest peak-to-peak NP voltage over the steady-state windows, V
\expandafter\def\csname shm@sim.step.wave.pct\endcsname{48.8}% branch switch (the owner change that changes the waveform most): share of the period that changes, %
\expandafter\def\csname shm@sim.step.jump.deg\endcsname{30.7}% branch switch: largest angle move at the switch, deg
\expandafter\def\csname shm@sim.step.peak.jump\endcsname{0.37}% branch switch: peak phase-current deviation from the new steady state, % of I1
\expandafter\def\csname shm@sim.step.pre.jump\endcsname{2.38}% branch switch: how far apart the two steady-state currents are, % of I1
\expandafter\def\csname shm@sim.step.offset.jump\endcsname{0.05}% branch switch: mean current offset over the first cycle after the switch, % of I1
\expandafter\def\csname shm@sim.step.peak.stay\endcsname{0.22}% the same Ma step on the old branch (no branch switch): peak phase-current deviation, % of I1
\expandafter\def\csname shm@sim.step.angle.deg\endcsname{49.2}% second branch switch, at the largest angle move: that move, deg
\expandafter\def\csname shm@sim.step.angle.wave.pct\endcsname{4.3}% second branch switch: share of the period that changes, %
\expandafter\def\csname shm@sim.step.angle.peak\endcsname{0.18}% second branch switch: peak phase-current deviation from the new steady state, % of I1
\expandafter\def\csname shm@sim.step.angle.pre\endcsname{0.92}% second branch switch: how far apart the two steady-state currents are, % of I1
\expandafter\def\csname shm@ni.R\endcsname{100}% rows (sets) of the tables evaluated under non-ideal switching
\expandafter\def\csname shm@ni.conds\endcsname{6}% non-ideal conditions evaluated over the whole range
\expandafter\def\csname shm@ni.ideal.opt.dmax\endcsname{0.100}% ideal switches: optimum, largest |THD49 - closed form|, pp
\expandafter\def\csname shm@ni.ideal.opt.b1\endcsname{0.06}% ideal switches: optimum, largest |fundamental error|, %
\expandafter\def\csname shm@ni.ideal.lut.mean\endcsname{3.87}% ideal switches: table, mean excess over the optimum, pp
\expandafter\def\csname shm@ni.ideal.lut.max\endcsname{62.9}% ideal switches: table, worst excess over the optimum, pp
\expandafter\def\csname shm@ni.ideal.br.mean\endcsname{0.057}% ideal switches: branch-aware table, mean excess, pp
\expandafter\def\csname shm@ni.ideal.br.max\endcsname{2.85}% ideal switches: branch-aware table, worst excess, pp
\expandafter\def\csname shm@ni.drops.opt.dmax\endcsname{0.153}% datasheet devices: optimum, largest |THD49 - closed form|, pp
\expandafter\def\csname shm@ni.drops.opt.b1\endcsname{0.55}% datasheet devices: optimum, largest |fundamental error|, %
\expandafter\def\csname shm@ni.drops.lut.mean\endcsname{3.89}% datasheet devices: table, mean excess over the optimum, pp
\expandafter\def\csname shm@ni.drops.lut.max\endcsname{63.3}% datasheet devices: table, worst excess over the optimum, pp
\expandafter\def\csname shm@ni.drops.br.mean\endcsname{0.056}% datasheet devices: branch-aware table, mean excess, pp
\expandafter\def\csname shm@ni.drops.br.max\endcsname{2.82}% datasheet devices: branch-aware table, worst excess, pp
\expandafter\def\csname shm@ni.dt.opt.dmax\endcsname{0.455}% + dead time 2 us: optimum, largest |THD49 - closed form|, pp
\expandafter\def\csname shm@ni.dt.opt.b1\endcsname{1.10}% + dead time 2 us: optimum, largest |fundamental error|, %
\expandafter\def\csname shm@ni.dt.lut.mean\endcsname{3.92}% + dead time 2 us: table, mean excess over the optimum, pp
\expandafter\def\csname shm@ni.dt.lut.max\endcsname{63.7}% + dead time 2 us: table, worst excess over the optimum, pp
\expandafter\def\csname shm@ni.dt.br.mean\endcsname{0.057}% + dead time 2 us: branch-aware table, mean excess, pp
\expandafter\def\csname shm@ni.dt.br.max\endcsname{2.77}% + dead time 2 us: branch-aware table, worst excess, pp
\expandafter\def\csname shm@ni.np.opt.dmax\endcsname{0.737}% + NP offset 2%: optimum, largest |THD49 - closed form|, pp
\expandafter\def\csname shm@ni.np.opt.b1\endcsname{0.55}% + NP offset 2%: optimum, largest |fundamental error|, %
\expandafter\def\csname shm@ni.np.lut.mean\endcsname{3.85}% + NP offset 2%: table, mean excess over the optimum, pp
\expandafter\def\csname shm@ni.np.lut.max\endcsname{62.9}% + NP offset 2%: table, worst excess over the optimum, pp
\expandafter\def\csname shm@ni.np.br.mean\endcsname{0.051}% + NP offset 2%: branch-aware table, mean excess, pp
\expandafter\def\csname shm@ni.np.br.max\endcsname{2.65}% + NP offset 2%: branch-aware table, worst excess, pp
\expandafter\def\csname shm@ni.tq.opt.dmax\endcsname{0.173}% + timer 1 us: optimum, largest |THD49 - closed form|, pp
\expandafter\def\csname shm@ni.tq.opt.b1\endcsname{0.61}% + timer 1 us: optimum, largest |fundamental error|, %
\expandafter\def\csname shm@ni.tq.lut.mean\endcsname{3.89}% + timer 1 us: table, mean excess over the optimum, pp
\expandafter\def\csname shm@ni.tq.lut.max\endcsname{63.3}% + timer 1 us: table, worst excess over the optimum, pp
\expandafter\def\csname shm@ni.tq.br.mean\endcsname{0.056}% + timer 1 us: branch-aware table, mean excess, pp
\expandafter\def\csname shm@ni.tq.br.max\endcsname{2.86}% + timer 1 us: branch-aware table, worst excess, pp
\expandafter\def\csname shm@ni.all.opt.dmax\endcsname{0.898}% all combined: optimum, largest |THD49 - closed form|, pp
\expandafter\def\csname shm@ni.all.opt.b1\endcsname{1.13}% all combined: optimum, largest |fundamental error|, %
\expandafter\def\csname shm@ni.all.lut.mean\endcsname{3.87}% all combined: table, mean excess over the optimum, pp
\expandafter\def\csname shm@ni.all.lut.max\endcsname{63.3}% all combined: table, worst excess over the optimum, pp
\expandafter\def\csname shm@ni.all.br.mean\endcsname{0.052}% all combined: branch-aware table, mean excess, pp
\expandafter\def\csname shm@ni.all.br.max\endcsname{2.65}% all combined: branch-aware table, worst excess, pp
\expandafter\def\csname shm@ni.opt.dmax\endcsname{0.898}% largest |THD49 - closed form| of the optimum under any condition, pp
\expandafter\def\csname shm@ni.lut.spread\endcsname{0.07}% spread of the table's mean excess across the conditions, pp
\expandafter\def\csname shm@ni.br.spread\endcsname{0.006}% spread of the branch-aware mean excess across the conditions, pp
\expandafter\def\csname shm@ni.ripple.pp\endcsname{7.3}% largest neutral-point ripple in the switching-function model, V peak to peak
\expandafter\def\csname shm@sim.ni.runs\endcsname{54}% circuit runs with datasheet switches that ran the commanded angles
\expandafter\def\csname shm@sim.ni.invalid\endcsname{0}% runs left out: commanded angles not run (must be 0)
\expandafter\def\csname shm@sim.ni.conds\endcsname{6}% non-ideal conditions simulated
\expandafter\def\csname shm@sim.ni.swfn.max\endcsname{0.013}% circuit minus switching-function model, optimal and branch-aware patterns, largest |difference|, pp
\expandafter\def\csname shm@sim.ni.swfn.lut.rel\endcsname{0.18}% the same for the table's angles, % of their THD49
\expandafter\def\csname shm@sim.ni.offset.max\endcsname{1.49}% largest |neutral-point offset| the circuit had in the analysis window, % of Vdc
\expandafter\def\csname shm@sim.ni.swfn.cmd.max\endcsname{0.606}% circuit minus switching-function model at the COMMANDED offset (not the measured one), largest |difference|, pp
\expandafter\def\csname shm@sim.ni.lut.rows\endcsname{100}% rows of the largest table simulated with datasheet switches
\expandafter\def\csname shm@sim.ni.ideal.opt.dev\endcsname{0.114}% ideal switches: optimal patterns, largest |circuit - closed form| THD49, pp
\expandafter\def\csname shm@sim.ni.ideal.b1\endcsname{0.04}% ideal switches: optimal patterns, largest |fundamental error| in the circuit, %
\expandafter\def\csname shm@sim.ni.ideal.ripple.pp\endcsname{24.3}% ideal switches: largest neutral-point voltage, V peak to peak
\expandafter\def\csname shm@sim.ni.ideal.offset.min\endcsname{-1.49}% ideal switches: smallest measured NP offset, % of Vdc
\expandafter\def\csname shm@sim.ni.ideal.offset.max\endcsname{-0.47}% ideal switches: largest measured NP offset, % of Vdc
\expandafter\def\csname shm@sim.ni.ideal.lut.circuit\endcsname{60.82}% ideal switches: table at its worst point, circuit THD49, %
\expandafter\def\csname shm@sim.ni.ideal.lut.ithd\endcsname{6.49}% ideal switches: the same, phase-current THD49, %
\expandafter\def\csname shm@sim.ni.ideal.branch.circuit\endcsname{1.13}% ideal switches: branch-aware table at that Ma, circuit THD49, %
\expandafter\def\csname shm@sim.ni.ideal.branch.ithd\endcsname{2.44}% ideal switches: the same, phase-current THD49, %
\expandafter\def\csname shm@sim.ni.drops.opt.dev\endcsname{0.149}% datasheet devices: optimal patterns, largest |circuit - closed form| THD49, pp
\expandafter\def\csname shm@sim.ni.drops.b1\endcsname{0.49}% datasheet devices: optimal patterns, largest |fundamental error| in the circuit, %
\expandafter\def\csname shm@sim.ni.drops.ripple.pp\endcsname{23.8}% datasheet devices: largest neutral-point voltage, V peak to peak
\expandafter\def\csname shm@sim.ni.drops.offset.min\endcsname{-1.44}% datasheet devices: smallest measured NP offset, % of Vdc
\expandafter\def\csname shm@sim.ni.drops.offset.max\endcsname{-0.45}% datasheet devices: largest measured NP offset, % of Vdc
\expandafter\def\csname shm@sim.ni.drops.lut.circuit\endcsname{61.14}% datasheet devices: table at its worst point, circuit THD49, %
\expandafter\def\csname shm@sim.ni.drops.lut.ithd\endcsname{6.49}% datasheet devices: the same, phase-current THD49, %
\expandafter\def\csname shm@sim.ni.drops.branch.circuit\endcsname{1.11}% datasheet devices: branch-aware table at that Ma, circuit THD49, %
\expandafter\def\csname shm@sim.ni.drops.branch.ithd\endcsname{2.37}% datasheet devices: the same, phase-current THD49, %
\expandafter\def\csname shm@sim.ni.dt.opt.dev\endcsname{0.180}% + dead time 2 us: optimal patterns, largest |circuit - closed form| THD49, pp
\expandafter\def\csname shm@sim.ni.dt.b1\endcsname{0.95}% + dead time 2 us: optimal patterns, largest |fundamental error| in the circuit, %
\expandafter\def\csname shm@sim.ni.dt.ripple.pp\endcsname{23.6}% + dead time 2 us: largest neutral-point voltage, V peak to peak
\expandafter\def\csname shm@sim.ni.dt.offset.min\endcsname{-1.41}% + dead time 2 us: smallest measured NP offset, % of Vdc
\expandafter\def\csname shm@sim.ni.dt.offset.max\endcsname{-0.45}% + dead time 2 us: largest measured NP offset, % of Vdc
\expandafter\def\csname shm@sim.ni.dt.lut.circuit\endcsname{61.57}% + dead time 2 us: table at its worst point, circuit THD49, %
\expandafter\def\csname shm@sim.ni.dt.lut.ithd\endcsname{6.52}% + dead time 2 us: the same, phase-current THD49, %
\expandafter\def\csname shm@sim.ni.dt.branch.circuit\endcsname{1.17}% + dead time 2 us: branch-aware table at that Ma, circuit THD49, %
\expandafter\def\csname shm@sim.ni.dt.branch.ithd\endcsname{2.36}% + dead time 2 us: the same, phase-current THD49, %
\expandafter\def\csname shm@sim.ni.np.opt.dev\endcsname{0.128}% + NP offset 2%: optimal patterns, largest |circuit - closed form| THD49, pp
\expandafter\def\csname shm@sim.ni.np.b1\endcsname{0.49}% + NP offset 2%: optimal patterns, largest |fundamental error| in the circuit, %
\expandafter\def\csname shm@sim.ni.np.ripple.pp\endcsname{27.2}% + NP offset 2%: largest neutral-point voltage, V peak to peak
\expandafter\def\csname shm@sim.ni.np.offset.min\endcsname{0.03}% + NP offset 2%: smallest measured NP offset, % of Vdc
\expandafter\def\csname shm@sim.ni.np.offset.max\endcsname{1.07}% + NP offset 2%: largest measured NP offset, % of Vdc
\expandafter\def\csname shm@sim.ni.np.lut.circuit\endcsname{61.14}% + NP offset 2%: table at its worst point, circuit THD49, %
\expandafter\def\csname shm@sim.ni.np.lut.ithd\endcsname{6.41}% + NP offset 2%: the same, phase-current THD49, %
\expandafter\def\csname shm@sim.ni.np.branch.circuit\endcsname{1.06}% + NP offset 2%: branch-aware table at that Ma, circuit THD49, %
\expandafter\def\csname shm@sim.ni.np.branch.ithd\endcsname{2.21}% + NP offset 2%: the same, phase-current THD49, %
\expandafter\def\csname shm@sim.ni.tq.opt.dev\endcsname{0.158}% + timer 1 us: optimal patterns, largest |circuit - closed form| THD49, pp
\expandafter\def\csname shm@sim.ni.tq.b1\endcsname{0.47}% + timer 1 us: optimal patterns, largest |fundamental error| in the circuit, %
\expandafter\def\csname shm@sim.ni.tq.ripple.pp\endcsname{23.9}% + timer 1 us: largest neutral-point voltage, V peak to peak
\expandafter\def\csname shm@sim.ni.tq.offset.min\endcsname{-1.44}% + timer 1 us: smallest measured NP offset, % of Vdc
\expandafter\def\csname shm@sim.ni.tq.offset.max\endcsname{-0.45}% + timer 1 us: largest measured NP offset, % of Vdc
\expandafter\def\csname shm@sim.ni.tq.lut.circuit\endcsname{61.17}% + timer 1 us: table at its worst point, circuit THD49, %
\expandafter\def\csname shm@sim.ni.tq.lut.ithd\endcsname{6.49}% + timer 1 us: the same, phase-current THD49, %
\expandafter\def\csname shm@sim.ni.tq.branch.circuit\endcsname{1.12}% + timer 1 us: branch-aware table at that Ma, circuit THD49, %
\expandafter\def\csname shm@sim.ni.tq.branch.ithd\endcsname{2.37}% + timer 1 us: the same, phase-current THD49, %
\expandafter\def\csname shm@sim.ni.all.opt.dev\endcsname{0.168}% all combined: optimal patterns, largest |circuit - closed form| THD49, pp
\expandafter\def\csname shm@sim.ni.all.b1\endcsname{0.93}% all combined: optimal patterns, largest |fundamental error| in the circuit, %
\expandafter\def\csname shm@sim.ni.all.ripple.pp\endcsname{27.0}% all combined: largest neutral-point voltage, V peak to peak
\expandafter\def\csname shm@sim.ni.all.offset.min\endcsname{0.06}% all combined: smallest measured NP offset, % of Vdc
\expandafter\def\csname shm@sim.ni.all.offset.max\endcsname{1.09}% all combined: largest measured NP offset, % of Vdc
\expandafter\def\csname shm@sim.ni.all.lut.circuit\endcsname{61.61}% all combined: table at its worst point, circuit THD49, %
\expandafter\def\csname shm@sim.ni.all.lut.ithd\endcsname{6.45}% all combined: the same, phase-current THD49, %
\expandafter\def\csname shm@sim.ni.all.branch.circuit\endcsname{1.12}% all combined: branch-aware table at that Ma, circuit THD49, %
\expandafter\def\csname shm@sim.ni.all.branch.ithd\endcsname{2.21}% all combined: the same, phase-current THD49, %
\expandafter\def\csname shm@time.total\endcsname{0.00}% hours for the whole run
\makeatother
 in the preamble, then \shm{key} in the text.
\makeatletter
\providecommand{\shm}[1]{\@ifundefined{shm@#1}{\textbf{??#1??}}{\csname shm@#1\endcsname}}
\expandafter\def\csname shm@angles\endcsname{15}% switching angles per quarter cycle
\expandafter\def\csname shm@orders\endcsname{16}% odd non-triplen orders in THD49 (5..49)
\expandafter\def\csname shm@ma.lo\endcsname{0.70}% lowest modulation index of the study
\expandafter\def\csname shm@ma.hi\endcsname{1.20}% highest modulation index of the study
\expandafter\def\csname shm@ma.step\endcsname{0.0025}% modulation-index step of the traced grid
\expandafter\def\csname shm@grid.points\endcsname{201}% grid points along Ma
\expandafter\def\csname shm@anchors\endcsname{11}% multistart anchors
\expandafter\def\csname shm@nstart\endcsname{150}% random starts per anchor
\expandafter\def\csname shm@distinct.mean\endcsname{104.2}% distinct local optima per anchor, mean
\expandafter\def\csname shm@distinct.max\endcsname{112}% distinct local optima at one anchor, max
\expandafter\def\csname shm@hitrate\endcsname{1.1}% % of converged random starts that reach the best optimum
\expandafter\def\csname shm@spread\endcsname{9.41}% median minus best THD49 among the distinct optima, pp
\expandafter\def\csname shm@branches\endcsname{76}% distinct continuous branches traced
\expandafter\def\csname shm@owners\endcsname{17}% branches that own part of the per-Ma optimum
\expandafter\def\csname shm@ownerchanges\endcsname{21}% owner changes of the per-Ma optimum over the grid
\expandafter\def\csname shm@jump.mean\endcsname{27.2}% mean jump of the optimal angle vector at an owner change, deg
\expandafter\def\csname shm@jump.max\endcsname{49.2}% largest jump of the optimal angle vector, deg
\expandafter\def\csname shm@stored\endcsname{13}% branches kept in the stored subset
\expandafter\def\csname shm@stored.gap.mean\endcsname{0.0004}% mean gap of the stored subset to the full envelope, pp
\expandafter\def\csname shm@stored.gap.max\endcsname{0.0282}% worst gap of the stored subset to the full envelope, pp
\expandafter\def\csname shm@coverage\endcsname{100.0}% % of grid points where some branch exists
\expandafter\def\csname shm@env.thd.mean\endcsname{2.13}% mean THD49 of the per-Ma optimum, %
\expandafter\def\csname shm@env.thd.min\endcsname{0.22}% lowest THD49 of the per-Ma optimum, %
\expandafter\def\csname shm@env.thd.max\endcsname{3.86}% highest THD49 of the per-Ma optimum, %
\expandafter\def\csname shm@nfe.branchset\endcsname{17696578}% objective evaluations to build the branch set
\expandafter\def\csname shm@time.branchset\endcsname{29.3}% minutes to build the branch set
\expandafter\def\csname shm@effangles.min\endcsname{12.0}% fewest effective angles on any branch (collapsed pairs removed)
\expandafter\def\csname shm@traces\endcsname{91}% continuation traces before merging (seeds and probes)
\expandafter\def\csname shm@probes\endcsname{50}% modulation indices of the probe multistart
\expandafter\def\csname shm@probe.nstart\endcsname{100}% random starts per probe
\expandafter\def\csname shm@probe.traces\endcsname{11}% traces added because a probe beat the envelope
\expandafter\def\csname shm@probe.hits\endcsname{2}% probes that found an optimum below the envelope
\expandafter\def\csname shm@fold.points\endcsname{24}% grid points searched just beyond the ends of owning branches
\expandafter\def\csname shm@fold.traces\endcsname{2}% traces added by that fold-end search
\expandafter\def\csname shm@dropped\endcsname{0}% traces of one or two grid points that own nothing, dropped
\expandafter\def\csname shm@distinct.mean.0p5\endcsname{102.5}% distinct optima per anchor, clustered at 0.5 deg, mean
\expandafter\def\csname shm@distinct.mean.2\endcsname{99.5}% distinct optima per anchor, clustered at 2 deg, mean
\expandafter\def\csname shm@distinct.levels\endcsname{95.6}% distinct THD49 levels (1e-3 pp apart) per anchor, mean
\expandafter\def\csname shm@distinct.within1pp\endcsname{10.4}% distinct optima within 1 pp of the best, per anchor, mean
\expandafter\def\csname shm@jump.wave.mean\endcsname{15.1}% share of the period in which the optimal voltage changes at an owner change, mean, %
\expandafter\def\csname shm@jump.wave.max\endcsname{48.8}% share of the period in which the optimal voltage changes at an owner change, largest, %
\expandafter\def\csname shm@jump.wave.min\endcsname{4.3}% share of the period in which the optimal voltage changes at an owner change, smallest, %
\expandafter\def\csname shm@jump.wave.at.maxangle\endcsname{4.3}% the same at the owner change with the largest angle move, %
\expandafter\def\csname shm@distinct.once\endcsname{83}% share of the distinct optima reached by a single random start, per anchor, mean, %
\expandafter\def\csname shm@distinct.chao1\endcsname{480}% Chao1 estimate of the total number of local optima per anchor, mean (the distinct count is a lower bound)
\expandafter\def\csname shm@lut.excess.R10\endcsname{21.9503}% LUT at R = 10: mean THD49 excess over the optimum, pp
\expandafter\def\csname shm@lut.excessmax.R10\endcsname{63.309}% LUT at R = 10: worst THD49 excess, pp
\expandafter\def\csname shm@lut.angmax.R10\endcsname{41.19}% LUT at R = 10: worst angle error, deg
\expandafter\def\csname shm@lut.fund.R10\endcsname{0.185}% LUT at R = 10: mean fundamental error, %
\expandafter\def\csname shm@lut.excess.R25\endcsname{13.3422}% LUT at R = 25: mean THD49 excess over the optimum, pp
\expandafter\def\csname shm@lut.excessmax.R25\endcsname{62.555}% LUT at R = 25: worst THD49 excess, pp
\expandafter\def\csname shm@lut.angmax.R25\endcsname{37.65}% LUT at R = 25: worst angle error, deg
\expandafter\def\csname shm@lut.fund.R25\endcsname{0.108}% LUT at R = 25: mean fundamental error, %
\expandafter\def\csname shm@lut.excess.R50\endcsname{7.6717}% LUT at R = 50: mean THD49 excess over the optimum, pp
\expandafter\def\csname shm@lut.excessmax.R50\endcsname{62.767}% LUT at R = 50: worst THD49 excess, pp
\expandafter\def\csname shm@lut.angmax.R50\endcsname{39.05}% LUT at R = 50: worst angle error, deg
\expandafter\def\csname shm@lut.fund.R50\endcsname{0.073}% LUT at R = 50: mean fundamental error, %
\expandafter\def\csname shm@lut.excess.R100\endcsname{3.7973}% LUT at R = 100: mean THD49 excess over the optimum, pp
\expandafter\def\csname shm@lut.excessmax.R100\endcsname{59.862}% LUT at R = 100: worst THD49 excess, pp
\expandafter\def\csname shm@lut.angmax.R100\endcsname{48.42}% LUT at R = 100: worst angle error, deg
\expandafter\def\csname shm@lut.fund.R100\endcsname{0.044}% LUT at R = 100: mean fundamental error, %
\expandafter\def\csname shm@nearest.excess.R10\endcsname{0.0610}% NEAREST at R = 10: mean THD49 excess over the optimum, pp
\expandafter\def\csname shm@nearest.excessmax.R10\endcsname{1.923}% NEAREST at R = 10: worst THD49 excess, pp
\expandafter\def\csname shm@nearest.angmax.R10\endcsname{52.70}% NEAREST at R = 10: worst angle error, deg
\expandafter\def\csname shm@nearest.fund.R10\endcsname{1.567}% NEAREST at R = 10: mean fundamental error, %
\expandafter\def\csname shm@nearest.excess.R25\endcsname{0.0100}% NEAREST at R = 25: mean THD49 excess over the optimum, pp
\expandafter\def\csname shm@nearest.excessmax.R25\endcsname{0.665}% NEAREST at R = 25: worst THD49 excess, pp
\expandafter\def\csname shm@nearest.angmax.R25\endcsname{40.62}% NEAREST at R = 25: worst angle error, deg
\expandafter\def\csname shm@nearest.fund.R25\endcsname{0.637}% NEAREST at R = 25: mean fundamental error, %
\expandafter\def\csname shm@nearest.excess.R50\endcsname{0.0005}% NEAREST at R = 50: mean THD49 excess over the optimum, pp
\expandafter\def\csname shm@nearest.excessmax.R50\endcsname{0.452}% NEAREST at R = 50: worst THD49 excess, pp
\expandafter\def\csname shm@nearest.angmax.R50\endcsname{49.20}% NEAREST at R = 50: worst angle error, deg
\expandafter\def\csname shm@nearest.fund.R50\endcsname{0.364}% NEAREST at R = 50: mean fundamental error, %
\expandafter\def\csname shm@nearest.excess.R100\endcsname{-0.0082}% NEAREST at R = 100: mean THD49 excess over the optimum, pp
\expandafter\def\csname shm@nearest.excessmax.R100\endcsname{0.150}% NEAREST at R = 100: worst THD49 excess, pp
\expandafter\def\csname shm@nearest.angmax.R100\endcsname{47.63}% NEAREST at R = 100: worst angle error, deg
\expandafter\def\csname shm@nearest.fund.R100\endcsname{0.269}% NEAREST at R = 100: mean fundamental error, %
\expandafter\def\csname shm@sqp.excess.R10\endcsname{4.7744}% LUT+SQP at R = 10: mean THD49 excess over the optimum, pp
\expandafter\def\csname shm@sqp.excessmax.R10\endcsname{31.525}% LUT+SQP at R = 10: worst THD49 excess, pp
\expandafter\def\csname shm@sqp.angmax.R10\endcsname{78.09}% LUT+SQP at R = 10: worst angle error, deg
\expandafter\def\csname shm@sqp.fund.R10\endcsname{0.003}% LUT+SQP at R = 10: mean fundamental error, %
\expandafter\def\csname shm@sqp.excess.R25\endcsname{2.4875}% LUT+SQP at R = 25: mean THD49 excess over the optimum, pp
\expandafter\def\csname shm@sqp.excessmax.R25\endcsname{30.501}% LUT+SQP at R = 25: worst THD49 excess, pp
\expandafter\def\csname shm@sqp.angmax.R25\endcsname{78.48}% LUT+SQP at R = 25: worst angle error, deg
\expandafter\def\csname shm@sqp.fund.R25\endcsname{0.000}% LUT+SQP at R = 25: mean fundamental error, %
\expandafter\def\csname shm@sqp.excess.R50\endcsname{1.3759}% LUT+SQP at R = 50: mean THD49 excess over the optimum, pp
\expandafter\def\csname shm@sqp.excessmax.R50\endcsname{29.799}% LUT+SQP at R = 50: worst THD49 excess, pp
\expandafter\def\csname shm@sqp.angmax.R50\endcsname{76.64}% LUT+SQP at R = 50: worst angle error, deg
\expandafter\def\csname shm@sqp.fund.R50\endcsname{0.000}% LUT+SQP at R = 50: mean fundamental error, %
\expandafter\def\csname shm@sqp.excess.R100\endcsname{0.5158}% LUT+SQP at R = 100: mean THD49 excess over the optimum, pp
\expandafter\def\csname shm@sqp.excessmax.R100\endcsname{29.182}% LUT+SQP at R = 100: worst THD49 excess, pp
\expandafter\def\csname shm@sqp.angmax.R100\endcsname{76.44}% LUT+SQP at R = 100: worst angle error, deg
\expandafter\def\csname shm@sqp.fund.R100\endcsname{0.000}% LUT+SQP at R = 100: mean fundamental error, %
\expandafter\def\csname shm@cont.excess.R10\endcsname{2.9051}% CONT at R = 10: mean THD49 excess over the optimum, pp
\expandafter\def\csname shm@cont.excessmax.R10\endcsname{7.882}% CONT at R = 10: worst THD49 excess, pp
\expandafter\def\csname shm@cont.angmax.R10\endcsname{49.71}% CONT at R = 10: worst angle error, deg
\expandafter\def\csname shm@cont.fund.R10\endcsname{0.040}% CONT at R = 10: mean fundamental error, %
\expandafter\def\csname shm@cont.excess.R25\endcsname{1.7159}% CONT at R = 25: mean THD49 excess over the optimum, pp
\expandafter\def\csname shm@cont.excessmax.R25\endcsname{6.730}% CONT at R = 25: worst THD49 excess, pp
\expandafter\def\csname shm@cont.angmax.R25\endcsname{48.59}% CONT at R = 25: worst angle error, deg
\expandafter\def\csname shm@cont.fund.R25\endcsname{0.007}% CONT at R = 25: mean fundamental error, %
\expandafter\def\csname shm@cont.excess.R50\endcsname{1.3021}% CONT at R = 50: mean THD49 excess over the optimum, pp
\expandafter\def\csname shm@cont.excessmax.R50\endcsname{5.197}% CONT at R = 50: worst THD49 excess, pp
\expandafter\def\csname shm@cont.angmax.R50\endcsname{48.43}% CONT at R = 50: worst angle error, deg
\expandafter\def\csname shm@cont.fund.R50\endcsname{0.003}% CONT at R = 50: mean fundamental error, %
\expandafter\def\csname shm@cont.excess.R100\endcsname{1.1490}% CONT at R = 100: mean THD49 excess over the optimum, pp
\expandafter\def\csname shm@cont.excessmax.R100\endcsname{3.979}% CONT at R = 100: worst THD49 excess, pp
\expandafter\def\csname shm@cont.angmax.R100\endcsname{48.42}% CONT at R = 100: worst angle error, deg
\expandafter\def\csname shm@cont.fund.R100\endcsname{0.002}% CONT at R = 100: mean fundamental error, %
\expandafter\def\csname shm@branch.excess.R10\endcsname{1.2992}% BRANCH at R = 10: mean THD49 excess over the optimum, pp
\expandafter\def\csname shm@branch.excessmax.R10\endcsname{13.031}% BRANCH at R = 10: worst THD49 excess, pp
\expandafter\def\csname shm@branch.angmax.R10\endcsname{52.62}% BRANCH at R = 10: worst angle error, deg
\expandafter\def\csname shm@branch.fund.R10\endcsname{0.004}% BRANCH at R = 10: mean fundamental error, %
\expandafter\def\csname shm@branch.excess.R25\endcsname{0.3848}% BRANCH at R = 25: mean THD49 excess over the optimum, pp
\expandafter\def\csname shm@branch.excessmax.R25\endcsname{10.243}% BRANCH at R = 25: worst THD49 excess, pp
\expandafter\def\csname shm@branch.angmax.R25\endcsname{52.62}% BRANCH at R = 25: worst angle error, deg
\expandafter\def\csname shm@branch.fund.R25\endcsname{0.018}% BRANCH at R = 25: mean fundamental error, %
\expandafter\def\csname shm@branch.excess.R50\endcsname{0.1626}% BRANCH at R = 50: mean THD49 excess over the optimum, pp
\expandafter\def\csname shm@branch.excessmax.R50\endcsname{6.860}% BRANCH at R = 50: worst THD49 excess, pp
\expandafter\def\csname shm@branch.angmax.R50\endcsname{52.62}% BRANCH at R = 50: worst angle error, deg
\expandafter\def\csname shm@branch.fund.R50\endcsname{0.003}% BRANCH at R = 50: mean fundamental error, %
\expandafter\def\csname shm@branch.excess.R100\endcsname{0.0706}% BRANCH at R = 100: mean THD49 excess over the optimum, pp
\expandafter\def\csname shm@branch.excessmax.R100\endcsname{2.023}% BRANCH at R = 100: worst THD49 excess, pp
\expandafter\def\csname shm@branch.angmax.R100\endcsname{50.94}% BRANCH at R = 100: worst angle error, deg
\expandafter\def\csname shm@branch.fund.R100\endcsname{0.003}% BRANCH at R = 100: mean fundamental error, %
\expandafter\def\csname shm@lut.floor.ratio\endcsname{0.17}% LUT mean excess at R = 100 divided by that at R = 10 (1 = no improvement at all)
\expandafter\def\csname shm@lut.worst.ratio\endcsname{0.95}% LUT worst excess at R = 100 divided by that at R = 10
\expandafter\def\csname shm@lut.slope\endcsname{-0.76}% log-log slope of the LUT mean excess vs R (resolution error alone: -2)
\expandafter\def\csname shm@branch.slope\endcsname{-2.00}% log-log slope of the branch-aware mean excess vs stored sets
\expandafter\def\csname shm@split.slope.same\endcsname{-2.78}% log-log slope of the same-owner (resolution) error vs R, median |excess|
\expandafter\def\csname shm@split.slope.straddle\endcsname{-0.05}% log-log slope of the straddling (structural) error vs R, median |excess|
\expandafter\def\csname shm@split.slope.same.mean\endcsname{-1.74}% log-log slope of the same-owner error vs R, mean
\expandafter\def\csname shm@split.slope.straddle.mean\endcsname{-0.03}% log-log slope of the straddling error vs R, mean
\expandafter\def\csname shm@split.frac.R10\endcsname{89.0}% % of test points whose bracketing rows straddle an owner change, R = 10
\expandafter\def\csname shm@split.straddle.R10\endcsname{24.661}% mean excess of straddling intervals, R = 10, pp
\expandafter\def\csname shm@split.same.R10\endcsname{0.00210}% mean excess of same-owner intervals, R = 10, pp
\expandafter\def\csname shm@split.straddle.median.R10\endcsname{19.070}% median |excess| of straddling intervals, R = 10, pp
\expandafter\def\csname shm@split.same.median.R10\endcsname{2.19e-03}% median |excess| of same-owner intervals, R = 10, pp
\expandafter\def\csname shm@split.same.p90.R10\endcsname{3.76e-03}% 90th percentile |excess| of same-owner intervals, R = 10, pp
\expandafter\def\csname shm@split.frac.R25\endcsname{46.0}% % of test points whose bracketing rows straddle an owner change, R = 25
\expandafter\def\csname shm@split.straddle.R25\endcsname{28.989}% mean excess of straddling intervals, R = 25, pp
\expandafter\def\csname shm@split.same.R25\endcsname{0.00123}% mean excess of same-owner intervals, R = 25, pp
\expandafter\def\csname shm@split.straddle.median.R25\endcsname{27.360}% median |excess| of straddling intervals, R = 25, pp
\expandafter\def\csname shm@split.same.median.R25\endcsname{3.63e-05}% median |excess| of same-owner intervals, R = 25, pp
\expandafter\def\csname shm@split.same.p90.R25\endcsname{1.95e-03}% 90th percentile |excess| of same-owner intervals, R = 25, pp
\expandafter\def\csname shm@split.frac.R50\endcsname{27.8}% % of test points whose bracketing rows straddle an owner change, R = 50
\expandafter\def\csname shm@split.straddle.R50\endcsname{27.581}% mean excess of straddling intervals, R = 50, pp
\expandafter\def\csname shm@split.same.R50\endcsname{0.00017}% mean excess of same-owner intervals, R = 50, pp
\expandafter\def\csname shm@split.straddle.median.R50\endcsname{24.542}% median |excess| of straddling intervals, R = 50, pp
\expandafter\def\csname shm@split.same.median.R50\endcsname{1.09e-05}% median |excess| of same-owner intervals, R = 50, pp
\expandafter\def\csname shm@split.same.p90.R50\endcsname{9.76e-05}% 90th percentile |excess| of same-owner intervals, R = 50, pp
\expandafter\def\csname shm@split.frac.R100\endcsname{16.8}% % of test points whose bracketing rows straddle an owner change, R = 100
\expandafter\def\csname shm@split.straddle.R100\endcsname{22.560}% mean excess of straddling intervals, R = 100, pp
\expandafter\def\csname shm@split.same.R100\endcsname{0.00005}% mean excess of same-owner intervals, R = 100, pp
\expandafter\def\csname shm@split.straddle.median.R100\endcsname{16.736}% median |excess| of straddling intervals, R = 100, pp
\expandafter\def\csname shm@split.same.median.R100\endcsname{3.22e-06}% median |excess| of same-owner intervals, R = 100, pp
\expandafter\def\csname shm@split.same.p90.R100\endcsname{1.29e-05}% 90th percentile |excess| of same-owner intervals, R = 100, pp
\expandafter\def\csname shm@split.straddle.Rmax\endcsname{22.560}% straddling-interval excess at R = 100, pp
\expandafter\def\csname shm@split.same.Rmax\endcsname{0.00005}% same-owner-interval excess at R = 100, pp
\expandafter\def\csname shm@basin.onowner\endcsname{61.4}% % of corrected points (R = 25) that land on the owning branch
\expandafter\def\csname shm@basin.onowner.near\endcsname{14.3}% % landing on the owning branch within two grid steps of an owner change
\expandafter\def\csname shm@basin.onowner.far\endcsname{79.5}% % landing on the owning branch farther from an owner change
\expandafter\def\csname shm@basin.excess.wrong\endcsname{6.438}% mean THD49 excess where the corrector landed on another branch, pp
\expandafter\def\csname shm@basin.ok.near\endcsname{28.6}% % of corrected points within 0.01 pp of the optimum, within two grid steps of an owner change
\expandafter\def\csname shm@basin.ok.far\endcsname{79.5}% % of corrected points within 0.01 pp of the optimum, farther from an owner change
\expandafter\def\csname shm@cont.cost.mean\endcsname{0.999}% mean THD49 cost of the continuity-constrained path, pp
\expandafter\def\csname shm@cont.cost.max\endcsname{2.282}% worst THD49 cost of the continuity-constrained path, pp
\expandafter\def\csname shm@cont.breaks\endcsname{0}% places where even the continuity-constrained path must jump
\expandafter\def\csname shm@cont.jumps.env\endcsname{20}% jumps (> 5 deg) the per-Ma optimum makes over the grid
\expandafter\def\csname shm@cont.jumps.zero\endcsname{11}% fewest jumps of a path whose mean THD cost is below 1e-4 pp (the rest only swap near-equal branches)
\expandafter\def\csname shm@cont.jumps.for.1\endcsname{0}% fewest jumps that keep the mean THD49 cost at or below 1 pp
\expandafter\def\csname shm@cont.jumps.for.0p1\endcsname{5}% fewest jumps that keep the mean THD49 cost at or below 0.1 pp
\expandafter\def\csname shm@cont.jumps.for.0p01\endcsname{8}% fewest jumps that keep the mean THD49 cost at or below 0.01 pp
\expandafter\def\csname shm@cont.cost.k1\endcsname{0.688}% mean THD49 cost of the best path allowed one jump, pp
\expandafter\def\csname shm@reg1.linear.K41\endcsname{0.0101}% 1-D regression, K = 41: angle MAE of linear, deg
\expandafter\def\csname shm@reg1.linear.K6\endcsname{0.1820}% 1-D regression, K = 6: angle MAE of linear, deg
\expandafter\def\csname shm@reg1.spline.K41\endcsname{0.0077}% 1-D regression, K = 41: angle MAE of spline, deg
\expandafter\def\csname shm@reg1.spline.K6\endcsname{0.1351}% 1-D regression, K = 6: angle MAE of spline, deg
\expandafter\def\csname shm@reg1.pchip.K41\endcsname{0.0068}% 1-D regression, K = 41: angle MAE of pchip, deg
\expandafter\def\csname shm@reg1.pchip.K6\endcsname{0.1359}% 1-D regression, K = 6: angle MAE of pchip, deg
\expandafter\def\csname shm@reg1.RBF.K41\endcsname{0.0092}% 1-D regression, K = 41: angle MAE of RBF, deg
\expandafter\def\csname shm@reg1.RBF.K6\endcsname{0.1741}% 1-D regression, K = 6: angle MAE of RBF, deg
\expandafter\def\csname shm@reg1.GP.K41\endcsname{0.0114}% 1-D regression, K = 41: angle MAE of GP, deg
\expandafter\def\csname shm@reg1.GP.K6\endcsname{0.2451}% 1-D regression, K = 6: angle MAE of GP, deg
\expandafter\def\csname shm@reg1.tanhLM.K41\endcsname{0.0119}% 1-D regression, K = 41: angle MAE of tanhLM, deg
\expandafter\def\csname shm@reg1.tanhLM.K6\endcsname{0.1463}% 1-D regression, K = 6: angle MAE of tanhLM, deg
\expandafter\def\csname shm@reg1.MLP.K41\endcsname{0.0441}% 1-D regression, K = 41: angle MAE of MLP, deg
\expandafter\def\csname shm@reg1.MLP.K6\endcsname{0.5077}% 1-D regression, K = 6: angle MAE of MLP, deg
\expandafter\def\csname shm@reg1.ANFIS.K41\endcsname{0.0140}% 1-D regression, K = 41: angle MAE of ANFIS, deg
\expandafter\def\csname shm@reg1.ANFIS.K6\endcsname{0.3298}% 1-D regression, K = 6: angle MAE of ANFIS, deg
\expandafter\def\csname shm@reg1.mlp.ratio\endcsname{6.4}% MLP angle MAE divided by the best interpolant (pchip) at K = 41
\expandafter\def\csname shm@reg1.mlp.seedcv\endcsname{23}% seed-to-seed coefficient of variation of the MLP MAE, %
\expandafter\def\csname shm@reg2.linear\endcsname{0.0555}% 2-D regression, largest lattice (64 points): angle MAE of linear, deg
\expandafter\def\csname shm@reg2.spline\endcsname{0.0384}% 2-D regression, largest lattice (64 points): angle MAE of spline, deg
\expandafter\def\csname shm@reg2.RBF\endcsname{0.0468}% 2-D regression, largest lattice (64 points): angle MAE of RBF, deg
\expandafter\def\csname shm@reg2.GP\endcsname{0.0598}% 2-D regression, largest lattice (64 points): angle MAE of GP, deg
\expandafter\def\csname shm@reg2.tanhLM\endcsname{0.0485}% 2-D regression, largest lattice (64 points): angle MAE of tanhLM, deg
\expandafter\def\csname shm@reg2.MLP\endcsname{0.0774}% 2-D regression, largest lattice (64 points): angle MAE of MLP, deg
\expandafter\def\csname shm@reg2.ANFIS\endcsname{0.0553}% 2-D regression, largest lattice (64 points): angle MAE of ANFIS, deg
\expandafter\def\csname shm@reg3.linear\endcsname{0.0189}% 3-D regression, largest lattice (216 points): angle MAE of linear, deg
\expandafter\def\csname shm@reg3.spline\endcsname{0.0120}% 3-D regression, largest lattice (216 points): angle MAE of spline, deg
\expandafter\def\csname shm@reg3.RBF\endcsname{0.0119}% 3-D regression, largest lattice (216 points): angle MAE of RBF, deg
\expandafter\def\csname shm@reg3.GP\endcsname{0.0208}% 3-D regression, largest lattice (216 points): angle MAE of GP, deg
\expandafter\def\csname shm@reg3.tanhLM\endcsname{0.0246}% 3-D regression, largest lattice (216 points): angle MAE of tanhLM, deg
\expandafter\def\csname shm@reg3.MLP\endcsname{0.0197}% 3-D regression, largest lattice (216 points): angle MAE of MLP, deg
\expandafter\def\csname shm@reg3.ANFIS\endcsname{0.0265}% 3-D regression, largest lattice (216 points): angle MAE of ANFIS, deg
\expandafter\def\csname shm@dstar.sel.K9\endcsname{4.33}% dimension from which a 9-per-axis label table outgrows the classifier
\expandafter\def\csname shm@dstar.reg.K9\endcsname{3.25}% dimension from which a 9-per-axis angle table outgrows one MLP per branch
\expandafter\def\csname shm@dstar.sel.K17\endcsname{3.36}% dimension from which a 17-per-axis label table outgrows the classifier
\expandafter\def\csname shm@dstar.reg.K17\endcsname{2.52}% dimension from which a 17-per-axis angle table outgrows one MLP per branch
\expandafter\def\csname shm@dstar.sel.K25\endcsname{2.95}% dimension from which a 25-per-axis label table outgrows the classifier
\expandafter\def\csname shm@dstar.reg.K25\endcsname{2.22}% dimension from which a 25-per-axis angle table outgrows one MLP per branch
\expandafter\def\csname shm@t.lut\endcsname{34.4}% table lookup + linear interpolation, microseconds per call
\expandafter\def\csname shm@t.pchip\endcsname{71.5}% pchip evaluation inside a branch, microseconds per call
\expandafter\def\csname shm@t.cls\endcsname{86.4}% softmax classifier forward pass, microseconds per call
\expandafter\def\csname shm@t.mlp\endcsname{267.8}% MLP forward pass, microseconds per call
\expandafter\def\csname shm@t.sqp\endcsname{30771.7}% SQP corrector (fmincon, warm start), microseconds per call
\expandafter\def\csname shm@dense.points\endcsname{61}% off-grid points of the independent multistart check
\expandafter\def\csname shm@dense.nstart\endcsname{200}% fresh random starts per off-grid point
\expandafter\def\csname shm@dense.gap.mean\endcsname{-0.4249}% stored branches minus independent multistart, mean, pp
\expandafter\def\csname shm@dense.gap.worst\endcsname{0.0022}% stored branches minus independent multistart, worst, pp
\expandafter\def\csname shm@dense.better\endcsname{32}% off-grid points where the stored branches beat the multistart
\expandafter\def\csname shm@dense.over1\endcsname{0}% off-grid points where the stored branches lose by more than 1 pp
\expandafter\def\csname shm@dense.over0p1\endcsname{0}% off-grid points where the stored branches lose by more than 0.1 pp
\expandafter\def\csname shm@dense.gap.worst.interp\endcsname{0.0268}% the same without the corrector: the table's own interpolation, worst, pp
\expandafter\def\csname shm@dense.better.max\endcsname{1.76}% largest margin by which the stored branches beat the fresh multistart, pp
\expandafter\def\csname shm@dense.interp.used\endcsname{0}% off-grid points whose stored value is a branch's interpolated angles (its corrector failed or left it)
\expandafter\def\csname shm@dense.corr.fail\endcsname{1}% off-grid points where the corrector of the best interpolated branch did not converge or left the branch
\expandafter\def\csname shm@spec.ma\endcsname{1.00}% Ma of the spectrum audit quoted
\expandafter\def\csname shm@spec.thd49\endcsname{2.44}% THD49 at that Ma, %
\expandafter\def\csname shm@spec.thd99\endcsname{25.7}% THD to the 99th at that Ma, %
\expandafter\def\csname shm@spec.thdfull\endcsname{34.2}% THD to the 999th at that Ma, %
\expandafter\def\csname shm@spec.wthd\endcsname{0.499}% weighted THD (1/n) to the 999th at that Ma, %
\expandafter\def\csname shm@spec.oob\endcsname{17.1}% largest single harmonic above the 49th, % of b1
\expandafter\def\csname shm@spec.oob.order\endcsname{55}% its order
\expandafter\def\csname shm@comply\endcsname{89.6}% % of the envelope meeting the per-order limits (0.2+32.5/n above the 25th)
\expandafter\def\csname shm@comply.alt\endcsname{35.8}% % meeting them with 0.2+12.5/n above the 25th
\expandafter\def\csname shm@bound.ub.model\endcsname{89.875}% the model's upper angle bound, 90 - theta_gap/2, deg
\expandafter\def\csname shm@bound.ub.min\endcsname{80}% tightest upper angle bound tested, deg
\expandafter\def\csname shm@bound.cost.mean\endcsname{0.825}% mean rise of the best THD49 from the model's bound to the tightest one, pp
\expandafter\def\csname shm@bound.cost.max\endcsname{1.688}% largest rise of the best THD49 from the model's bound to the tightest one, pp
\expandafter\def\csname shm@bound.cost89.max\endcsname{1.030}% largest rise of the best THD49 if the angles were kept below 89 deg, pp
\expandafter\def\csname shm@sw.dev\endcsname{1.6e-15}% worst relative deviation, closed form vs exact Fourier integral of the gate pattern
\expandafter\def\csname shm@sw.fftdev\endcsname{0.0096}% worst |THD49 sampled FFT - model| at zero dead time, pp
\expandafter\def\csname shm@sw.dthd.2us\endcsname{0.092}% largest change of THD49 of the optimal patterns caused by a 2 us dead time, pp
\expandafter\def\csname shm@sw.dthd.4us\endcsname{0.089}% largest change of THD49 of the optimal patterns caused by a 4 us dead time, pp
\expandafter\def\csname shm@sw.loss.1us\endcsname{0.0}% % of the envelope whose narrowest pulse is shorter than a 1 us dead time
\expandafter\def\csname shm@sw.loss.2us\endcsname{0.0}% % of the envelope whose narrowest pulse is shorter than a 2 us dead time
\expandafter\def\csname shm@sw.loss.4us\endcsname{0.0}% % of the envelope whose narrowest pulse is shorter than a 4 us dead time
\expandafter\def\csname shm@sw.loss.8us\endcsname{0.0}% % of the envelope whose narrowest pulse is shorter than a 8 us dead time
\expandafter\def\csname shm@sim.cases\endcsname{9}% steady-state cases whose circuit ran the commanded angles
\expandafter\def\csname shm@sim.invalid\endcsname{0}% cases left out: the circuit did not run the commanded angles (must be 0)
\expandafter\def\csname shm@sim.vdc\endcsname{800}% DC-link voltage of the circuit model, V
\expandafter\def\csname shm@sim.ts.us\endcsname{1}% time step of the logged line-to-line voltage, us
\expandafter\def\csname shm@sim.opt.dev.max\endcsname{0.098}% optimal patterns: largest |circuit - analytic| THD49, pp
\expandafter\def\csname shm@sim.opt.samp.max\endcsname{0.004}% optimal patterns: largest |sampled ideal - analytic| THD49 (the time step), pp
\expandafter\def\csname shm@sim.opt.circ.max\endcsname{0.101}% optimal patterns: largest |circuit - sampled ideal| THD49 (what the circuit adds), pp
\expandafter\def\csname shm@sim.opt.ma.dev\endcsname{0.23}% optimal patterns: largest relative error of the circuit's fundamental, %
\expandafter\def\csname shm@sim.lut.rows\endcsname{100}% rows of the largest table simulated at its worst point
\expandafter\def\csname shm@sim.lut.ma\endcsname{0.7925}% Ma of that worst point
\expandafter\def\csname shm@sim.lut.model\endcsname{60.90}% that table's interpolated angles: analytic THD49, %
\expandafter\def\csname shm@sim.lut.circuit\endcsname{60.92}% that table's interpolated angles: circuit THD49, %
\expandafter\def\csname shm@sim.lut.ithd\endcsname{5.53}% that table's interpolated angles: phase-current THD49, %
\expandafter\def\csname shm@sim.branch.model\endcsname{1.04}% branch-aware table at the same Ma: analytic THD49, %
\expandafter\def\csname shm@sim.branch.circuit\endcsname{0.98}% branch-aware table at the same Ma: circuit THD49, %
\expandafter\def\csname shm@sim.branch.ithd\endcsname{0.08}% branch-aware table at the same Ma: phase-current THD49, %
\expandafter\def\csname shm@sim.np.pp.max\endcsname{13.6}% largest peak-to-peak NP voltage over the steady-state windows, V
\expandafter\def\csname shm@sim.step.wave.pct\endcsname{48.8}% branch switch (the owner change that changes the waveform most): share of the period that changes, %
\expandafter\def\csname shm@sim.step.jump.deg\endcsname{30.7}% branch switch: largest angle move at the switch, deg
\expandafter\def\csname shm@sim.step.peak.jump\endcsname{0.37}% branch switch: peak phase-current deviation from the new steady state, % of I1
\expandafter\def\csname shm@sim.step.pre.jump\endcsname{2.38}% branch switch: how far apart the two steady-state currents are, % of I1
\expandafter\def\csname shm@sim.step.offset.jump\endcsname{0.05}% branch switch: mean current offset over the first cycle after the switch, % of I1
\expandafter\def\csname shm@sim.step.peak.stay\endcsname{0.22}% the same Ma step on the old branch (no branch switch): peak phase-current deviation, % of I1
\expandafter\def\csname shm@sim.step.angle.deg\endcsname{49.2}% second branch switch, at the largest angle move: that move, deg
\expandafter\def\csname shm@sim.step.angle.wave.pct\endcsname{4.3}% second branch switch: share of the period that changes, %
\expandafter\def\csname shm@sim.step.angle.peak\endcsname{0.18}% second branch switch: peak phase-current deviation from the new steady state, % of I1
\expandafter\def\csname shm@sim.step.angle.pre\endcsname{0.92}% second branch switch: how far apart the two steady-state currents are, % of I1
\expandafter\def\csname shm@ni.R\endcsname{100}% rows (sets) of the tables evaluated under non-ideal switching
\expandafter\def\csname shm@ni.conds\endcsname{6}% non-ideal conditions evaluated over the whole range
\expandafter\def\csname shm@ni.ideal.opt.dmax\endcsname{0.100}% ideal switches: optimum, largest |THD49 - closed form|, pp
\expandafter\def\csname shm@ni.ideal.opt.b1\endcsname{0.06}% ideal switches: optimum, largest |fundamental error|, %
\expandafter\def\csname shm@ni.ideal.lut.mean\endcsname{3.87}% ideal switches: table, mean excess over the optimum, pp
\expandafter\def\csname shm@ni.ideal.lut.max\endcsname{62.9}% ideal switches: table, worst excess over the optimum, pp
\expandafter\def\csname shm@ni.ideal.br.mean\endcsname{0.057}% ideal switches: branch-aware table, mean excess, pp
\expandafter\def\csname shm@ni.ideal.br.max\endcsname{2.85}% ideal switches: branch-aware table, worst excess, pp
\expandafter\def\csname shm@ni.drops.opt.dmax\endcsname{0.153}% datasheet devices: optimum, largest |THD49 - closed form|, pp
\expandafter\def\csname shm@ni.drops.opt.b1\endcsname{0.55}% datasheet devices: optimum, largest |fundamental error|, %
\expandafter\def\csname shm@ni.drops.lut.mean\endcsname{3.89}% datasheet devices: table, mean excess over the optimum, pp
\expandafter\def\csname shm@ni.drops.lut.max\endcsname{63.3}% datasheet devices: table, worst excess over the optimum, pp
\expandafter\def\csname shm@ni.drops.br.mean\endcsname{0.056}% datasheet devices: branch-aware table, mean excess, pp
\expandafter\def\csname shm@ni.drops.br.max\endcsname{2.82}% datasheet devices: branch-aware table, worst excess, pp
\expandafter\def\csname shm@ni.dt.opt.dmax\endcsname{0.455}% + dead time 2 us: optimum, largest |THD49 - closed form|, pp
\expandafter\def\csname shm@ni.dt.opt.b1\endcsname{1.10}% + dead time 2 us: optimum, largest |fundamental error|, %
\expandafter\def\csname shm@ni.dt.lut.mean\endcsname{3.92}% + dead time 2 us: table, mean excess over the optimum, pp
\expandafter\def\csname shm@ni.dt.lut.max\endcsname{63.7}% + dead time 2 us: table, worst excess over the optimum, pp
\expandafter\def\csname shm@ni.dt.br.mean\endcsname{0.057}% + dead time 2 us: branch-aware table, mean excess, pp
\expandafter\def\csname shm@ni.dt.br.max\endcsname{2.77}% + dead time 2 us: branch-aware table, worst excess, pp
\expandafter\def\csname shm@ni.np.opt.dmax\endcsname{0.737}% + NP offset 2%: optimum, largest |THD49 - closed form|, pp
\expandafter\def\csname shm@ni.np.opt.b1\endcsname{0.55}% + NP offset 2%: optimum, largest |fundamental error|, %
\expandafter\def\csname shm@ni.np.lut.mean\endcsname{3.85}% + NP offset 2%: table, mean excess over the optimum, pp
\expandafter\def\csname shm@ni.np.lut.max\endcsname{62.9}% + NP offset 2%: table, worst excess over the optimum, pp
\expandafter\def\csname shm@ni.np.br.mean\endcsname{0.051}% + NP offset 2%: branch-aware table, mean excess, pp
\expandafter\def\csname shm@ni.np.br.max\endcsname{2.65}% + NP offset 2%: branch-aware table, worst excess, pp
\expandafter\def\csname shm@ni.tq.opt.dmax\endcsname{0.173}% + timer 1 us: optimum, largest |THD49 - closed form|, pp
\expandafter\def\csname shm@ni.tq.opt.b1\endcsname{0.61}% + timer 1 us: optimum, largest |fundamental error|, %
\expandafter\def\csname shm@ni.tq.lut.mean\endcsname{3.89}% + timer 1 us: table, mean excess over the optimum, pp
\expandafter\def\csname shm@ni.tq.lut.max\endcsname{63.3}% + timer 1 us: table, worst excess over the optimum, pp
\expandafter\def\csname shm@ni.tq.br.mean\endcsname{0.056}% + timer 1 us: branch-aware table, mean excess, pp
\expandafter\def\csname shm@ni.tq.br.max\endcsname{2.86}% + timer 1 us: branch-aware table, worst excess, pp
\expandafter\def\csname shm@ni.all.opt.dmax\endcsname{0.898}% all combined: optimum, largest |THD49 - closed form|, pp
\expandafter\def\csname shm@ni.all.opt.b1\endcsname{1.13}% all combined: optimum, largest |fundamental error|, %
\expandafter\def\csname shm@ni.all.lut.mean\endcsname{3.87}% all combined: table, mean excess over the optimum, pp
\expandafter\def\csname shm@ni.all.lut.max\endcsname{63.3}% all combined: table, worst excess over the optimum, pp
\expandafter\def\csname shm@ni.all.br.mean\endcsname{0.052}% all combined: branch-aware table, mean excess, pp
\expandafter\def\csname shm@ni.all.br.max\endcsname{2.65}% all combined: branch-aware table, worst excess, pp
\expandafter\def\csname shm@ni.opt.dmax\endcsname{0.898}% largest |THD49 - closed form| of the optimum under any condition, pp
\expandafter\def\csname shm@ni.lut.spread\endcsname{0.07}% spread of the table's mean excess across the conditions, pp
\expandafter\def\csname shm@ni.br.spread\endcsname{0.006}% spread of the branch-aware mean excess across the conditions, pp
\expandafter\def\csname shm@ni.ripple.pp\endcsname{7.3}% largest neutral-point ripple in the switching-function model, V peak to peak
\expandafter\def\csname shm@sim.ni.runs\endcsname{54}% circuit runs with datasheet switches that ran the commanded angles
\expandafter\def\csname shm@sim.ni.invalid\endcsname{0}% runs left out: commanded angles not run (must be 0)
\expandafter\def\csname shm@sim.ni.conds\endcsname{6}% non-ideal conditions simulated
\expandafter\def\csname shm@sim.ni.swfn.max\endcsname{0.013}% circuit minus switching-function model, optimal and branch-aware patterns, largest |difference|, pp
\expandafter\def\csname shm@sim.ni.swfn.lut.rel\endcsname{0.18}% the same for the table's angles, % of their THD49
\expandafter\def\csname shm@sim.ni.offset.max\endcsname{1.49}% largest |neutral-point offset| the circuit had in the analysis window, % of Vdc
\expandafter\def\csname shm@sim.ni.swfn.cmd.max\endcsname{0.606}% circuit minus switching-function model at the COMMANDED offset (not the measured one), largest |difference|, pp
\expandafter\def\csname shm@sim.ni.lut.rows\endcsname{100}% rows of the largest table simulated with datasheet switches
\expandafter\def\csname shm@sim.ni.ideal.opt.dev\endcsname{0.114}% ideal switches: optimal patterns, largest |circuit - closed form| THD49, pp
\expandafter\def\csname shm@sim.ni.ideal.b1\endcsname{0.04}% ideal switches: optimal patterns, largest |fundamental error| in the circuit, %
\expandafter\def\csname shm@sim.ni.ideal.ripple.pp\endcsname{24.3}% ideal switches: largest neutral-point voltage, V peak to peak
\expandafter\def\csname shm@sim.ni.ideal.offset.min\endcsname{-1.49}% ideal switches: smallest measured NP offset, % of Vdc
\expandafter\def\csname shm@sim.ni.ideal.offset.max\endcsname{-0.47}% ideal switches: largest measured NP offset, % of Vdc
\expandafter\def\csname shm@sim.ni.ideal.lut.circuit\endcsname{60.82}% ideal switches: table at its worst point, circuit THD49, %
\expandafter\def\csname shm@sim.ni.ideal.lut.ithd\endcsname{6.49}% ideal switches: the same, phase-current THD49, %
\expandafter\def\csname shm@sim.ni.ideal.branch.circuit\endcsname{1.13}% ideal switches: branch-aware table at that Ma, circuit THD49, %
\expandafter\def\csname shm@sim.ni.ideal.branch.ithd\endcsname{2.44}% ideal switches: the same, phase-current THD49, %
\expandafter\def\csname shm@sim.ni.drops.opt.dev\endcsname{0.149}% datasheet devices: optimal patterns, largest |circuit - closed form| THD49, pp
\expandafter\def\csname shm@sim.ni.drops.b1\endcsname{0.49}% datasheet devices: optimal patterns, largest |fundamental error| in the circuit, %
\expandafter\def\csname shm@sim.ni.drops.ripple.pp\endcsname{23.8}% datasheet devices: largest neutral-point voltage, V peak to peak
\expandafter\def\csname shm@sim.ni.drops.offset.min\endcsname{-1.44}% datasheet devices: smallest measured NP offset, % of Vdc
\expandafter\def\csname shm@sim.ni.drops.offset.max\endcsname{-0.45}% datasheet devices: largest measured NP offset, % of Vdc
\expandafter\def\csname shm@sim.ni.drops.lut.circuit\endcsname{61.14}% datasheet devices: table at its worst point, circuit THD49, %
\expandafter\def\csname shm@sim.ni.drops.lut.ithd\endcsname{6.49}% datasheet devices: the same, phase-current THD49, %
\expandafter\def\csname shm@sim.ni.drops.branch.circuit\endcsname{1.11}% datasheet devices: branch-aware table at that Ma, circuit THD49, %
\expandafter\def\csname shm@sim.ni.drops.branch.ithd\endcsname{2.37}% datasheet devices: the same, phase-current THD49, %
\expandafter\def\csname shm@sim.ni.dt.opt.dev\endcsname{0.180}% + dead time 2 us: optimal patterns, largest |circuit - closed form| THD49, pp
\expandafter\def\csname shm@sim.ni.dt.b1\endcsname{0.95}% + dead time 2 us: optimal patterns, largest |fundamental error| in the circuit, %
\expandafter\def\csname shm@sim.ni.dt.ripple.pp\endcsname{23.6}% + dead time 2 us: largest neutral-point voltage, V peak to peak
\expandafter\def\csname shm@sim.ni.dt.offset.min\endcsname{-1.41}% + dead time 2 us: smallest measured NP offset, % of Vdc
\expandafter\def\csname shm@sim.ni.dt.offset.max\endcsname{-0.45}% + dead time 2 us: largest measured NP offset, % of Vdc
\expandafter\def\csname shm@sim.ni.dt.lut.circuit\endcsname{61.57}% + dead time 2 us: table at its worst point, circuit THD49, %
\expandafter\def\csname shm@sim.ni.dt.lut.ithd\endcsname{6.52}% + dead time 2 us: the same, phase-current THD49, %
\expandafter\def\csname shm@sim.ni.dt.branch.circuit\endcsname{1.17}% + dead time 2 us: branch-aware table at that Ma, circuit THD49, %
\expandafter\def\csname shm@sim.ni.dt.branch.ithd\endcsname{2.36}% + dead time 2 us: the same, phase-current THD49, %
\expandafter\def\csname shm@sim.ni.np.opt.dev\endcsname{0.128}% + NP offset 2%: optimal patterns, largest |circuit - closed form| THD49, pp
\expandafter\def\csname shm@sim.ni.np.b1\endcsname{0.49}% + NP offset 2%: optimal patterns, largest |fundamental error| in the circuit, %
\expandafter\def\csname shm@sim.ni.np.ripple.pp\endcsname{27.2}% + NP offset 2%: largest neutral-point voltage, V peak to peak
\expandafter\def\csname shm@sim.ni.np.offset.min\endcsname{0.03}% + NP offset 2%: smallest measured NP offset, % of Vdc
\expandafter\def\csname shm@sim.ni.np.offset.max\endcsname{1.07}% + NP offset 2%: largest measured NP offset, % of Vdc
\expandafter\def\csname shm@sim.ni.np.lut.circuit\endcsname{61.14}% + NP offset 2%: table at its worst point, circuit THD49, %
\expandafter\def\csname shm@sim.ni.np.lut.ithd\endcsname{6.41}% + NP offset 2%: the same, phase-current THD49, %
\expandafter\def\csname shm@sim.ni.np.branch.circuit\endcsname{1.06}% + NP offset 2%: branch-aware table at that Ma, circuit THD49, %
\expandafter\def\csname shm@sim.ni.np.branch.ithd\endcsname{2.21}% + NP offset 2%: the same, phase-current THD49, %
\expandafter\def\csname shm@sim.ni.tq.opt.dev\endcsname{0.158}% + timer 1 us: optimal patterns, largest |circuit - closed form| THD49, pp
\expandafter\def\csname shm@sim.ni.tq.b1\endcsname{0.47}% + timer 1 us: optimal patterns, largest |fundamental error| in the circuit, %
\expandafter\def\csname shm@sim.ni.tq.ripple.pp\endcsname{23.9}% + timer 1 us: largest neutral-point voltage, V peak to peak
\expandafter\def\csname shm@sim.ni.tq.offset.min\endcsname{-1.44}% + timer 1 us: smallest measured NP offset, % of Vdc
\expandafter\def\csname shm@sim.ni.tq.offset.max\endcsname{-0.45}% + timer 1 us: largest measured NP offset, % of Vdc
\expandafter\def\csname shm@sim.ni.tq.lut.circuit\endcsname{61.17}% + timer 1 us: table at its worst point, circuit THD49, %
\expandafter\def\csname shm@sim.ni.tq.lut.ithd\endcsname{6.49}% + timer 1 us: the same, phase-current THD49, %
\expandafter\def\csname shm@sim.ni.tq.branch.circuit\endcsname{1.12}% + timer 1 us: branch-aware table at that Ma, circuit THD49, %
\expandafter\def\csname shm@sim.ni.tq.branch.ithd\endcsname{2.37}% + timer 1 us: the same, phase-current THD49, %
\expandafter\def\csname shm@sim.ni.all.opt.dev\endcsname{0.168}% all combined: optimal patterns, largest |circuit - closed form| THD49, pp
\expandafter\def\csname shm@sim.ni.all.b1\endcsname{0.93}% all combined: optimal patterns, largest |fundamental error| in the circuit, %
\expandafter\def\csname shm@sim.ni.all.ripple.pp\endcsname{27.0}% all combined: largest neutral-point voltage, V peak to peak
\expandafter\def\csname shm@sim.ni.all.offset.min\endcsname{0.06}% all combined: smallest measured NP offset, % of Vdc
\expandafter\def\csname shm@sim.ni.all.offset.max\endcsname{1.09}% all combined: largest measured NP offset, % of Vdc
\expandafter\def\csname shm@sim.ni.all.lut.circuit\endcsname{61.61}% all combined: table at its worst point, circuit THD49, %
\expandafter\def\csname shm@sim.ni.all.lut.ithd\endcsname{6.45}% all combined: the same, phase-current THD49, %
\expandafter\def\csname shm@sim.ni.all.branch.circuit\endcsname{1.12}% all combined: branch-aware table at that Ma, circuit THD49, %
\expandafter\def\csname shm@sim.ni.all.branch.ithd\endcsname{2.21}% all combined: the same, phase-current THD49, %
\expandafter\def\csname shm@time.total\endcsname{0.00}% hours for the whole run
\makeatother
 in the preamble, then \shm{key} in the text.
\makeatletter
\providecommand{\shm}[1]{\@ifundefined{shm@#1}{\textbf{??#1??}}{\csname shm@#1\endcsname}}
\expandafter\def\csname shm@angles\endcsname{15}% switching angles per quarter cycle
\expandafter\def\csname shm@orders\endcsname{16}% odd non-triplen orders in THD49 (5..49)
\expandafter\def\csname shm@ma.lo\endcsname{0.70}% lowest modulation index of the study
\expandafter\def\csname shm@ma.hi\endcsname{1.20}% highest modulation index of the study
\expandafter\def\csname shm@ma.step\endcsname{0.0025}% modulation-index step of the traced grid
\expandafter\def\csname shm@grid.points\endcsname{201}% grid points along Ma
\expandafter\def\csname shm@anchors\endcsname{11}% multistart anchors
\expandafter\def\csname shm@nstart\endcsname{150}% random starts per anchor
\expandafter\def\csname shm@distinct.mean\endcsname{104.2}% distinct local optima per anchor, mean
\expandafter\def\csname shm@distinct.max\endcsname{112}% distinct local optima at one anchor, max
\expandafter\def\csname shm@hitrate\endcsname{1.1}% % of converged random starts that reach the best optimum
\expandafter\def\csname shm@spread\endcsname{9.41}% median minus best THD49 among the distinct optima, pp
\expandafter\def\csname shm@branches\endcsname{76}% distinct continuous branches traced
\expandafter\def\csname shm@owners\endcsname{17}% branches that own part of the per-Ma optimum
\expandafter\def\csname shm@ownerchanges\endcsname{21}% owner changes of the per-Ma optimum over the grid
\expandafter\def\csname shm@jump.mean\endcsname{27.2}% mean jump of the optimal angle vector at an owner change, deg
\expandafter\def\csname shm@jump.max\endcsname{49.2}% largest jump of the optimal angle vector, deg
\expandafter\def\csname shm@stored\endcsname{13}% branches kept in the stored subset
\expandafter\def\csname shm@stored.gap.mean\endcsname{0.0004}% mean gap of the stored subset to the full envelope, pp
\expandafter\def\csname shm@stored.gap.max\endcsname{0.0282}% worst gap of the stored subset to the full envelope, pp
\expandafter\def\csname shm@coverage\endcsname{100.0}% % of grid points where some branch exists
\expandafter\def\csname shm@env.thd.mean\endcsname{2.13}% mean THD49 of the per-Ma optimum, %
\expandafter\def\csname shm@env.thd.min\endcsname{0.22}% lowest THD49 of the per-Ma optimum, %
\expandafter\def\csname shm@env.thd.max\endcsname{3.86}% highest THD49 of the per-Ma optimum, %
\expandafter\def\csname shm@nfe.branchset\endcsname{17696578}% objective evaluations to build the branch set
\expandafter\def\csname shm@time.branchset\endcsname{29.3}% minutes to build the branch set
\expandafter\def\csname shm@effangles.min\endcsname{12.0}% fewest effective angles on any branch (collapsed pairs removed)
\expandafter\def\csname shm@traces\endcsname{91}% continuation traces before merging (seeds and probes)
\expandafter\def\csname shm@probes\endcsname{50}% modulation indices of the probe multistart
\expandafter\def\csname shm@probe.nstart\endcsname{100}% random starts per probe
\expandafter\def\csname shm@probe.traces\endcsname{11}% traces added because a probe beat the envelope
\expandafter\def\csname shm@probe.hits\endcsname{2}% probes that found an optimum below the envelope
\expandafter\def\csname shm@fold.points\endcsname{24}% grid points searched just beyond the ends of owning branches
\expandafter\def\csname shm@fold.traces\endcsname{2}% traces added by that fold-end search
\expandafter\def\csname shm@dropped\endcsname{0}% traces of one or two grid points that own nothing, dropped
\expandafter\def\csname shm@distinct.mean.0p5\endcsname{102.5}% distinct optima per anchor, clustered at 0.5 deg, mean
\expandafter\def\csname shm@distinct.mean.2\endcsname{99.5}% distinct optima per anchor, clustered at 2 deg, mean
\expandafter\def\csname shm@distinct.levels\endcsname{95.6}% distinct THD49 levels (1e-3 pp apart) per anchor, mean
\expandafter\def\csname shm@distinct.within1pp\endcsname{10.4}% distinct optima within 1 pp of the best, per anchor, mean
\expandafter\def\csname shm@jump.wave.mean\endcsname{15.1}% share of the period in which the optimal voltage changes at an owner change, mean, %
\expandafter\def\csname shm@jump.wave.max\endcsname{48.8}% share of the period in which the optimal voltage changes at an owner change, largest, %
\expandafter\def\csname shm@jump.wave.min\endcsname{4.3}% share of the period in which the optimal voltage changes at an owner change, smallest, %
\expandafter\def\csname shm@jump.wave.at.maxangle\endcsname{4.3}% the same at the owner change with the largest angle move, %
\expandafter\def\csname shm@distinct.once\endcsname{83}% share of the distinct optima reached by a single random start, per anchor, mean, %
\expandafter\def\csname shm@distinct.chao1\endcsname{480}% Chao1 estimate of the total number of local optima per anchor, mean (the distinct count is a lower bound)
\expandafter\def\csname shm@lut.excess.R10\endcsname{21.9503}% LUT at R = 10: mean THD49 excess over the optimum, pp
\expandafter\def\csname shm@lut.excessmax.R10\endcsname{63.309}% LUT at R = 10: worst THD49 excess, pp
\expandafter\def\csname shm@lut.angmax.R10\endcsname{41.19}% LUT at R = 10: worst angle error, deg
\expandafter\def\csname shm@lut.fund.R10\endcsname{0.185}% LUT at R = 10: mean fundamental error, %
\expandafter\def\csname shm@lut.excess.R25\endcsname{13.3422}% LUT at R = 25: mean THD49 excess over the optimum, pp
\expandafter\def\csname shm@lut.excessmax.R25\endcsname{62.555}% LUT at R = 25: worst THD49 excess, pp
\expandafter\def\csname shm@lut.angmax.R25\endcsname{37.65}% LUT at R = 25: worst angle error, deg
\expandafter\def\csname shm@lut.fund.R25\endcsname{0.108}% LUT at R = 25: mean fundamental error, %
\expandafter\def\csname shm@lut.excess.R50\endcsname{7.6717}% LUT at R = 50: mean THD49 excess over the optimum, pp
\expandafter\def\csname shm@lut.excessmax.R50\endcsname{62.767}% LUT at R = 50: worst THD49 excess, pp
\expandafter\def\csname shm@lut.angmax.R50\endcsname{39.05}% LUT at R = 50: worst angle error, deg
\expandafter\def\csname shm@lut.fund.R50\endcsname{0.073}% LUT at R = 50: mean fundamental error, %
\expandafter\def\csname shm@lut.excess.R100\endcsname{3.7973}% LUT at R = 100: mean THD49 excess over the optimum, pp
\expandafter\def\csname shm@lut.excessmax.R100\endcsname{59.862}% LUT at R = 100: worst THD49 excess, pp
\expandafter\def\csname shm@lut.angmax.R100\endcsname{48.42}% LUT at R = 100: worst angle error, deg
\expandafter\def\csname shm@lut.fund.R100\endcsname{0.044}% LUT at R = 100: mean fundamental error, %
\expandafter\def\csname shm@nearest.excess.R10\endcsname{0.0610}% NEAREST at R = 10: mean THD49 excess over the optimum, pp
\expandafter\def\csname shm@nearest.excessmax.R10\endcsname{1.923}% NEAREST at R = 10: worst THD49 excess, pp
\expandafter\def\csname shm@nearest.angmax.R10\endcsname{52.70}% NEAREST at R = 10: worst angle error, deg
\expandafter\def\csname shm@nearest.fund.R10\endcsname{1.567}% NEAREST at R = 10: mean fundamental error, %
\expandafter\def\csname shm@nearest.excess.R25\endcsname{0.0100}% NEAREST at R = 25: mean THD49 excess over the optimum, pp
\expandafter\def\csname shm@nearest.excessmax.R25\endcsname{0.665}% NEAREST at R = 25: worst THD49 excess, pp
\expandafter\def\csname shm@nearest.angmax.R25\endcsname{40.62}% NEAREST at R = 25: worst angle error, deg
\expandafter\def\csname shm@nearest.fund.R25\endcsname{0.637}% NEAREST at R = 25: mean fundamental error, %
\expandafter\def\csname shm@nearest.excess.R50\endcsname{0.0005}% NEAREST at R = 50: mean THD49 excess over the optimum, pp
\expandafter\def\csname shm@nearest.excessmax.R50\endcsname{0.452}% NEAREST at R = 50: worst THD49 excess, pp
\expandafter\def\csname shm@nearest.angmax.R50\endcsname{49.20}% NEAREST at R = 50: worst angle error, deg
\expandafter\def\csname shm@nearest.fund.R50\endcsname{0.364}% NEAREST at R = 50: mean fundamental error, %
\expandafter\def\csname shm@nearest.excess.R100\endcsname{-0.0082}% NEAREST at R = 100: mean THD49 excess over the optimum, pp
\expandafter\def\csname shm@nearest.excessmax.R100\endcsname{0.150}% NEAREST at R = 100: worst THD49 excess, pp
\expandafter\def\csname shm@nearest.angmax.R100\endcsname{47.63}% NEAREST at R = 100: worst angle error, deg
\expandafter\def\csname shm@nearest.fund.R100\endcsname{0.269}% NEAREST at R = 100: mean fundamental error, %
\expandafter\def\csname shm@sqp.excess.R10\endcsname{4.7744}% LUT+SQP at R = 10: mean THD49 excess over the optimum, pp
\expandafter\def\csname shm@sqp.excessmax.R10\endcsname{31.525}% LUT+SQP at R = 10: worst THD49 excess, pp
\expandafter\def\csname shm@sqp.angmax.R10\endcsname{78.09}% LUT+SQP at R = 10: worst angle error, deg
\expandafter\def\csname shm@sqp.fund.R10\endcsname{0.003}% LUT+SQP at R = 10: mean fundamental error, %
\expandafter\def\csname shm@sqp.excess.R25\endcsname{2.4875}% LUT+SQP at R = 25: mean THD49 excess over the optimum, pp
\expandafter\def\csname shm@sqp.excessmax.R25\endcsname{30.501}% LUT+SQP at R = 25: worst THD49 excess, pp
\expandafter\def\csname shm@sqp.angmax.R25\endcsname{78.48}% LUT+SQP at R = 25: worst angle error, deg
\expandafter\def\csname shm@sqp.fund.R25\endcsname{0.000}% LUT+SQP at R = 25: mean fundamental error, %
\expandafter\def\csname shm@sqp.excess.R50\endcsname{1.3759}% LUT+SQP at R = 50: mean THD49 excess over the optimum, pp
\expandafter\def\csname shm@sqp.excessmax.R50\endcsname{29.799}% LUT+SQP at R = 50: worst THD49 excess, pp
\expandafter\def\csname shm@sqp.angmax.R50\endcsname{76.64}% LUT+SQP at R = 50: worst angle error, deg
\expandafter\def\csname shm@sqp.fund.R50\endcsname{0.000}% LUT+SQP at R = 50: mean fundamental error, %
\expandafter\def\csname shm@sqp.excess.R100\endcsname{0.5158}% LUT+SQP at R = 100: mean THD49 excess over the optimum, pp
\expandafter\def\csname shm@sqp.excessmax.R100\endcsname{29.182}% LUT+SQP at R = 100: worst THD49 excess, pp
\expandafter\def\csname shm@sqp.angmax.R100\endcsname{76.44}% LUT+SQP at R = 100: worst angle error, deg
\expandafter\def\csname shm@sqp.fund.R100\endcsname{0.000}% LUT+SQP at R = 100: mean fundamental error, %
\expandafter\def\csname shm@cont.excess.R10\endcsname{2.9051}% CONT at R = 10: mean THD49 excess over the optimum, pp
\expandafter\def\csname shm@cont.excessmax.R10\endcsname{7.882}% CONT at R = 10: worst THD49 excess, pp
\expandafter\def\csname shm@cont.angmax.R10\endcsname{49.71}% CONT at R = 10: worst angle error, deg
\expandafter\def\csname shm@cont.fund.R10\endcsname{0.040}% CONT at R = 10: mean fundamental error, %
\expandafter\def\csname shm@cont.excess.R25\endcsname{1.7159}% CONT at R = 25: mean THD49 excess over the optimum, pp
\expandafter\def\csname shm@cont.excessmax.R25\endcsname{6.730}% CONT at R = 25: worst THD49 excess, pp
\expandafter\def\csname shm@cont.angmax.R25\endcsname{48.59}% CONT at R = 25: worst angle error, deg
\expandafter\def\csname shm@cont.fund.R25\endcsname{0.007}% CONT at R = 25: mean fundamental error, %
\expandafter\def\csname shm@cont.excess.R50\endcsname{1.3021}% CONT at R = 50: mean THD49 excess over the optimum, pp
\expandafter\def\csname shm@cont.excessmax.R50\endcsname{5.197}% CONT at R = 50: worst THD49 excess, pp
\expandafter\def\csname shm@cont.angmax.R50\endcsname{48.43}% CONT at R = 50: worst angle error, deg
\expandafter\def\csname shm@cont.fund.R50\endcsname{0.003}% CONT at R = 50: mean fundamental error, %
\expandafter\def\csname shm@cont.excess.R100\endcsname{1.1490}% CONT at R = 100: mean THD49 excess over the optimum, pp
\expandafter\def\csname shm@cont.excessmax.R100\endcsname{3.979}% CONT at R = 100: worst THD49 excess, pp
\expandafter\def\csname shm@cont.angmax.R100\endcsname{48.42}% CONT at R = 100: worst angle error, deg
\expandafter\def\csname shm@cont.fund.R100\endcsname{0.002}% CONT at R = 100: mean fundamental error, %
\expandafter\def\csname shm@branch.excess.R10\endcsname{1.2992}% BRANCH at R = 10: mean THD49 excess over the optimum, pp
\expandafter\def\csname shm@branch.excessmax.R10\endcsname{13.031}% BRANCH at R = 10: worst THD49 excess, pp
\expandafter\def\csname shm@branch.angmax.R10\endcsname{52.62}% BRANCH at R = 10: worst angle error, deg
\expandafter\def\csname shm@branch.fund.R10\endcsname{0.004}% BRANCH at R = 10: mean fundamental error, %
\expandafter\def\csname shm@branch.excess.R25\endcsname{0.3848}% BRANCH at R = 25: mean THD49 excess over the optimum, pp
\expandafter\def\csname shm@branch.excessmax.R25\endcsname{10.243}% BRANCH at R = 25: worst THD49 excess, pp
\expandafter\def\csname shm@branch.angmax.R25\endcsname{52.62}% BRANCH at R = 25: worst angle error, deg
\expandafter\def\csname shm@branch.fund.R25\endcsname{0.018}% BRANCH at R = 25: mean fundamental error, %
\expandafter\def\csname shm@branch.excess.R50\endcsname{0.1626}% BRANCH at R = 50: mean THD49 excess over the optimum, pp
\expandafter\def\csname shm@branch.excessmax.R50\endcsname{6.860}% BRANCH at R = 50: worst THD49 excess, pp
\expandafter\def\csname shm@branch.angmax.R50\endcsname{52.62}% BRANCH at R = 50: worst angle error, deg
\expandafter\def\csname shm@branch.fund.R50\endcsname{0.003}% BRANCH at R = 50: mean fundamental error, %
\expandafter\def\csname shm@branch.excess.R100\endcsname{0.0706}% BRANCH at R = 100: mean THD49 excess over the optimum, pp
\expandafter\def\csname shm@branch.excessmax.R100\endcsname{2.023}% BRANCH at R = 100: worst THD49 excess, pp
\expandafter\def\csname shm@branch.angmax.R100\endcsname{50.94}% BRANCH at R = 100: worst angle error, deg
\expandafter\def\csname shm@branch.fund.R100\endcsname{0.003}% BRANCH at R = 100: mean fundamental error, %
\expandafter\def\csname shm@lut.floor.ratio\endcsname{0.17}% LUT mean excess at R = 100 divided by that at R = 10 (1 = no improvement at all)
\expandafter\def\csname shm@lut.worst.ratio\endcsname{0.95}% LUT worst excess at R = 100 divided by that at R = 10
\expandafter\def\csname shm@lut.slope\endcsname{-0.76}% log-log slope of the LUT mean excess vs R (resolution error alone: -2)
\expandafter\def\csname shm@branch.slope\endcsname{-2.00}% log-log slope of the branch-aware mean excess vs stored sets
\expandafter\def\csname shm@split.slope.same\endcsname{-2.78}% log-log slope of the same-owner (resolution) error vs R, median |excess|
\expandafter\def\csname shm@split.slope.straddle\endcsname{-0.05}% log-log slope of the straddling (structural) error vs R, median |excess|
\expandafter\def\csname shm@split.slope.same.mean\endcsname{-1.74}% log-log slope of the same-owner error vs R, mean
\expandafter\def\csname shm@split.slope.straddle.mean\endcsname{-0.03}% log-log slope of the straddling error vs R, mean
\expandafter\def\csname shm@split.frac.R10\endcsname{89.0}% % of test points whose bracketing rows straddle an owner change, R = 10
\expandafter\def\csname shm@split.straddle.R10\endcsname{24.661}% mean excess of straddling intervals, R = 10, pp
\expandafter\def\csname shm@split.same.R10\endcsname{0.00210}% mean excess of same-owner intervals, R = 10, pp
\expandafter\def\csname shm@split.straddle.median.R10\endcsname{19.070}% median |excess| of straddling intervals, R = 10, pp
\expandafter\def\csname shm@split.same.median.R10\endcsname{2.19e-03}% median |excess| of same-owner intervals, R = 10, pp
\expandafter\def\csname shm@split.same.p90.R10\endcsname{3.76e-03}% 90th percentile |excess| of same-owner intervals, R = 10, pp
\expandafter\def\csname shm@split.frac.R25\endcsname{46.0}% % of test points whose bracketing rows straddle an owner change, R = 25
\expandafter\def\csname shm@split.straddle.R25\endcsname{28.989}% mean excess of straddling intervals, R = 25, pp
\expandafter\def\csname shm@split.same.R25\endcsname{0.00123}% mean excess of same-owner intervals, R = 25, pp
\expandafter\def\csname shm@split.straddle.median.R25\endcsname{27.360}% median |excess| of straddling intervals, R = 25, pp
\expandafter\def\csname shm@split.same.median.R25\endcsname{3.63e-05}% median |excess| of same-owner intervals, R = 25, pp
\expandafter\def\csname shm@split.same.p90.R25\endcsname{1.95e-03}% 90th percentile |excess| of same-owner intervals, R = 25, pp
\expandafter\def\csname shm@split.frac.R50\endcsname{27.8}% % of test points whose bracketing rows straddle an owner change, R = 50
\expandafter\def\csname shm@split.straddle.R50\endcsname{27.581}% mean excess of straddling intervals, R = 50, pp
\expandafter\def\csname shm@split.same.R50\endcsname{0.00017}% mean excess of same-owner intervals, R = 50, pp
\expandafter\def\csname shm@split.straddle.median.R50\endcsname{24.542}% median |excess| of straddling intervals, R = 50, pp
\expandafter\def\csname shm@split.same.median.R50\endcsname{1.09e-05}% median |excess| of same-owner intervals, R = 50, pp
\expandafter\def\csname shm@split.same.p90.R50\endcsname{9.76e-05}% 90th percentile |excess| of same-owner intervals, R = 50, pp
\expandafter\def\csname shm@split.frac.R100\endcsname{16.8}% % of test points whose bracketing rows straddle an owner change, R = 100
\expandafter\def\csname shm@split.straddle.R100\endcsname{22.560}% mean excess of straddling intervals, R = 100, pp
\expandafter\def\csname shm@split.same.R100\endcsname{0.00005}% mean excess of same-owner intervals, R = 100, pp
\expandafter\def\csname shm@split.straddle.median.R100\endcsname{16.736}% median |excess| of straddling intervals, R = 100, pp
\expandafter\def\csname shm@split.same.median.R100\endcsname{3.22e-06}% median |excess| of same-owner intervals, R = 100, pp
\expandafter\def\csname shm@split.same.p90.R100\endcsname{1.29e-05}% 90th percentile |excess| of same-owner intervals, R = 100, pp
\expandafter\def\csname shm@split.straddle.Rmax\endcsname{22.560}% straddling-interval excess at R = 100, pp
\expandafter\def\csname shm@split.same.Rmax\endcsname{0.00005}% same-owner-interval excess at R = 100, pp
\expandafter\def\csname shm@basin.onowner\endcsname{61.4}% % of corrected points (R = 25) that land on the owning branch
\expandafter\def\csname shm@basin.onowner.near\endcsname{14.3}% % landing on the owning branch within two grid steps of an owner change
\expandafter\def\csname shm@basin.onowner.far\endcsname{79.5}% % landing on the owning branch farther from an owner change
\expandafter\def\csname shm@basin.excess.wrong\endcsname{6.438}% mean THD49 excess where the corrector landed on another branch, pp
\expandafter\def\csname shm@basin.ok.near\endcsname{28.6}% % of corrected points within 0.01 pp of the optimum, within two grid steps of an owner change
\expandafter\def\csname shm@basin.ok.far\endcsname{79.5}% % of corrected points within 0.01 pp of the optimum, farther from an owner change
\expandafter\def\csname shm@cont.cost.mean\endcsname{0.999}% mean THD49 cost of the continuity-constrained path, pp
\expandafter\def\csname shm@cont.cost.max\endcsname{2.282}% worst THD49 cost of the continuity-constrained path, pp
\expandafter\def\csname shm@cont.breaks\endcsname{0}% places where even the continuity-constrained path must jump
\expandafter\def\csname shm@cont.jumps.env\endcsname{20}% jumps (> 5 deg) the per-Ma optimum makes over the grid
\expandafter\def\csname shm@cont.jumps.zero\endcsname{11}% fewest jumps of a path whose mean THD cost is below 1e-4 pp (the rest only swap near-equal branches)
\expandafter\def\csname shm@cont.jumps.for.1\endcsname{0}% fewest jumps that keep the mean THD49 cost at or below 1 pp
\expandafter\def\csname shm@cont.jumps.for.0p1\endcsname{5}% fewest jumps that keep the mean THD49 cost at or below 0.1 pp
\expandafter\def\csname shm@cont.jumps.for.0p01\endcsname{8}% fewest jumps that keep the mean THD49 cost at or below 0.01 pp
\expandafter\def\csname shm@cont.cost.k1\endcsname{0.688}% mean THD49 cost of the best path allowed one jump, pp
\expandafter\def\csname shm@reg1.linear.K41\endcsname{0.0101}% 1-D regression, K = 41: angle MAE of linear, deg
\expandafter\def\csname shm@reg1.linear.K6\endcsname{0.1820}% 1-D regression, K = 6: angle MAE of linear, deg
\expandafter\def\csname shm@reg1.spline.K41\endcsname{0.0077}% 1-D regression, K = 41: angle MAE of spline, deg
\expandafter\def\csname shm@reg1.spline.K6\endcsname{0.1351}% 1-D regression, K = 6: angle MAE of spline, deg
\expandafter\def\csname shm@reg1.pchip.K41\endcsname{0.0068}% 1-D regression, K = 41: angle MAE of pchip, deg
\expandafter\def\csname shm@reg1.pchip.K6\endcsname{0.1359}% 1-D regression, K = 6: angle MAE of pchip, deg
\expandafter\def\csname shm@reg1.RBF.K41\endcsname{0.0092}% 1-D regression, K = 41: angle MAE of RBF, deg
\expandafter\def\csname shm@reg1.RBF.K6\endcsname{0.1741}% 1-D regression, K = 6: angle MAE of RBF, deg
\expandafter\def\csname shm@reg1.GP.K41\endcsname{0.0114}% 1-D regression, K = 41: angle MAE of GP, deg
\expandafter\def\csname shm@reg1.GP.K6\endcsname{0.2451}% 1-D regression, K = 6: angle MAE of GP, deg
\expandafter\def\csname shm@reg1.tanhLM.K41\endcsname{0.0119}% 1-D regression, K = 41: angle MAE of tanhLM, deg
\expandafter\def\csname shm@reg1.tanhLM.K6\endcsname{0.1463}% 1-D regression, K = 6: angle MAE of tanhLM, deg
\expandafter\def\csname shm@reg1.MLP.K41\endcsname{0.0441}% 1-D regression, K = 41: angle MAE of MLP, deg
\expandafter\def\csname shm@reg1.MLP.K6\endcsname{0.5077}% 1-D regression, K = 6: angle MAE of MLP, deg
\expandafter\def\csname shm@reg1.ANFIS.K41\endcsname{0.0140}% 1-D regression, K = 41: angle MAE of ANFIS, deg
\expandafter\def\csname shm@reg1.ANFIS.K6\endcsname{0.3298}% 1-D regression, K = 6: angle MAE of ANFIS, deg
\expandafter\def\csname shm@reg1.mlp.ratio\endcsname{6.4}% MLP angle MAE divided by the best interpolant (pchip) at K = 41
\expandafter\def\csname shm@reg1.mlp.seedcv\endcsname{23}% seed-to-seed coefficient of variation of the MLP MAE, %
\expandafter\def\csname shm@reg2.linear\endcsname{0.0555}% 2-D regression, largest lattice (64 points): angle MAE of linear, deg
\expandafter\def\csname shm@reg2.spline\endcsname{0.0384}% 2-D regression, largest lattice (64 points): angle MAE of spline, deg
\expandafter\def\csname shm@reg2.RBF\endcsname{0.0468}% 2-D regression, largest lattice (64 points): angle MAE of RBF, deg
\expandafter\def\csname shm@reg2.GP\endcsname{0.0598}% 2-D regression, largest lattice (64 points): angle MAE of GP, deg
\expandafter\def\csname shm@reg2.tanhLM\endcsname{0.0485}% 2-D regression, largest lattice (64 points): angle MAE of tanhLM, deg
\expandafter\def\csname shm@reg2.MLP\endcsname{0.0774}% 2-D regression, largest lattice (64 points): angle MAE of MLP, deg
\expandafter\def\csname shm@reg2.ANFIS\endcsname{0.0553}% 2-D regression, largest lattice (64 points): angle MAE of ANFIS, deg
\expandafter\def\csname shm@reg3.linear\endcsname{0.0189}% 3-D regression, largest lattice (216 points): angle MAE of linear, deg
\expandafter\def\csname shm@reg3.spline\endcsname{0.0120}% 3-D regression, largest lattice (216 points): angle MAE of spline, deg
\expandafter\def\csname shm@reg3.RBF\endcsname{0.0119}% 3-D regression, largest lattice (216 points): angle MAE of RBF, deg
\expandafter\def\csname shm@reg3.GP\endcsname{0.0208}% 3-D regression, largest lattice (216 points): angle MAE of GP, deg
\expandafter\def\csname shm@reg3.tanhLM\endcsname{0.0246}% 3-D regression, largest lattice (216 points): angle MAE of tanhLM, deg
\expandafter\def\csname shm@reg3.MLP\endcsname{0.0197}% 3-D regression, largest lattice (216 points): angle MAE of MLP, deg
\expandafter\def\csname shm@reg3.ANFIS\endcsname{0.0265}% 3-D regression, largest lattice (216 points): angle MAE of ANFIS, deg
\expandafter\def\csname shm@dstar.sel.K9\endcsname{4.33}% dimension from which a 9-per-axis label table outgrows the classifier
\expandafter\def\csname shm@dstar.reg.K9\endcsname{3.25}% dimension from which a 9-per-axis angle table outgrows one MLP per branch
\expandafter\def\csname shm@dstar.sel.K17\endcsname{3.36}% dimension from which a 17-per-axis label table outgrows the classifier
\expandafter\def\csname shm@dstar.reg.K17\endcsname{2.52}% dimension from which a 17-per-axis angle table outgrows one MLP per branch
\expandafter\def\csname shm@dstar.sel.K25\endcsname{2.95}% dimension from which a 25-per-axis label table outgrows the classifier
\expandafter\def\csname shm@dstar.reg.K25\endcsname{2.22}% dimension from which a 25-per-axis angle table outgrows one MLP per branch
\expandafter\def\csname shm@t.lut\endcsname{34.4}% table lookup + linear interpolation, microseconds per call
\expandafter\def\csname shm@t.pchip\endcsname{71.5}% pchip evaluation inside a branch, microseconds per call
\expandafter\def\csname shm@t.cls\endcsname{86.4}% softmax classifier forward pass, microseconds per call
\expandafter\def\csname shm@t.mlp\endcsname{267.8}% MLP forward pass, microseconds per call
\expandafter\def\csname shm@t.sqp\endcsname{30771.7}% SQP corrector (fmincon, warm start), microseconds per call
\expandafter\def\csname shm@dense.points\endcsname{61}% off-grid points of the independent multistart check
\expandafter\def\csname shm@dense.nstart\endcsname{200}% fresh random starts per off-grid point
\expandafter\def\csname shm@dense.gap.mean\endcsname{-0.4249}% stored branches minus independent multistart, mean, pp
\expandafter\def\csname shm@dense.gap.worst\endcsname{0.0022}% stored branches minus independent multistart, worst, pp
\expandafter\def\csname shm@dense.better\endcsname{32}% off-grid points where the stored branches beat the multistart
\expandafter\def\csname shm@dense.over1\endcsname{0}% off-grid points where the stored branches lose by more than 1 pp
\expandafter\def\csname shm@dense.over0p1\endcsname{0}% off-grid points where the stored branches lose by more than 0.1 pp
\expandafter\def\csname shm@dense.gap.worst.interp\endcsname{0.0268}% the same without the corrector: the table's own interpolation, worst, pp
\expandafter\def\csname shm@dense.better.max\endcsname{1.76}% largest margin by which the stored branches beat the fresh multistart, pp
\expandafter\def\csname shm@dense.interp.used\endcsname{0}% off-grid points whose stored value is a branch's interpolated angles (its corrector failed or left it)
\expandafter\def\csname shm@dense.corr.fail\endcsname{1}% off-grid points where the corrector of the best interpolated branch did not converge or left the branch
\expandafter\def\csname shm@spec.ma\endcsname{1.00}% Ma of the spectrum audit quoted
\expandafter\def\csname shm@spec.thd49\endcsname{2.44}% THD49 at that Ma, %
\expandafter\def\csname shm@spec.thd99\endcsname{25.7}% THD to the 99th at that Ma, %
\expandafter\def\csname shm@spec.thdfull\endcsname{34.2}% THD to the 999th at that Ma, %
\expandafter\def\csname shm@spec.wthd\endcsname{0.499}% weighted THD (1/n) to the 999th at that Ma, %
\expandafter\def\csname shm@spec.oob\endcsname{17.1}% largest single harmonic above the 49th, % of b1
\expandafter\def\csname shm@spec.oob.order\endcsname{55}% its order
\expandafter\def\csname shm@comply\endcsname{89.6}% % of the envelope meeting the per-order limits (0.2+32.5/n above the 25th)
\expandafter\def\csname shm@comply.alt\endcsname{35.8}% % meeting them with 0.2+12.5/n above the 25th
\expandafter\def\csname shm@bound.ub.model\endcsname{89.875}% the model's upper angle bound, 90 - theta_gap/2, deg
\expandafter\def\csname shm@bound.ub.min\endcsname{80}% tightest upper angle bound tested, deg
\expandafter\def\csname shm@bound.cost.mean\endcsname{0.825}% mean rise of the best THD49 from the model's bound to the tightest one, pp
\expandafter\def\csname shm@bound.cost.max\endcsname{1.688}% largest rise of the best THD49 from the model's bound to the tightest one, pp
\expandafter\def\csname shm@bound.cost89.max\endcsname{1.030}% largest rise of the best THD49 if the angles were kept below 89 deg, pp
\expandafter\def\csname shm@sw.dev\endcsname{1.6e-15}% worst relative deviation, closed form vs exact Fourier integral of the gate pattern
\expandafter\def\csname shm@sw.fftdev\endcsname{0.0096}% worst |THD49 sampled FFT - model| at zero dead time, pp
\expandafter\def\csname shm@sw.dthd.2us\endcsname{0.092}% largest change of THD49 of the optimal patterns caused by a 2 us dead time, pp
\expandafter\def\csname shm@sw.dthd.4us\endcsname{0.089}% largest change of THD49 of the optimal patterns caused by a 4 us dead time, pp
\expandafter\def\csname shm@sw.loss.1us\endcsname{0.0}% % of the envelope whose narrowest pulse is shorter than a 1 us dead time
\expandafter\def\csname shm@sw.loss.2us\endcsname{0.0}% % of the envelope whose narrowest pulse is shorter than a 2 us dead time
\expandafter\def\csname shm@sw.loss.4us\endcsname{0.0}% % of the envelope whose narrowest pulse is shorter than a 4 us dead time
\expandafter\def\csname shm@sw.loss.8us\endcsname{0.0}% % of the envelope whose narrowest pulse is shorter than a 8 us dead time
\expandafter\def\csname shm@sim.cases\endcsname{9}% steady-state cases whose circuit ran the commanded angles
\expandafter\def\csname shm@sim.invalid\endcsname{0}% cases left out: the circuit did not run the commanded angles (must be 0)
\expandafter\def\csname shm@sim.vdc\endcsname{800}% DC-link voltage of the circuit model, V
\expandafter\def\csname shm@sim.ts.us\endcsname{1}% time step of the logged line-to-line voltage, us
\expandafter\def\csname shm@sim.opt.dev.max\endcsname{0.098}% optimal patterns: largest |circuit - analytic| THD49, pp
\expandafter\def\csname shm@sim.opt.samp.max\endcsname{0.004}% optimal patterns: largest |sampled ideal - analytic| THD49 (the time step), pp
\expandafter\def\csname shm@sim.opt.circ.max\endcsname{0.101}% optimal patterns: largest |circuit - sampled ideal| THD49 (what the circuit adds), pp
\expandafter\def\csname shm@sim.opt.ma.dev\endcsname{0.23}% optimal patterns: largest relative error of the circuit's fundamental, %
\expandafter\def\csname shm@sim.lut.rows\endcsname{100}% rows of the largest table simulated at its worst point
\expandafter\def\csname shm@sim.lut.ma\endcsname{0.7925}% Ma of that worst point
\expandafter\def\csname shm@sim.lut.model\endcsname{60.90}% that table's interpolated angles: analytic THD49, %
\expandafter\def\csname shm@sim.lut.circuit\endcsname{60.92}% that table's interpolated angles: circuit THD49, %
\expandafter\def\csname shm@sim.lut.ithd\endcsname{5.53}% that table's interpolated angles: phase-current THD49, %
\expandafter\def\csname shm@sim.branch.model\endcsname{1.04}% branch-aware table at the same Ma: analytic THD49, %
\expandafter\def\csname shm@sim.branch.circuit\endcsname{0.98}% branch-aware table at the same Ma: circuit THD49, %
\expandafter\def\csname shm@sim.branch.ithd\endcsname{0.08}% branch-aware table at the same Ma: phase-current THD49, %
\expandafter\def\csname shm@sim.np.pp.max\endcsname{13.6}% largest peak-to-peak NP voltage over the steady-state windows, V
\expandafter\def\csname shm@sim.step.wave.pct\endcsname{48.8}% branch switch (the owner change that changes the waveform most): share of the period that changes, %
\expandafter\def\csname shm@sim.step.jump.deg\endcsname{30.7}% branch switch: largest angle move at the switch, deg
\expandafter\def\csname shm@sim.step.peak.jump\endcsname{0.37}% branch switch: peak phase-current deviation from the new steady state, % of I1
\expandafter\def\csname shm@sim.step.pre.jump\endcsname{2.38}% branch switch: how far apart the two steady-state currents are, % of I1
\expandafter\def\csname shm@sim.step.offset.jump\endcsname{0.05}% branch switch: mean current offset over the first cycle after the switch, % of I1
\expandafter\def\csname shm@sim.step.peak.stay\endcsname{0.22}% the same Ma step on the old branch (no branch switch): peak phase-current deviation, % of I1
\expandafter\def\csname shm@sim.step.angle.deg\endcsname{49.2}% second branch switch, at the largest angle move: that move, deg
\expandafter\def\csname shm@sim.step.angle.wave.pct\endcsname{4.3}% second branch switch: share of the period that changes, %
\expandafter\def\csname shm@sim.step.angle.peak\endcsname{0.18}% second branch switch: peak phase-current deviation from the new steady state, % of I1
\expandafter\def\csname shm@sim.step.angle.pre\endcsname{0.92}% second branch switch: how far apart the two steady-state currents are, % of I1
\expandafter\def\csname shm@ni.R\endcsname{100}% rows (sets) of the tables evaluated under non-ideal switching
\expandafter\def\csname shm@ni.conds\endcsname{6}% non-ideal conditions evaluated over the whole range
\expandafter\def\csname shm@ni.ideal.opt.dmax\endcsname{0.100}% ideal switches: optimum, largest |THD49 - closed form|, pp
\expandafter\def\csname shm@ni.ideal.opt.b1\endcsname{0.06}% ideal switches: optimum, largest |fundamental error|, %
\expandafter\def\csname shm@ni.ideal.lut.mean\endcsname{3.87}% ideal switches: table, mean excess over the optimum, pp
\expandafter\def\csname shm@ni.ideal.lut.max\endcsname{62.9}% ideal switches: table, worst excess over the optimum, pp
\expandafter\def\csname shm@ni.ideal.br.mean\endcsname{0.057}% ideal switches: branch-aware table, mean excess, pp
\expandafter\def\csname shm@ni.ideal.br.max\endcsname{2.85}% ideal switches: branch-aware table, worst excess, pp
\expandafter\def\csname shm@ni.drops.opt.dmax\endcsname{0.153}% datasheet devices: optimum, largest |THD49 - closed form|, pp
\expandafter\def\csname shm@ni.drops.opt.b1\endcsname{0.55}% datasheet devices: optimum, largest |fundamental error|, %
\expandafter\def\csname shm@ni.drops.lut.mean\endcsname{3.89}% datasheet devices: table, mean excess over the optimum, pp
\expandafter\def\csname shm@ni.drops.lut.max\endcsname{63.3}% datasheet devices: table, worst excess over the optimum, pp
\expandafter\def\csname shm@ni.drops.br.mean\endcsname{0.056}% datasheet devices: branch-aware table, mean excess, pp
\expandafter\def\csname shm@ni.drops.br.max\endcsname{2.82}% datasheet devices: branch-aware table, worst excess, pp
\expandafter\def\csname shm@ni.dt.opt.dmax\endcsname{0.455}% + dead time 2 us: optimum, largest |THD49 - closed form|, pp
\expandafter\def\csname shm@ni.dt.opt.b1\endcsname{1.10}% + dead time 2 us: optimum, largest |fundamental error|, %
\expandafter\def\csname shm@ni.dt.lut.mean\endcsname{3.92}% + dead time 2 us: table, mean excess over the optimum, pp
\expandafter\def\csname shm@ni.dt.lut.max\endcsname{63.7}% + dead time 2 us: table, worst excess over the optimum, pp
\expandafter\def\csname shm@ni.dt.br.mean\endcsname{0.057}% + dead time 2 us: branch-aware table, mean excess, pp
\expandafter\def\csname shm@ni.dt.br.max\endcsname{2.77}% + dead time 2 us: branch-aware table, worst excess, pp
\expandafter\def\csname shm@ni.np.opt.dmax\endcsname{0.737}% + NP offset 2%: optimum, largest |THD49 - closed form|, pp
\expandafter\def\csname shm@ni.np.opt.b1\endcsname{0.55}% + NP offset 2%: optimum, largest |fundamental error|, %
\expandafter\def\csname shm@ni.np.lut.mean\endcsname{3.85}% + NP offset 2%: table, mean excess over the optimum, pp
\expandafter\def\csname shm@ni.np.lut.max\endcsname{62.9}% + NP offset 2%: table, worst excess over the optimum, pp
\expandafter\def\csname shm@ni.np.br.mean\endcsname{0.051}% + NP offset 2%: branch-aware table, mean excess, pp
\expandafter\def\csname shm@ni.np.br.max\endcsname{2.65}% + NP offset 2%: branch-aware table, worst excess, pp
\expandafter\def\csname shm@ni.tq.opt.dmax\endcsname{0.173}% + timer 1 us: optimum, largest |THD49 - closed form|, pp
\expandafter\def\csname shm@ni.tq.opt.b1\endcsname{0.61}% + timer 1 us: optimum, largest |fundamental error|, %
\expandafter\def\csname shm@ni.tq.lut.mean\endcsname{3.89}% + timer 1 us: table, mean excess over the optimum, pp
\expandafter\def\csname shm@ni.tq.lut.max\endcsname{63.3}% + timer 1 us: table, worst excess over the optimum, pp
\expandafter\def\csname shm@ni.tq.br.mean\endcsname{0.056}% + timer 1 us: branch-aware table, mean excess, pp
\expandafter\def\csname shm@ni.tq.br.max\endcsname{2.86}% + timer 1 us: branch-aware table, worst excess, pp
\expandafter\def\csname shm@ni.all.opt.dmax\endcsname{0.898}% all combined: optimum, largest |THD49 - closed form|, pp
\expandafter\def\csname shm@ni.all.opt.b1\endcsname{1.13}% all combined: optimum, largest |fundamental error|, %
\expandafter\def\csname shm@ni.all.lut.mean\endcsname{3.87}% all combined: table, mean excess over the optimum, pp
\expandafter\def\csname shm@ni.all.lut.max\endcsname{63.3}% all combined: table, worst excess over the optimum, pp
\expandafter\def\csname shm@ni.all.br.mean\endcsname{0.052}% all combined: branch-aware table, mean excess, pp
\expandafter\def\csname shm@ni.all.br.max\endcsname{2.65}% all combined: branch-aware table, worst excess, pp
\expandafter\def\csname shm@ni.opt.dmax\endcsname{0.898}% largest |THD49 - closed form| of the optimum under any condition, pp
\expandafter\def\csname shm@ni.lut.spread\endcsname{0.07}% spread of the table's mean excess across the conditions, pp
\expandafter\def\csname shm@ni.br.spread\endcsname{0.006}% spread of the branch-aware mean excess across the conditions, pp
\expandafter\def\csname shm@ni.ripple.pp\endcsname{7.3}% largest neutral-point ripple in the switching-function model, V peak to peak
\expandafter\def\csname shm@sim.ni.runs\endcsname{54}% circuit runs with datasheet switches that ran the commanded angles
\expandafter\def\csname shm@sim.ni.invalid\endcsname{0}% runs left out: commanded angles not run (must be 0)
\expandafter\def\csname shm@sim.ni.conds\endcsname{6}% non-ideal conditions simulated
\expandafter\def\csname shm@sim.ni.swfn.max\endcsname{0.013}% circuit minus switching-function model, optimal and branch-aware patterns, largest |difference|, pp
\expandafter\def\csname shm@sim.ni.swfn.lut.rel\endcsname{0.18}% the same for the table's angles, % of their THD49
\expandafter\def\csname shm@sim.ni.offset.max\endcsname{1.49}% largest |neutral-point offset| the circuit had in the analysis window, % of Vdc
\expandafter\def\csname shm@sim.ni.swfn.cmd.max\endcsname{0.606}% circuit minus switching-function model at the COMMANDED offset (not the measured one), largest |difference|, pp
\expandafter\def\csname shm@sim.ni.lut.rows\endcsname{100}% rows of the largest table simulated with datasheet switches
\expandafter\def\csname shm@sim.ni.ideal.opt.dev\endcsname{0.114}% ideal switches: optimal patterns, largest |circuit - closed form| THD49, pp
\expandafter\def\csname shm@sim.ni.ideal.b1\endcsname{0.04}% ideal switches: optimal patterns, largest |fundamental error| in the circuit, %
\expandafter\def\csname shm@sim.ni.ideal.ripple.pp\endcsname{24.3}% ideal switches: largest neutral-point voltage, V peak to peak
\expandafter\def\csname shm@sim.ni.ideal.offset.min\endcsname{-1.49}% ideal switches: smallest measured NP offset, % of Vdc
\expandafter\def\csname shm@sim.ni.ideal.offset.max\endcsname{-0.47}% ideal switches: largest measured NP offset, % of Vdc
\expandafter\def\csname shm@sim.ni.ideal.lut.circuit\endcsname{60.82}% ideal switches: table at its worst point, circuit THD49, %
\expandafter\def\csname shm@sim.ni.ideal.lut.ithd\endcsname{6.49}% ideal switches: the same, phase-current THD49, %
\expandafter\def\csname shm@sim.ni.ideal.branch.circuit\endcsname{1.13}% ideal switches: branch-aware table at that Ma, circuit THD49, %
\expandafter\def\csname shm@sim.ni.ideal.branch.ithd\endcsname{2.44}% ideal switches: the same, phase-current THD49, %
\expandafter\def\csname shm@sim.ni.drops.opt.dev\endcsname{0.149}% datasheet devices: optimal patterns, largest |circuit - closed form| THD49, pp
\expandafter\def\csname shm@sim.ni.drops.b1\endcsname{0.49}% datasheet devices: optimal patterns, largest |fundamental error| in the circuit, %
\expandafter\def\csname shm@sim.ni.drops.ripple.pp\endcsname{23.8}% datasheet devices: largest neutral-point voltage, V peak to peak
\expandafter\def\csname shm@sim.ni.drops.offset.min\endcsname{-1.44}% datasheet devices: smallest measured NP offset, % of Vdc
\expandafter\def\csname shm@sim.ni.drops.offset.max\endcsname{-0.45}% datasheet devices: largest measured NP offset, % of Vdc
\expandafter\def\csname shm@sim.ni.drops.lut.circuit\endcsname{61.14}% datasheet devices: table at its worst point, circuit THD49, %
\expandafter\def\csname shm@sim.ni.drops.lut.ithd\endcsname{6.49}% datasheet devices: the same, phase-current THD49, %
\expandafter\def\csname shm@sim.ni.drops.branch.circuit\endcsname{1.11}% datasheet devices: branch-aware table at that Ma, circuit THD49, %
\expandafter\def\csname shm@sim.ni.drops.branch.ithd\endcsname{2.37}% datasheet devices: the same, phase-current THD49, %
\expandafter\def\csname shm@sim.ni.dt.opt.dev\endcsname{0.180}% + dead time 2 us: optimal patterns, largest |circuit - closed form| THD49, pp
\expandafter\def\csname shm@sim.ni.dt.b1\endcsname{0.95}% + dead time 2 us: optimal patterns, largest |fundamental error| in the circuit, %
\expandafter\def\csname shm@sim.ni.dt.ripple.pp\endcsname{23.6}% + dead time 2 us: largest neutral-point voltage, V peak to peak
\expandafter\def\csname shm@sim.ni.dt.offset.min\endcsname{-1.41}% + dead time 2 us: smallest measured NP offset, % of Vdc
\expandafter\def\csname shm@sim.ni.dt.offset.max\endcsname{-0.45}% + dead time 2 us: largest measured NP offset, % of Vdc
\expandafter\def\csname shm@sim.ni.dt.lut.circuit\endcsname{61.57}% + dead time 2 us: table at its worst point, circuit THD49, %
\expandafter\def\csname shm@sim.ni.dt.lut.ithd\endcsname{6.52}% + dead time 2 us: the same, phase-current THD49, %
\expandafter\def\csname shm@sim.ni.dt.branch.circuit\endcsname{1.17}% + dead time 2 us: branch-aware table at that Ma, circuit THD49, %
\expandafter\def\csname shm@sim.ni.dt.branch.ithd\endcsname{2.36}% + dead time 2 us: the same, phase-current THD49, %
\expandafter\def\csname shm@sim.ni.np.opt.dev\endcsname{0.128}% + NP offset 2%: optimal patterns, largest |circuit - closed form| THD49, pp
\expandafter\def\csname shm@sim.ni.np.b1\endcsname{0.49}% + NP offset 2%: optimal patterns, largest |fundamental error| in the circuit, %
\expandafter\def\csname shm@sim.ni.np.ripple.pp\endcsname{27.2}% + NP offset 2%: largest neutral-point voltage, V peak to peak
\expandafter\def\csname shm@sim.ni.np.offset.min\endcsname{0.03}% + NP offset 2%: smallest measured NP offset, % of Vdc
\expandafter\def\csname shm@sim.ni.np.offset.max\endcsname{1.07}% + NP offset 2%: largest measured NP offset, % of Vdc
\expandafter\def\csname shm@sim.ni.np.lut.circuit\endcsname{61.14}% + NP offset 2%: table at its worst point, circuit THD49, %
\expandafter\def\csname shm@sim.ni.np.lut.ithd\endcsname{6.41}% + NP offset 2%: the same, phase-current THD49, %
\expandafter\def\csname shm@sim.ni.np.branch.circuit\endcsname{1.06}% + NP offset 2%: branch-aware table at that Ma, circuit THD49, %
\expandafter\def\csname shm@sim.ni.np.branch.ithd\endcsname{2.21}% + NP offset 2%: the same, phase-current THD49, %
\expandafter\def\csname shm@sim.ni.tq.opt.dev\endcsname{0.158}% + timer 1 us: optimal patterns, largest |circuit - closed form| THD49, pp
\expandafter\def\csname shm@sim.ni.tq.b1\endcsname{0.47}% + timer 1 us: optimal patterns, largest |fundamental error| in the circuit, %
\expandafter\def\csname shm@sim.ni.tq.ripple.pp\endcsname{23.9}% + timer 1 us: largest neutral-point voltage, V peak to peak
\expandafter\def\csname shm@sim.ni.tq.offset.min\endcsname{-1.44}% + timer 1 us: smallest measured NP offset, % of Vdc
\expandafter\def\csname shm@sim.ni.tq.offset.max\endcsname{-0.45}% + timer 1 us: largest measured NP offset, % of Vdc
\expandafter\def\csname shm@sim.ni.tq.lut.circuit\endcsname{61.17}% + timer 1 us: table at its worst point, circuit THD49, %
\expandafter\def\csname shm@sim.ni.tq.lut.ithd\endcsname{6.49}% + timer 1 us: the same, phase-current THD49, %
\expandafter\def\csname shm@sim.ni.tq.branch.circuit\endcsname{1.12}% + timer 1 us: branch-aware table at that Ma, circuit THD49, %
\expandafter\def\csname shm@sim.ni.tq.branch.ithd\endcsname{2.37}% + timer 1 us: the same, phase-current THD49, %
\expandafter\def\csname shm@sim.ni.all.opt.dev\endcsname{0.168}% all combined: optimal patterns, largest |circuit - closed form| THD49, pp
\expandafter\def\csname shm@sim.ni.all.b1\endcsname{0.93}% all combined: optimal patterns, largest |fundamental error| in the circuit, %
\expandafter\def\csname shm@sim.ni.all.ripple.pp\endcsname{27.0}% all combined: largest neutral-point voltage, V peak to peak
\expandafter\def\csname shm@sim.ni.all.offset.min\endcsname{0.06}% all combined: smallest measured NP offset, % of Vdc
\expandafter\def\csname shm@sim.ni.all.offset.max\endcsname{1.09}% all combined: largest measured NP offset, % of Vdc
\expandafter\def\csname shm@sim.ni.all.lut.circuit\endcsname{61.61}% all combined: table at its worst point, circuit THD49, %
\expandafter\def\csname shm@sim.ni.all.lut.ithd\endcsname{6.45}% all combined: the same, phase-current THD49, %
\expandafter\def\csname shm@sim.ni.all.branch.circuit\endcsname{1.12}% all combined: branch-aware table at that Ma, circuit THD49, %
\expandafter\def\csname shm@sim.ni.all.branch.ithd\endcsname{2.21}% all combined: the same, phase-current THD49, %
\expandafter\def\csname shm@time.total\endcsname{0.00}% hours for the whole run
\makeatother
 in the preamble, then \shm{key} in the text.
\makeatletter
\providecommand{\shm}[1]{\@ifundefined{shm@#1}{\textbf{??#1??}}{\csname shm@#1\endcsname}}
\expandafter\def\csname shm@angles\endcsname{15}% switching angles per quarter cycle
\expandafter\def\csname shm@orders\endcsname{16}% odd non-triplen orders in THD49 (5..49)
\expandafter\def\csname shm@ma.lo\endcsname{0.70}% lowest modulation index of the study
\expandafter\def\csname shm@ma.hi\endcsname{1.20}% highest modulation index of the study
\expandafter\def\csname shm@ma.step\endcsname{0.0025}% modulation-index step of the traced grid
\expandafter\def\csname shm@grid.points\endcsname{201}% grid points along Ma
\expandafter\def\csname shm@anchors\endcsname{11}% multistart anchors
\expandafter\def\csname shm@nstart\endcsname{150}% random starts per anchor
\expandafter\def\csname shm@distinct.mean\endcsname{104.2}% distinct local optima per anchor, mean
\expandafter\def\csname shm@distinct.max\endcsname{112}% distinct local optima at one anchor, max
\expandafter\def\csname shm@hitrate\endcsname{1.1}% % of converged random starts that reach the best optimum
\expandafter\def\csname shm@spread\endcsname{9.41}% median minus best THD49 among the distinct optima, pp
\expandafter\def\csname shm@branches\endcsname{76}% distinct continuous branches traced
\expandafter\def\csname shm@owners\endcsname{17}% branches that own part of the per-Ma optimum
\expandafter\def\csname shm@ownerchanges\endcsname{21}% owner changes of the per-Ma optimum over the grid
\expandafter\def\csname shm@jump.mean\endcsname{27.2}% mean jump of the optimal angle vector at an owner change, deg
\expandafter\def\csname shm@jump.max\endcsname{49.2}% largest jump of the optimal angle vector, deg
\expandafter\def\csname shm@stored\endcsname{13}% branches kept in the stored subset
\expandafter\def\csname shm@stored.gap.mean\endcsname{0.0004}% mean gap of the stored subset to the full envelope, pp
\expandafter\def\csname shm@stored.gap.max\endcsname{0.0282}% worst gap of the stored subset to the full envelope, pp
\expandafter\def\csname shm@coverage\endcsname{100.0}% % of grid points where some branch exists
\expandafter\def\csname shm@env.thd.mean\endcsname{2.13}% mean THD49 of the per-Ma optimum, %
\expandafter\def\csname shm@env.thd.min\endcsname{0.22}% lowest THD49 of the per-Ma optimum, %
\expandafter\def\csname shm@env.thd.max\endcsname{3.86}% highest THD49 of the per-Ma optimum, %
\expandafter\def\csname shm@nfe.branchset\endcsname{17696578}% objective evaluations to build the branch set
\expandafter\def\csname shm@time.branchset\endcsname{29.3}% minutes to build the branch set
\expandafter\def\csname shm@effangles.min\endcsname{12.0}% fewest effective angles on any branch (collapsed pairs removed)
\expandafter\def\csname shm@traces\endcsname{91}% continuation traces before merging (seeds and probes)
\expandafter\def\csname shm@probes\endcsname{50}% modulation indices of the probe multistart
\expandafter\def\csname shm@probe.nstart\endcsname{100}% random starts per probe
\expandafter\def\csname shm@probe.traces\endcsname{11}% traces added because a probe beat the envelope
\expandafter\def\csname shm@probe.hits\endcsname{2}% probes that found an optimum below the envelope
\expandafter\def\csname shm@fold.points\endcsname{24}% grid points searched just beyond the ends of owning branches
\expandafter\def\csname shm@fold.traces\endcsname{2}% traces added by that fold-end search
\expandafter\def\csname shm@dropped\endcsname{0}% traces of one or two grid points that own nothing, dropped
\expandafter\def\csname shm@distinct.mean.0p5\endcsname{102.5}% distinct optima per anchor, clustered at 0.5 deg, mean
\expandafter\def\csname shm@distinct.mean.2\endcsname{99.5}% distinct optima per anchor, clustered at 2 deg, mean
\expandafter\def\csname shm@distinct.levels\endcsname{95.6}% distinct THD49 levels (1e-3 pp apart) per anchor, mean
\expandafter\def\csname shm@distinct.within1pp\endcsname{10.4}% distinct optima within 1 pp of the best, per anchor, mean
\expandafter\def\csname shm@jump.wave.mean\endcsname{15.1}% share of the period in which the optimal voltage changes at an owner change, mean, %
\expandafter\def\csname shm@jump.wave.max\endcsname{48.8}% share of the period in which the optimal voltage changes at an owner change, largest, %
\expandafter\def\csname shm@jump.wave.min\endcsname{4.3}% share of the period in which the optimal voltage changes at an owner change, smallest, %
\expandafter\def\csname shm@jump.wave.at.maxangle\endcsname{4.3}% the same at the owner change with the largest angle move, %
\expandafter\def\csname shm@distinct.once\endcsname{83}% share of the distinct optima reached by a single random start, per anchor, mean, %
\expandafter\def\csname shm@distinct.chao1\endcsname{480}% Chao1 estimate of the total number of local optima per anchor, mean (the distinct count is a lower bound)
\expandafter\def\csname shm@lut.excess.R10\endcsname{21.9503}% LUT at R = 10: mean THD49 excess over the optimum, pp
\expandafter\def\csname shm@lut.excessmax.R10\endcsname{63.309}% LUT at R = 10: worst THD49 excess, pp
\expandafter\def\csname shm@lut.angmax.R10\endcsname{41.19}% LUT at R = 10: worst angle error, deg
\expandafter\def\csname shm@lut.fund.R10\endcsname{0.185}% LUT at R = 10: mean fundamental error, %
\expandafter\def\csname shm@lut.excess.R25\endcsname{13.3422}% LUT at R = 25: mean THD49 excess over the optimum, pp
\expandafter\def\csname shm@lut.excessmax.R25\endcsname{62.555}% LUT at R = 25: worst THD49 excess, pp
\expandafter\def\csname shm@lut.angmax.R25\endcsname{37.65}% LUT at R = 25: worst angle error, deg
\expandafter\def\csname shm@lut.fund.R25\endcsname{0.108}% LUT at R = 25: mean fundamental error, %
\expandafter\def\csname shm@lut.excess.R50\endcsname{7.6717}% LUT at R = 50: mean THD49 excess over the optimum, pp
\expandafter\def\csname shm@lut.excessmax.R50\endcsname{62.767}% LUT at R = 50: worst THD49 excess, pp
\expandafter\def\csname shm@lut.angmax.R50\endcsname{39.05}% LUT at R = 50: worst angle error, deg
\expandafter\def\csname shm@lut.fund.R50\endcsname{0.073}% LUT at R = 50: mean fundamental error, %
\expandafter\def\csname shm@lut.excess.R100\endcsname{3.7973}% LUT at R = 100: mean THD49 excess over the optimum, pp
\expandafter\def\csname shm@lut.excessmax.R100\endcsname{59.862}% LUT at R = 100: worst THD49 excess, pp
\expandafter\def\csname shm@lut.angmax.R100\endcsname{48.42}% LUT at R = 100: worst angle error, deg
\expandafter\def\csname shm@lut.fund.R100\endcsname{0.044}% LUT at R = 100: mean fundamental error, %
\expandafter\def\csname shm@nearest.excess.R10\endcsname{0.0610}% NEAREST at R = 10: mean THD49 excess over the optimum, pp
\expandafter\def\csname shm@nearest.excessmax.R10\endcsname{1.923}% NEAREST at R = 10: worst THD49 excess, pp
\expandafter\def\csname shm@nearest.angmax.R10\endcsname{52.70}% NEAREST at R = 10: worst angle error, deg
\expandafter\def\csname shm@nearest.fund.R10\endcsname{1.567}% NEAREST at R = 10: mean fundamental error, %
\expandafter\def\csname shm@nearest.excess.R25\endcsname{0.0100}% NEAREST at R = 25: mean THD49 excess over the optimum, pp
\expandafter\def\csname shm@nearest.excessmax.R25\endcsname{0.665}% NEAREST at R = 25: worst THD49 excess, pp
\expandafter\def\csname shm@nearest.angmax.R25\endcsname{40.62}% NEAREST at R = 25: worst angle error, deg
\expandafter\def\csname shm@nearest.fund.R25\endcsname{0.637}% NEAREST at R = 25: mean fundamental error, %
\expandafter\def\csname shm@nearest.excess.R50\endcsname{0.0005}% NEAREST at R = 50: mean THD49 excess over the optimum, pp
\expandafter\def\csname shm@nearest.excessmax.R50\endcsname{0.452}% NEAREST at R = 50: worst THD49 excess, pp
\expandafter\def\csname shm@nearest.angmax.R50\endcsname{49.20}% NEAREST at R = 50: worst angle error, deg
\expandafter\def\csname shm@nearest.fund.R50\endcsname{0.364}% NEAREST at R = 50: mean fundamental error, %
\expandafter\def\csname shm@nearest.excess.R100\endcsname{-0.0082}% NEAREST at R = 100: mean THD49 excess over the optimum, pp
\expandafter\def\csname shm@nearest.excessmax.R100\endcsname{0.150}% NEAREST at R = 100: worst THD49 excess, pp
\expandafter\def\csname shm@nearest.angmax.R100\endcsname{47.63}% NEAREST at R = 100: worst angle error, deg
\expandafter\def\csname shm@nearest.fund.R100\endcsname{0.269}% NEAREST at R = 100: mean fundamental error, %
\expandafter\def\csname shm@sqp.excess.R10\endcsname{4.7744}% LUT+SQP at R = 10: mean THD49 excess over the optimum, pp
\expandafter\def\csname shm@sqp.excessmax.R10\endcsname{31.525}% LUT+SQP at R = 10: worst THD49 excess, pp
\expandafter\def\csname shm@sqp.angmax.R10\endcsname{78.09}% LUT+SQP at R = 10: worst angle error, deg
\expandafter\def\csname shm@sqp.fund.R10\endcsname{0.003}% LUT+SQP at R = 10: mean fundamental error, %
\expandafter\def\csname shm@sqp.excess.R25\endcsname{2.4875}% LUT+SQP at R = 25: mean THD49 excess over the optimum, pp
\expandafter\def\csname shm@sqp.excessmax.R25\endcsname{30.501}% LUT+SQP at R = 25: worst THD49 excess, pp
\expandafter\def\csname shm@sqp.angmax.R25\endcsname{78.48}% LUT+SQP at R = 25: worst angle error, deg
\expandafter\def\csname shm@sqp.fund.R25\endcsname{0.000}% LUT+SQP at R = 25: mean fundamental error, %
\expandafter\def\csname shm@sqp.excess.R50\endcsname{1.3759}% LUT+SQP at R = 50: mean THD49 excess over the optimum, pp
\expandafter\def\csname shm@sqp.excessmax.R50\endcsname{29.799}% LUT+SQP at R = 50: worst THD49 excess, pp
\expandafter\def\csname shm@sqp.angmax.R50\endcsname{76.64}% LUT+SQP at R = 50: worst angle error, deg
\expandafter\def\csname shm@sqp.fund.R50\endcsname{0.000}% LUT+SQP at R = 50: mean fundamental error, %
\expandafter\def\csname shm@sqp.excess.R100\endcsname{0.5158}% LUT+SQP at R = 100: mean THD49 excess over the optimum, pp
\expandafter\def\csname shm@sqp.excessmax.R100\endcsname{29.182}% LUT+SQP at R = 100: worst THD49 excess, pp
\expandafter\def\csname shm@sqp.angmax.R100\endcsname{76.44}% LUT+SQP at R = 100: worst angle error, deg
\expandafter\def\csname shm@sqp.fund.R100\endcsname{0.000}% LUT+SQP at R = 100: mean fundamental error, %
\expandafter\def\csname shm@cont.excess.R10\endcsname{2.9051}% CONT at R = 10: mean THD49 excess over the optimum, pp
\expandafter\def\csname shm@cont.excessmax.R10\endcsname{7.882}% CONT at R = 10: worst THD49 excess, pp
\expandafter\def\csname shm@cont.angmax.R10\endcsname{49.71}% CONT at R = 10: worst angle error, deg
\expandafter\def\csname shm@cont.fund.R10\endcsname{0.040}% CONT at R = 10: mean fundamental error, %
\expandafter\def\csname shm@cont.excess.R25\endcsname{1.7159}% CONT at R = 25: mean THD49 excess over the optimum, pp
\expandafter\def\csname shm@cont.excessmax.R25\endcsname{6.730}% CONT at R = 25: worst THD49 excess, pp
\expandafter\def\csname shm@cont.angmax.R25\endcsname{48.59}% CONT at R = 25: worst angle error, deg
\expandafter\def\csname shm@cont.fund.R25\endcsname{0.007}% CONT at R = 25: mean fundamental error, %
\expandafter\def\csname shm@cont.excess.R50\endcsname{1.3021}% CONT at R = 50: mean THD49 excess over the optimum, pp
\expandafter\def\csname shm@cont.excessmax.R50\endcsname{5.197}% CONT at R = 50: worst THD49 excess, pp
\expandafter\def\csname shm@cont.angmax.R50\endcsname{48.43}% CONT at R = 50: worst angle error, deg
\expandafter\def\csname shm@cont.fund.R50\endcsname{0.003}% CONT at R = 50: mean fundamental error, %
\expandafter\def\csname shm@cont.excess.R100\endcsname{1.1490}% CONT at R = 100: mean THD49 excess over the optimum, pp
\expandafter\def\csname shm@cont.excessmax.R100\endcsname{3.979}% CONT at R = 100: worst THD49 excess, pp
\expandafter\def\csname shm@cont.angmax.R100\endcsname{48.42}% CONT at R = 100: worst angle error, deg
\expandafter\def\csname shm@cont.fund.R100\endcsname{0.002}% CONT at R = 100: mean fundamental error, %
\expandafter\def\csname shm@branch.excess.R10\endcsname{1.2992}% BRANCH at R = 10: mean THD49 excess over the optimum, pp
\expandafter\def\csname shm@branch.excessmax.R10\endcsname{13.031}% BRANCH at R = 10: worst THD49 excess, pp
\expandafter\def\csname shm@branch.angmax.R10\endcsname{52.62}% BRANCH at R = 10: worst angle error, deg
\expandafter\def\csname shm@branch.fund.R10\endcsname{0.004}% BRANCH at R = 10: mean fundamental error, %
\expandafter\def\csname shm@branch.excess.R25\endcsname{0.3848}% BRANCH at R = 25: mean THD49 excess over the optimum, pp
\expandafter\def\csname shm@branch.excessmax.R25\endcsname{10.243}% BRANCH at R = 25: worst THD49 excess, pp
\expandafter\def\csname shm@branch.angmax.R25\endcsname{52.62}% BRANCH at R = 25: worst angle error, deg
\expandafter\def\csname shm@branch.fund.R25\endcsname{0.018}% BRANCH at R = 25: mean fundamental error, %
\expandafter\def\csname shm@branch.excess.R50\endcsname{0.1626}% BRANCH at R = 50: mean THD49 excess over the optimum, pp
\expandafter\def\csname shm@branch.excessmax.R50\endcsname{6.860}% BRANCH at R = 50: worst THD49 excess, pp
\expandafter\def\csname shm@branch.angmax.R50\endcsname{52.62}% BRANCH at R = 50: worst angle error, deg
\expandafter\def\csname shm@branch.fund.R50\endcsname{0.003}% BRANCH at R = 50: mean fundamental error, %
\expandafter\def\csname shm@branch.excess.R100\endcsname{0.0706}% BRANCH at R = 100: mean THD49 excess over the optimum, pp
\expandafter\def\csname shm@branch.excessmax.R100\endcsname{2.023}% BRANCH at R = 100: worst THD49 excess, pp
\expandafter\def\csname shm@branch.angmax.R100\endcsname{50.94}% BRANCH at R = 100: worst angle error, deg
\expandafter\def\csname shm@branch.fund.R100\endcsname{0.003}% BRANCH at R = 100: mean fundamental error, %
\expandafter\def\csname shm@lut.floor.ratio\endcsname{0.17}% LUT mean excess at R = 100 divided by that at R = 10 (1 = no improvement at all)
\expandafter\def\csname shm@lut.worst.ratio\endcsname{0.95}% LUT worst excess at R = 100 divided by that at R = 10
\expandafter\def\csname shm@lut.slope\endcsname{-0.76}% log-log slope of the LUT mean excess vs R (resolution error alone: -2)
\expandafter\def\csname shm@branch.slope\endcsname{-2.00}% log-log slope of the branch-aware mean excess vs stored sets
\expandafter\def\csname shm@split.slope.same\endcsname{-2.78}% log-log slope of the same-owner (resolution) error vs R, median |excess|
\expandafter\def\csname shm@split.slope.straddle\endcsname{-0.05}% log-log slope of the straddling (structural) error vs R, median |excess|
\expandafter\def\csname shm@split.slope.same.mean\endcsname{-1.74}% log-log slope of the same-owner error vs R, mean
\expandafter\def\csname shm@split.slope.straddle.mean\endcsname{-0.03}% log-log slope of the straddling error vs R, mean
\expandafter\def\csname shm@split.frac.R10\endcsname{89.0}% % of test points whose bracketing rows straddle an owner change, R = 10
\expandafter\def\csname shm@split.straddle.R10\endcsname{24.661}% mean excess of straddling intervals, R = 10, pp
\expandafter\def\csname shm@split.same.R10\endcsname{0.00210}% mean excess of same-owner intervals, R = 10, pp
\expandafter\def\csname shm@split.straddle.median.R10\endcsname{19.070}% median |excess| of straddling intervals, R = 10, pp
\expandafter\def\csname shm@split.same.median.R10\endcsname{2.19e-03}% median |excess| of same-owner intervals, R = 10, pp
\expandafter\def\csname shm@split.same.p90.R10\endcsname{3.76e-03}% 90th percentile |excess| of same-owner intervals, R = 10, pp
\expandafter\def\csname shm@split.frac.R25\endcsname{46.0}% % of test points whose bracketing rows straddle an owner change, R = 25
\expandafter\def\csname shm@split.straddle.R25\endcsname{28.989}% mean excess of straddling intervals, R = 25, pp
\expandafter\def\csname shm@split.same.R25\endcsname{0.00123}% mean excess of same-owner intervals, R = 25, pp
\expandafter\def\csname shm@split.straddle.median.R25\endcsname{27.360}% median |excess| of straddling intervals, R = 25, pp
\expandafter\def\csname shm@split.same.median.R25\endcsname{3.63e-05}% median |excess| of same-owner intervals, R = 25, pp
\expandafter\def\csname shm@split.same.p90.R25\endcsname{1.95e-03}% 90th percentile |excess| of same-owner intervals, R = 25, pp
\expandafter\def\csname shm@split.frac.R50\endcsname{27.8}% % of test points whose bracketing rows straddle an owner change, R = 50
\expandafter\def\csname shm@split.straddle.R50\endcsname{27.581}% mean excess of straddling intervals, R = 50, pp
\expandafter\def\csname shm@split.same.R50\endcsname{0.00017}% mean excess of same-owner intervals, R = 50, pp
\expandafter\def\csname shm@split.straddle.median.R50\endcsname{24.542}% median |excess| of straddling intervals, R = 50, pp
\expandafter\def\csname shm@split.same.median.R50\endcsname{1.09e-05}% median |excess| of same-owner intervals, R = 50, pp
\expandafter\def\csname shm@split.same.p90.R50\endcsname{9.76e-05}% 90th percentile |excess| of same-owner intervals, R = 50, pp
\expandafter\def\csname shm@split.frac.R100\endcsname{16.8}% % of test points whose bracketing rows straddle an owner change, R = 100
\expandafter\def\csname shm@split.straddle.R100\endcsname{22.560}% mean excess of straddling intervals, R = 100, pp
\expandafter\def\csname shm@split.same.R100\endcsname{0.00005}% mean excess of same-owner intervals, R = 100, pp
\expandafter\def\csname shm@split.straddle.median.R100\endcsname{16.736}% median |excess| of straddling intervals, R = 100, pp
\expandafter\def\csname shm@split.same.median.R100\endcsname{3.22e-06}% median |excess| of same-owner intervals, R = 100, pp
\expandafter\def\csname shm@split.same.p90.R100\endcsname{1.29e-05}% 90th percentile |excess| of same-owner intervals, R = 100, pp
\expandafter\def\csname shm@split.straddle.Rmax\endcsname{22.560}% straddling-interval excess at R = 100, pp
\expandafter\def\csname shm@split.same.Rmax\endcsname{0.00005}% same-owner-interval excess at R = 100, pp
\expandafter\def\csname shm@basin.onowner\endcsname{61.4}% % of corrected points (R = 25) that land on the owning branch
\expandafter\def\csname shm@basin.onowner.near\endcsname{14.3}% % landing on the owning branch within two grid steps of an owner change
\expandafter\def\csname shm@basin.onowner.far\endcsname{79.5}% % landing on the owning branch farther from an owner change
\expandafter\def\csname shm@basin.excess.wrong\endcsname{6.438}% mean THD49 excess where the corrector landed on another branch, pp
\expandafter\def\csname shm@basin.ok.near\endcsname{28.6}% % of corrected points within 0.01 pp of the optimum, within two grid steps of an owner change
\expandafter\def\csname shm@basin.ok.far\endcsname{79.5}% % of corrected points within 0.01 pp of the optimum, farther from an owner change
\expandafter\def\csname shm@cont.cost.mean\endcsname{0.999}% mean THD49 cost of the continuity-constrained path, pp
\expandafter\def\csname shm@cont.cost.max\endcsname{2.282}% worst THD49 cost of the continuity-constrained path, pp
\expandafter\def\csname shm@cont.breaks\endcsname{0}% places where even the continuity-constrained path must jump
\expandafter\def\csname shm@cont.jumps.env\endcsname{20}% jumps (> 5 deg) the per-Ma optimum makes over the grid
\expandafter\def\csname shm@cont.jumps.zero\endcsname{11}% fewest jumps of a path whose mean THD cost is below 1e-4 pp (the rest only swap near-equal branches)
\expandafter\def\csname shm@cont.jumps.for.1\endcsname{0}% fewest jumps that keep the mean THD49 cost at or below 1 pp
\expandafter\def\csname shm@cont.jumps.for.0p1\endcsname{5}% fewest jumps that keep the mean THD49 cost at or below 0.1 pp
\expandafter\def\csname shm@cont.jumps.for.0p01\endcsname{8}% fewest jumps that keep the mean THD49 cost at or below 0.01 pp
\expandafter\def\csname shm@cont.cost.k1\endcsname{0.688}% mean THD49 cost of the best path allowed one jump, pp
\expandafter\def\csname shm@reg1.linear.K41\endcsname{0.0101}% 1-D regression, K = 41: angle MAE of linear, deg
\expandafter\def\csname shm@reg1.linear.K6\endcsname{0.1820}% 1-D regression, K = 6: angle MAE of linear, deg
\expandafter\def\csname shm@reg1.spline.K41\endcsname{0.0077}% 1-D regression, K = 41: angle MAE of spline, deg
\expandafter\def\csname shm@reg1.spline.K6\endcsname{0.1351}% 1-D regression, K = 6: angle MAE of spline, deg
\expandafter\def\csname shm@reg1.pchip.K41\endcsname{0.0068}% 1-D regression, K = 41: angle MAE of pchip, deg
\expandafter\def\csname shm@reg1.pchip.K6\endcsname{0.1359}% 1-D regression, K = 6: angle MAE of pchip, deg
\expandafter\def\csname shm@reg1.RBF.K41\endcsname{0.0092}% 1-D regression, K = 41: angle MAE of RBF, deg
\expandafter\def\csname shm@reg1.RBF.K6\endcsname{0.1741}% 1-D regression, K = 6: angle MAE of RBF, deg
\expandafter\def\csname shm@reg1.GP.K41\endcsname{0.0114}% 1-D regression, K = 41: angle MAE of GP, deg
\expandafter\def\csname shm@reg1.GP.K6\endcsname{0.2451}% 1-D regression, K = 6: angle MAE of GP, deg
\expandafter\def\csname shm@reg1.tanhLM.K41\endcsname{0.0119}% 1-D regression, K = 41: angle MAE of tanhLM, deg
\expandafter\def\csname shm@reg1.tanhLM.K6\endcsname{0.1463}% 1-D regression, K = 6: angle MAE of tanhLM, deg
\expandafter\def\csname shm@reg1.MLP.K41\endcsname{0.0441}% 1-D regression, K = 41: angle MAE of MLP, deg
\expandafter\def\csname shm@reg1.MLP.K6\endcsname{0.5077}% 1-D regression, K = 6: angle MAE of MLP, deg
\expandafter\def\csname shm@reg1.ANFIS.K41\endcsname{0.0140}% 1-D regression, K = 41: angle MAE of ANFIS, deg
\expandafter\def\csname shm@reg1.ANFIS.K6\endcsname{0.3298}% 1-D regression, K = 6: angle MAE of ANFIS, deg
\expandafter\def\csname shm@reg1.mlp.ratio\endcsname{6.4}% MLP angle MAE divided by the best interpolant (pchip) at K = 41
\expandafter\def\csname shm@reg1.mlp.seedcv\endcsname{23}% seed-to-seed coefficient of variation of the MLP MAE, %
\expandafter\def\csname shm@reg2.linear\endcsname{0.0555}% 2-D regression, largest lattice (64 points): angle MAE of linear, deg
\expandafter\def\csname shm@reg2.spline\endcsname{0.0384}% 2-D regression, largest lattice (64 points): angle MAE of spline, deg
\expandafter\def\csname shm@reg2.RBF\endcsname{0.0468}% 2-D regression, largest lattice (64 points): angle MAE of RBF, deg
\expandafter\def\csname shm@reg2.GP\endcsname{0.0598}% 2-D regression, largest lattice (64 points): angle MAE of GP, deg
\expandafter\def\csname shm@reg2.tanhLM\endcsname{0.0485}% 2-D regression, largest lattice (64 points): angle MAE of tanhLM, deg
\expandafter\def\csname shm@reg2.MLP\endcsname{0.0774}% 2-D regression, largest lattice (64 points): angle MAE of MLP, deg
\expandafter\def\csname shm@reg2.ANFIS\endcsname{0.0553}% 2-D regression, largest lattice (64 points): angle MAE of ANFIS, deg
\expandafter\def\csname shm@reg3.linear\endcsname{0.0189}% 3-D regression, largest lattice (216 points): angle MAE of linear, deg
\expandafter\def\csname shm@reg3.spline\endcsname{0.0120}% 3-D regression, largest lattice (216 points): angle MAE of spline, deg
\expandafter\def\csname shm@reg3.RBF\endcsname{0.0119}% 3-D regression, largest lattice (216 points): angle MAE of RBF, deg
\expandafter\def\csname shm@reg3.GP\endcsname{0.0208}% 3-D regression, largest lattice (216 points): angle MAE of GP, deg
\expandafter\def\csname shm@reg3.tanhLM\endcsname{0.0246}% 3-D regression, largest lattice (216 points): angle MAE of tanhLM, deg
\expandafter\def\csname shm@reg3.MLP\endcsname{0.0197}% 3-D regression, largest lattice (216 points): angle MAE of MLP, deg
\expandafter\def\csname shm@reg3.ANFIS\endcsname{0.0265}% 3-D regression, largest lattice (216 points): angle MAE of ANFIS, deg
\expandafter\def\csname shm@dstar.sel.K9\endcsname{4.33}% dimension from which a 9-per-axis label table outgrows the classifier
\expandafter\def\csname shm@dstar.reg.K9\endcsname{3.25}% dimension from which a 9-per-axis angle table outgrows one MLP per branch
\expandafter\def\csname shm@dstar.sel.K17\endcsname{3.36}% dimension from which a 17-per-axis label table outgrows the classifier
\expandafter\def\csname shm@dstar.reg.K17\endcsname{2.52}% dimension from which a 17-per-axis angle table outgrows one MLP per branch
\expandafter\def\csname shm@dstar.sel.K25\endcsname{2.95}% dimension from which a 25-per-axis label table outgrows the classifier
\expandafter\def\csname shm@dstar.reg.K25\endcsname{2.22}% dimension from which a 25-per-axis angle table outgrows one MLP per branch
\expandafter\def\csname shm@t.lut\endcsname{34.4}% table lookup + linear interpolation, microseconds per call
\expandafter\def\csname shm@t.pchip\endcsname{71.5}% pchip evaluation inside a branch, microseconds per call
\expandafter\def\csname shm@t.cls\endcsname{86.4}% softmax classifier forward pass, microseconds per call
\expandafter\def\csname shm@t.mlp\endcsname{267.8}% MLP forward pass, microseconds per call
\expandafter\def\csname shm@t.sqp\endcsname{30771.7}% SQP corrector (fmincon, warm start), microseconds per call
\expandafter\def\csname shm@dense.points\endcsname{61}% off-grid points of the independent multistart check
\expandafter\def\csname shm@dense.nstart\endcsname{200}% fresh random starts per off-grid point
\expandafter\def\csname shm@dense.gap.mean\endcsname{-0.4249}% stored branches minus independent multistart, mean, pp
\expandafter\def\csname shm@dense.gap.worst\endcsname{0.0022}% stored branches minus independent multistart, worst, pp
\expandafter\def\csname shm@dense.better\endcsname{32}% off-grid points where the stored branches beat the multistart
\expandafter\def\csname shm@dense.over1\endcsname{0}% off-grid points where the stored branches lose by more than 1 pp
\expandafter\def\csname shm@dense.over0p1\endcsname{0}% off-grid points where the stored branches lose by more than 0.1 pp
\expandafter\def\csname shm@dense.gap.worst.interp\endcsname{0.0268}% the same without the corrector: the table's own interpolation, worst, pp
\expandafter\def\csname shm@dense.better.max\endcsname{1.76}% largest margin by which the stored branches beat the fresh multistart, pp
\expandafter\def\csname shm@dense.interp.used\endcsname{0}% off-grid points whose stored value is a branch's interpolated angles (its corrector failed or left it)
\expandafter\def\csname shm@dense.corr.fail\endcsname{1}% off-grid points where the corrector of the best interpolated branch did not converge or left the branch
\expandafter\def\csname shm@spec.ma\endcsname{1.00}% Ma of the spectrum audit quoted
\expandafter\def\csname shm@spec.thd49\endcsname{2.44}% THD49 at that Ma, %
\expandafter\def\csname shm@spec.thd99\endcsname{25.7}% THD to the 99th at that Ma, %
\expandafter\def\csname shm@spec.thdfull\endcsname{34.2}% THD to the 999th at that Ma, %
\expandafter\def\csname shm@spec.wthd\endcsname{0.499}% weighted THD (1/n) to the 999th at that Ma, %
\expandafter\def\csname shm@spec.oob\endcsname{17.1}% largest single harmonic above the 49th, % of b1
\expandafter\def\csname shm@spec.oob.order\endcsname{55}% its order
\expandafter\def\csname shm@comply\endcsname{89.6}% % of the envelope meeting the per-order limits (0.2+32.5/n above the 25th)
\expandafter\def\csname shm@comply.alt\endcsname{35.8}% % meeting them with 0.2+12.5/n above the 25th
\expandafter\def\csname shm@bound.ub.model\endcsname{89.875}% the model's upper angle bound, 90 - theta_gap/2, deg
\expandafter\def\csname shm@bound.ub.min\endcsname{80}% tightest upper angle bound tested, deg
\expandafter\def\csname shm@bound.cost.mean\endcsname{0.825}% mean rise of the best THD49 from the model's bound to the tightest one, pp
\expandafter\def\csname shm@bound.cost.max\endcsname{1.688}% largest rise of the best THD49 from the model's bound to the tightest one, pp
\expandafter\def\csname shm@bound.cost89.max\endcsname{1.030}% largest rise of the best THD49 if the angles were kept below 89 deg, pp
\expandafter\def\csname shm@sw.dev\endcsname{1.6e-15}% worst relative deviation, closed form vs exact Fourier integral of the gate pattern
\expandafter\def\csname shm@sw.fftdev\endcsname{0.0096}% worst |THD49 sampled FFT - model| at zero dead time, pp
\expandafter\def\csname shm@sw.dthd.2us\endcsname{0.092}% largest change of THD49 of the optimal patterns caused by a 2 us dead time, pp
\expandafter\def\csname shm@sw.dthd.4us\endcsname{0.089}% largest change of THD49 of the optimal patterns caused by a 4 us dead time, pp
\expandafter\def\csname shm@sw.loss.1us\endcsname{0.0}% % of the envelope whose narrowest pulse is shorter than a 1 us dead time
\expandafter\def\csname shm@sw.loss.2us\endcsname{0.0}% % of the envelope whose narrowest pulse is shorter than a 2 us dead time
\expandafter\def\csname shm@sw.loss.4us\endcsname{0.0}% % of the envelope whose narrowest pulse is shorter than a 4 us dead time
\expandafter\def\csname shm@sw.loss.8us\endcsname{0.0}% % of the envelope whose narrowest pulse is shorter than a 8 us dead time
\expandafter\def\csname shm@sim.cases\endcsname{9}% steady-state cases whose circuit ran the commanded angles
\expandafter\def\csname shm@sim.invalid\endcsname{0}% cases left out: the circuit did not run the commanded angles (must be 0)
\expandafter\def\csname shm@sim.vdc\endcsname{800}% DC-link voltage of the circuit model, V
\expandafter\def\csname shm@sim.ts.us\endcsname{1}% time step of the logged line-to-line voltage, us
\expandafter\def\csname shm@sim.opt.dev.max\endcsname{0.098}% optimal patterns: largest |circuit - analytic| THD49, pp
\expandafter\def\csname shm@sim.opt.samp.max\endcsname{0.004}% optimal patterns: largest |sampled ideal - analytic| THD49 (the time step), pp
\expandafter\def\csname shm@sim.opt.circ.max\endcsname{0.101}% optimal patterns: largest |circuit - sampled ideal| THD49 (what the circuit adds), pp
\expandafter\def\csname shm@sim.opt.ma.dev\endcsname{0.23}% optimal patterns: largest relative error of the circuit's fundamental, %
\expandafter\def\csname shm@sim.lut.rows\endcsname{100}% rows of the largest table simulated at its worst point
\expandafter\def\csname shm@sim.lut.ma\endcsname{0.7925}% Ma of that worst point
\expandafter\def\csname shm@sim.lut.model\endcsname{60.90}% that table's interpolated angles: analytic THD49, %
\expandafter\def\csname shm@sim.lut.circuit\endcsname{60.92}% that table's interpolated angles: circuit THD49, %
\expandafter\def\csname shm@sim.lut.ithd\endcsname{5.53}% that table's interpolated angles: phase-current THD49, %
\expandafter\def\csname shm@sim.branch.model\endcsname{1.04}% branch-aware table at the same Ma: analytic THD49, %
\expandafter\def\csname shm@sim.branch.circuit\endcsname{0.98}% branch-aware table at the same Ma: circuit THD49, %
\expandafter\def\csname shm@sim.branch.ithd\endcsname{0.08}% branch-aware table at the same Ma: phase-current THD49, %
\expandafter\def\csname shm@sim.np.pp.max\endcsname{13.6}% largest peak-to-peak NP voltage over the steady-state windows, V
\expandafter\def\csname shm@sim.step.wave.pct\endcsname{48.8}% branch switch (the owner change that changes the waveform most): share of the period that changes, %
\expandafter\def\csname shm@sim.step.jump.deg\endcsname{30.7}% branch switch: largest angle move at the switch, deg
\expandafter\def\csname shm@sim.step.peak.jump\endcsname{0.37}% branch switch: peak phase-current deviation from the new steady state, % of I1
\expandafter\def\csname shm@sim.step.pre.jump\endcsname{2.38}% branch switch: how far apart the two steady-state currents are, % of I1
\expandafter\def\csname shm@sim.step.offset.jump\endcsname{0.05}% branch switch: mean current offset over the first cycle after the switch, % of I1
\expandafter\def\csname shm@sim.step.peak.stay\endcsname{0.22}% the same Ma step on the old branch (no branch switch): peak phase-current deviation, % of I1
\expandafter\def\csname shm@sim.step.angle.deg\endcsname{49.2}% second branch switch, at the largest angle move: that move, deg
\expandafter\def\csname shm@sim.step.angle.wave.pct\endcsname{4.3}% second branch switch: share of the period that changes, %
\expandafter\def\csname shm@sim.step.angle.peak\endcsname{0.18}% second branch switch: peak phase-current deviation from the new steady state, % of I1
\expandafter\def\csname shm@sim.step.angle.pre\endcsname{0.92}% second branch switch: how far apart the two steady-state currents are, % of I1
\expandafter\def\csname shm@ni.R\endcsname{100}% rows (sets) of the tables evaluated under non-ideal switching
\expandafter\def\csname shm@ni.conds\endcsname{6}% non-ideal conditions evaluated over the whole range
\expandafter\def\csname shm@ni.ideal.opt.dmax\endcsname{0.100}% ideal switches: optimum, largest |THD49 - closed form|, pp
\expandafter\def\csname shm@ni.ideal.opt.b1\endcsname{0.06}% ideal switches: optimum, largest |fundamental error|, %
\expandafter\def\csname shm@ni.ideal.lut.mean\endcsname{3.87}% ideal switches: table, mean excess over the optimum, pp
\expandafter\def\csname shm@ni.ideal.lut.max\endcsname{62.9}% ideal switches: table, worst excess over the optimum, pp
\expandafter\def\csname shm@ni.ideal.br.mean\endcsname{0.057}% ideal switches: branch-aware table, mean excess, pp
\expandafter\def\csname shm@ni.ideal.br.max\endcsname{2.85}% ideal switches: branch-aware table, worst excess, pp
\expandafter\def\csname shm@ni.drops.opt.dmax\endcsname{0.153}% datasheet devices: optimum, largest |THD49 - closed form|, pp
\expandafter\def\csname shm@ni.drops.opt.b1\endcsname{0.55}% datasheet devices: optimum, largest |fundamental error|, %
\expandafter\def\csname shm@ni.drops.lut.mean\endcsname{3.89}% datasheet devices: table, mean excess over the optimum, pp
\expandafter\def\csname shm@ni.drops.lut.max\endcsname{63.3}% datasheet devices: table, worst excess over the optimum, pp
\expandafter\def\csname shm@ni.drops.br.mean\endcsname{0.056}% datasheet devices: branch-aware table, mean excess, pp
\expandafter\def\csname shm@ni.drops.br.max\endcsname{2.82}% datasheet devices: branch-aware table, worst excess, pp
\expandafter\def\csname shm@ni.dt.opt.dmax\endcsname{0.455}% + dead time 2 us: optimum, largest |THD49 - closed form|, pp
\expandafter\def\csname shm@ni.dt.opt.b1\endcsname{1.10}% + dead time 2 us: optimum, largest |fundamental error|, %
\expandafter\def\csname shm@ni.dt.lut.mean\endcsname{3.92}% + dead time 2 us: table, mean excess over the optimum, pp
\expandafter\def\csname shm@ni.dt.lut.max\endcsname{63.7}% + dead time 2 us: table, worst excess over the optimum, pp
\expandafter\def\csname shm@ni.dt.br.mean\endcsname{0.057}% + dead time 2 us: branch-aware table, mean excess, pp
\expandafter\def\csname shm@ni.dt.br.max\endcsname{2.77}% + dead time 2 us: branch-aware table, worst excess, pp
\expandafter\def\csname shm@ni.np.opt.dmax\endcsname{0.737}% + NP offset 2%: optimum, largest |THD49 - closed form|, pp
\expandafter\def\csname shm@ni.np.opt.b1\endcsname{0.55}% + NP offset 2%: optimum, largest |fundamental error|, %
\expandafter\def\csname shm@ni.np.lut.mean\endcsname{3.85}% + NP offset 2%: table, mean excess over the optimum, pp
\expandafter\def\csname shm@ni.np.lut.max\endcsname{62.9}% + NP offset 2%: table, worst excess over the optimum, pp
\expandafter\def\csname shm@ni.np.br.mean\endcsname{0.051}% + NP offset 2%: branch-aware table, mean excess, pp
\expandafter\def\csname shm@ni.np.br.max\endcsname{2.65}% + NP offset 2%: branch-aware table, worst excess, pp
\expandafter\def\csname shm@ni.tq.opt.dmax\endcsname{0.173}% + timer 1 us: optimum, largest |THD49 - closed form|, pp
\expandafter\def\csname shm@ni.tq.opt.b1\endcsname{0.61}% + timer 1 us: optimum, largest |fundamental error|, %
\expandafter\def\csname shm@ni.tq.lut.mean\endcsname{3.89}% + timer 1 us: table, mean excess over the optimum, pp
\expandafter\def\csname shm@ni.tq.lut.max\endcsname{63.3}% + timer 1 us: table, worst excess over the optimum, pp
\expandafter\def\csname shm@ni.tq.br.mean\endcsname{0.056}% + timer 1 us: branch-aware table, mean excess, pp
\expandafter\def\csname shm@ni.tq.br.max\endcsname{2.86}% + timer 1 us: branch-aware table, worst excess, pp
\expandafter\def\csname shm@ni.all.opt.dmax\endcsname{0.898}% all combined: optimum, largest |THD49 - closed form|, pp
\expandafter\def\csname shm@ni.all.opt.b1\endcsname{1.13}% all combined: optimum, largest |fundamental error|, %
\expandafter\def\csname shm@ni.all.lut.mean\endcsname{3.87}% all combined: table, mean excess over the optimum, pp
\expandafter\def\csname shm@ni.all.lut.max\endcsname{63.3}% all combined: table, worst excess over the optimum, pp
\expandafter\def\csname shm@ni.all.br.mean\endcsname{0.052}% all combined: branch-aware table, mean excess, pp
\expandafter\def\csname shm@ni.all.br.max\endcsname{2.65}% all combined: branch-aware table, worst excess, pp
\expandafter\def\csname shm@ni.opt.dmax\endcsname{0.898}% largest |THD49 - closed form| of the optimum under any condition, pp
\expandafter\def\csname shm@ni.lut.spread\endcsname{0.07}% spread of the table's mean excess across the conditions, pp
\expandafter\def\csname shm@ni.br.spread\endcsname{0.006}% spread of the branch-aware mean excess across the conditions, pp
\expandafter\def\csname shm@ni.ripple.pp\endcsname{7.3}% largest neutral-point ripple in the switching-function model, V peak to peak
\expandafter\def\csname shm@sim.ni.runs\endcsname{54}% circuit runs with datasheet switches that ran the commanded angles
\expandafter\def\csname shm@sim.ni.invalid\endcsname{0}% runs left out: commanded angles not run (must be 0)
\expandafter\def\csname shm@sim.ni.conds\endcsname{6}% non-ideal conditions simulated
\expandafter\def\csname shm@sim.ni.swfn.max\endcsname{0.013}% circuit minus switching-function model, optimal and branch-aware patterns, largest |difference|, pp
\expandafter\def\csname shm@sim.ni.swfn.lut.rel\endcsname{0.18}% the same for the table's angles, % of their THD49
\expandafter\def\csname shm@sim.ni.offset.max\endcsname{1.49}% largest |neutral-point offset| the circuit had in the analysis window, % of Vdc
\expandafter\def\csname shm@sim.ni.swfn.cmd.max\endcsname{0.606}% circuit minus switching-function model at the COMMANDED offset (not the measured one), largest |difference|, pp
\expandafter\def\csname shm@sim.ni.lut.rows\endcsname{100}% rows of the largest table simulated with datasheet switches
\expandafter\def\csname shm@sim.ni.ideal.opt.dev\endcsname{0.114}% ideal switches: optimal patterns, largest |circuit - closed form| THD49, pp
\expandafter\def\csname shm@sim.ni.ideal.b1\endcsname{0.04}% ideal switches: optimal patterns, largest |fundamental error| in the circuit, %
\expandafter\def\csname shm@sim.ni.ideal.ripple.pp\endcsname{24.3}% ideal switches: largest neutral-point voltage, V peak to peak
\expandafter\def\csname shm@sim.ni.ideal.offset.min\endcsname{-1.49}% ideal switches: smallest measured NP offset, % of Vdc
\expandafter\def\csname shm@sim.ni.ideal.offset.max\endcsname{-0.47}% ideal switches: largest measured NP offset, % of Vdc
\expandafter\def\csname shm@sim.ni.ideal.lut.circuit\endcsname{60.82}% ideal switches: table at its worst point, circuit THD49, %
\expandafter\def\csname shm@sim.ni.ideal.lut.ithd\endcsname{6.49}% ideal switches: the same, phase-current THD49, %
\expandafter\def\csname shm@sim.ni.ideal.branch.circuit\endcsname{1.13}% ideal switches: branch-aware table at that Ma, circuit THD49, %
\expandafter\def\csname shm@sim.ni.ideal.branch.ithd\endcsname{2.44}% ideal switches: the same, phase-current THD49, %
\expandafter\def\csname shm@sim.ni.drops.opt.dev\endcsname{0.149}% datasheet devices: optimal patterns, largest |circuit - closed form| THD49, pp
\expandafter\def\csname shm@sim.ni.drops.b1\endcsname{0.49}% datasheet devices: optimal patterns, largest |fundamental error| in the circuit, %
\expandafter\def\csname shm@sim.ni.drops.ripple.pp\endcsname{23.8}% datasheet devices: largest neutral-point voltage, V peak to peak
\expandafter\def\csname shm@sim.ni.drops.offset.min\endcsname{-1.44}% datasheet devices: smallest measured NP offset, % of Vdc
\expandafter\def\csname shm@sim.ni.drops.offset.max\endcsname{-0.45}% datasheet devices: largest measured NP offset, % of Vdc
\expandafter\def\csname shm@sim.ni.drops.lut.circuit\endcsname{61.14}% datasheet devices: table at its worst point, circuit THD49, %
\expandafter\def\csname shm@sim.ni.drops.lut.ithd\endcsname{6.49}% datasheet devices: the same, phase-current THD49, %
\expandafter\def\csname shm@sim.ni.drops.branch.circuit\endcsname{1.11}% datasheet devices: branch-aware table at that Ma, circuit THD49, %
\expandafter\def\csname shm@sim.ni.drops.branch.ithd\endcsname{2.37}% datasheet devices: the same, phase-current THD49, %
\expandafter\def\csname shm@sim.ni.dt.opt.dev\endcsname{0.180}% + dead time 2 us: optimal patterns, largest |circuit - closed form| THD49, pp
\expandafter\def\csname shm@sim.ni.dt.b1\endcsname{0.95}% + dead time 2 us: optimal patterns, largest |fundamental error| in the circuit, %
\expandafter\def\csname shm@sim.ni.dt.ripple.pp\endcsname{23.6}% + dead time 2 us: largest neutral-point voltage, V peak to peak
\expandafter\def\csname shm@sim.ni.dt.offset.min\endcsname{-1.41}% + dead time 2 us: smallest measured NP offset, % of Vdc
\expandafter\def\csname shm@sim.ni.dt.offset.max\endcsname{-0.45}% + dead time 2 us: largest measured NP offset, % of Vdc
\expandafter\def\csname shm@sim.ni.dt.lut.circuit\endcsname{61.57}% + dead time 2 us: table at its worst point, circuit THD49, %
\expandafter\def\csname shm@sim.ni.dt.lut.ithd\endcsname{6.52}% + dead time 2 us: the same, phase-current THD49, %
\expandafter\def\csname shm@sim.ni.dt.branch.circuit\endcsname{1.17}% + dead time 2 us: branch-aware table at that Ma, circuit THD49, %
\expandafter\def\csname shm@sim.ni.dt.branch.ithd\endcsname{2.36}% + dead time 2 us: the same, phase-current THD49, %
\expandafter\def\csname shm@sim.ni.np.opt.dev\endcsname{0.128}% + NP offset 2%: optimal patterns, largest |circuit - closed form| THD49, pp
\expandafter\def\csname shm@sim.ni.np.b1\endcsname{0.49}% + NP offset 2%: optimal patterns, largest |fundamental error| in the circuit, %
\expandafter\def\csname shm@sim.ni.np.ripple.pp\endcsname{27.2}% + NP offset 2%: largest neutral-point voltage, V peak to peak
\expandafter\def\csname shm@sim.ni.np.offset.min\endcsname{0.03}% + NP offset 2%: smallest measured NP offset, % of Vdc
\expandafter\def\csname shm@sim.ni.np.offset.max\endcsname{1.07}% + NP offset 2%: largest measured NP offset, % of Vdc
\expandafter\def\csname shm@sim.ni.np.lut.circuit\endcsname{61.14}% + NP offset 2%: table at its worst point, circuit THD49, %
\expandafter\def\csname shm@sim.ni.np.lut.ithd\endcsname{6.41}% + NP offset 2%: the same, phase-current THD49, %
\expandafter\def\csname shm@sim.ni.np.branch.circuit\endcsname{1.06}% + NP offset 2%: branch-aware table at that Ma, circuit THD49, %
\expandafter\def\csname shm@sim.ni.np.branch.ithd\endcsname{2.21}% + NP offset 2%: the same, phase-current THD49, %
\expandafter\def\csname shm@sim.ni.tq.opt.dev\endcsname{0.158}% + timer 1 us: optimal patterns, largest |circuit - closed form| THD49, pp
\expandafter\def\csname shm@sim.ni.tq.b1\endcsname{0.47}% + timer 1 us: optimal patterns, largest |fundamental error| in the circuit, %
\expandafter\def\csname shm@sim.ni.tq.ripple.pp\endcsname{23.9}% + timer 1 us: largest neutral-point voltage, V peak to peak
\expandafter\def\csname shm@sim.ni.tq.offset.min\endcsname{-1.44}% + timer 1 us: smallest measured NP offset, % of Vdc
\expandafter\def\csname shm@sim.ni.tq.offset.max\endcsname{-0.45}% + timer 1 us: largest measured NP offset, % of Vdc
\expandafter\def\csname shm@sim.ni.tq.lut.circuit\endcsname{61.17}% + timer 1 us: table at its worst point, circuit THD49, %
\expandafter\def\csname shm@sim.ni.tq.lut.ithd\endcsname{6.49}% + timer 1 us: the same, phase-current THD49, %
\expandafter\def\csname shm@sim.ni.tq.branch.circuit\endcsname{1.12}% + timer 1 us: branch-aware table at that Ma, circuit THD49, %
\expandafter\def\csname shm@sim.ni.tq.branch.ithd\endcsname{2.37}% + timer 1 us: the same, phase-current THD49, %
\expandafter\def\csname shm@sim.ni.all.opt.dev\endcsname{0.168}% all combined: optimal patterns, largest |circuit - closed form| THD49, pp
\expandafter\def\csname shm@sim.ni.all.b1\endcsname{0.93}% all combined: optimal patterns, largest |fundamental error| in the circuit, %
\expandafter\def\csname shm@sim.ni.all.ripple.pp\endcsname{27.0}% all combined: largest neutral-point voltage, V peak to peak
\expandafter\def\csname shm@sim.ni.all.offset.min\endcsname{0.06}% all combined: smallest measured NP offset, % of Vdc
\expandafter\def\csname shm@sim.ni.all.offset.max\endcsname{1.09}% all combined: largest measured NP offset, % of Vdc
\expandafter\def\csname shm@sim.ni.all.lut.circuit\endcsname{61.61}% all combined: table at its worst point, circuit THD49, %
\expandafter\def\csname shm@sim.ni.all.lut.ithd\endcsname{6.45}% all combined: the same, phase-current THD49, %
\expandafter\def\csname shm@sim.ni.all.branch.circuit\endcsname{1.12}% all combined: branch-aware table at that Ma, circuit THD49, %
\expandafter\def\csname shm@sim.ni.all.branch.ithd\endcsname{2.21}% all combined: the same, phase-current THD49, %
\expandafter\def\csname shm@time.total\endcsname{0.00}% hours for the whole run
\makeatother
